% =============================================================================
% ACC 2027 CMG--PES paper: the six nominal controller families of Sec. VI-B
% equations, derivations, code, and architecture of each family
% =============================================================================
% Build from this folder (no Perl needed):
%   ..\..\.venv\Scripts\python.exe tools\build_family_tables.py     (from the repository root, see the script)
%   pdflatex ACC2027_controller_families
%   bibtex   ACC2027_controller_families
%   pdflatex ACC2027_controller_families   (x2)
% Paper numbers "P-(n)" refer to the manuscript PDF of 2026-09-30 (equations 1-30).
% Every listing is read from the repository file at the quoted lines (\lstinputlisting); every number in the
% generated/ tables is computed by tools/build_family_tables.py from the code or the A5 run outputs.
% =============================================================================
\documentclass[11pt]{article}
\usepackage[letterpaper,margin=0.85in]{geometry}
\usepackage{amsmath,amssymb,amsthm,bm,mathtools}
\usepackage{booktabs,array,longtable}
\usepackage[T1]{fontenc}
\usepackage{lmodern}
\usepackage{xcolor}
\usepackage{listings}
\usepackage{enumitem}
\usepackage{float}
\usepackage{needspace}
\usepackage{caption}
\captionsetup{font=small}
\captionsetup[table]{position=top,skip=5pt}
\usepackage{graphicx}
\usepackage{tikz}
\usetikzlibrary{arrows.meta,positioning,fit,backgrounds,calc}
\usepackage[obeyspaces,spaces]{url}
\urlstyle{tt}
\makeatletter\g@addto@macro\UrlBreaks{\do\_\do\.\do\-}\makeatother
\usepackage[colorlinks=true,linkcolor=blue!50!black,citecolor=green!35!black,urlcolor=blue!50!black]{hyperref}

\newcommand{\srcroot}{../..}
\newcommand{\T}{^{\mathsf T}}
\newcommand{\sat}{\operatorname{sat}}
\newcommand{\clip}{\operatorname{clip}}
\newcommand{\sgnp}{\operatorname{sgn}_{+}}
\newcommand{\vecpart}{\operatorname{vec}}
\newcommand{\Peq}[1]{P-(#1)}
\newcommand{\code}[1]{\path{#1}}
\newcommand{\ccode}[1]{\texttt{\detokenize{#1}}}% for captions (moving arguments)
\makeatletter\newcommand{\tabinput}[1]{\@@input #1 }\makeatother

\theoremstyle{remark}
\newtheorem*{status}{Status}

\definecolor{codebg}{gray}{0.97}
\definecolor{notebg}{RGB}{255,248,225}
\lstdefinestyle{py}{language=Python,basicstyle=\ttfamily\scriptsize,
  numbers=left,numberstyle=\tiny\color{gray},numbersep=6pt,frame=single,
  backgroundcolor=\color{codebg},rulecolor=\color{gray!60},breaklines=true,
  columns=fullflexible,keepspaces=true,showstringspaces=false,
  keywordstyle=\color{blue!60!black},commentstyle=\color{green!40!black},
  stringstyle=\color{red!50!black},captionpos=t,xleftmargin=1.5em,framexleftmargin=1.5em}
\newcommand{\listing}[4]{% file, first, last, title
  \Needspace{8\baselineskip}%
  \lstinputlisting[style=py,firstline=#2,lastline=#3,firstnumber=#2,
    title={\scriptsize\texttt{\detokenize{#1}}\,:\,#2--#3\quad #4}]{\srcroot/#1}}
\newcommand{\implnote}[1]{\par\smallskip\noindent\colorbox{notebg}{\parbox{\dimexpr\linewidth-2\fboxsep}{\small\textbf{Implementation note.} #1}}\par\smallskip}

% Release switch: the recursive-backstepping extension (Sec. bstrue) is kept out of the shared release until reviewed.
%   full version:     pdflatex ACC2027_controller_families
%   release version:  pdflatex -jobname=ACC2027_controller_families_release "\def\releasebuild{}% =============================================================================
% ACC 2027 CMG--PES paper: the six nominal controller families of Sec. VI-B
% equations, derivations, code, and architecture of each family
% =============================================================================
% Build from this folder (no Perl needed):
%   ..\..\.venv\Scripts\python.exe tools\build_family_tables.py     (from the repository root, see the script)
%   pdflatex ACC2027_controller_families
%   bibtex   ACC2027_controller_families
%   pdflatex ACC2027_controller_families   (x2)
% Paper numbers "P-(n)" refer to the manuscript PDF of 2026-09-30 (equations 1-30).
% Every listing is read from the repository file at the quoted lines (\lstinputlisting); every number in the
% generated/ tables is computed by tools/build_family_tables.py from the code or the A5 run outputs.
% =============================================================================
\documentclass[11pt]{article}
\usepackage[letterpaper,margin=0.85in]{geometry}
\usepackage{amsmath,amssymb,amsthm,bm,mathtools}
\usepackage{booktabs,array,longtable}
\usepackage[T1]{fontenc}
\usepackage{lmodern}
\usepackage{xcolor}
\usepackage{listings}
\usepackage{enumitem}
\usepackage{float}
\usepackage{needspace}
\usepackage{caption}
\captionsetup{font=small}
\captionsetup[table]{position=top,skip=5pt}
\usepackage{graphicx}
\usepackage{tikz}
\usetikzlibrary{arrows.meta,positioning,fit,backgrounds,calc}
\usepackage[obeyspaces,spaces]{url}
\urlstyle{tt}
\makeatletter\g@addto@macro\UrlBreaks{\do\_\do\.\do\-}\makeatother
\usepackage[colorlinks=true,linkcolor=blue!50!black,citecolor=green!35!black,urlcolor=blue!50!black]{hyperref}

\newcommand{\srcroot}{../..}
\newcommand{\T}{^{\mathsf T}}
\newcommand{\sat}{\operatorname{sat}}
\newcommand{\clip}{\operatorname{clip}}
\newcommand{\sgnp}{\operatorname{sgn}_{+}}
\newcommand{\vecpart}{\operatorname{vec}}
\newcommand{\Peq}[1]{P-(#1)}
\newcommand{\code}[1]{\path{#1}}
\newcommand{\ccode}[1]{\texttt{\detokenize{#1}}}% for captions (moving arguments)
\makeatletter\newcommand{\tabinput}[1]{\@@input #1 }\makeatother

\theoremstyle{remark}
\newtheorem*{status}{Status}

\definecolor{codebg}{gray}{0.97}
\definecolor{notebg}{RGB}{255,248,225}
\lstdefinestyle{py}{language=Python,basicstyle=\ttfamily\scriptsize,
  numbers=left,numberstyle=\tiny\color{gray},numbersep=6pt,frame=single,
  backgroundcolor=\color{codebg},rulecolor=\color{gray!60},breaklines=true,
  columns=fullflexible,keepspaces=true,showstringspaces=false,
  keywordstyle=\color{blue!60!black},commentstyle=\color{green!40!black},
  stringstyle=\color{red!50!black},captionpos=t,xleftmargin=1.5em,framexleftmargin=1.5em}
\newcommand{\listing}[4]{% file, first, last, title
  \Needspace{8\baselineskip}%
  \lstinputlisting[style=py,firstline=#2,lastline=#3,firstnumber=#2,
    title={\scriptsize\texttt{\detokenize{#1}}\,:\,#2--#3\quad #4}]{\srcroot/#1}}
\newcommand{\implnote}[1]{\par\smallskip\noindent\colorbox{notebg}{\parbox{\dimexpr\linewidth-2\fboxsep}{\small\textbf{Implementation note.} #1}}\par\smallskip}

% Release switch: the recursive-backstepping extension (Sec. bstrue) is kept out of the shared release until reviewed.
%   full version:     pdflatex ACC2027_controller_families
%   release version:  pdflatex -jobname=ACC2027_controller_families_release "\def\releasebuild{}% =============================================================================
% ACC 2027 CMG--PES paper: the six nominal controller families of Sec. VI-B
% equations, derivations, code, and architecture of each family
% =============================================================================
% Build from this folder (no Perl needed):
%   ..\..\.venv\Scripts\python.exe tools\build_family_tables.py     (from the repository root, see the script)
%   pdflatex ACC2027_controller_families
%   bibtex   ACC2027_controller_families
%   pdflatex ACC2027_controller_families   (x2)
% Paper numbers "P-(n)" refer to the manuscript PDF of 2026-09-30 (equations 1-30).
% Every listing is read from the repository file at the quoted lines (\lstinputlisting); every number in the
% generated/ tables is computed by tools/build_family_tables.py from the code or the A5 run outputs.
% =============================================================================
\documentclass[11pt]{article}
\usepackage[letterpaper,margin=0.85in]{geometry}
\usepackage{amsmath,amssymb,amsthm,bm,mathtools}
\usepackage{booktabs,array,longtable}
\usepackage[T1]{fontenc}
\usepackage{lmodern}
\usepackage{xcolor}
\usepackage{listings}
\usepackage{enumitem}
\usepackage{float}
\usepackage{needspace}
\usepackage{caption}
\captionsetup{font=small}
\captionsetup[table]{position=top,skip=5pt}
\usepackage{graphicx}
\usepackage{tikz}
\usetikzlibrary{arrows.meta,positioning,fit,backgrounds,calc}
\usepackage[obeyspaces,spaces]{url}
\urlstyle{tt}
\makeatletter\g@addto@macro\UrlBreaks{\do\_\do\.\do\-}\makeatother
\usepackage[colorlinks=true,linkcolor=blue!50!black,citecolor=green!35!black,urlcolor=blue!50!black]{hyperref}

\newcommand{\srcroot}{../..}
\newcommand{\T}{^{\mathsf T}}
\newcommand{\sat}{\operatorname{sat}}
\newcommand{\clip}{\operatorname{clip}}
\newcommand{\sgnp}{\operatorname{sgn}_{+}}
\newcommand{\vecpart}{\operatorname{vec}}
\newcommand{\Peq}[1]{P-(#1)}
\newcommand{\code}[1]{\path{#1}}
\newcommand{\ccode}[1]{\texttt{\detokenize{#1}}}% for captions (moving arguments)
\makeatletter\newcommand{\tabinput}[1]{\@@input #1 }\makeatother

\theoremstyle{remark}
\newtheorem*{status}{Status}

\definecolor{codebg}{gray}{0.97}
\definecolor{notebg}{RGB}{255,248,225}
\lstdefinestyle{py}{language=Python,basicstyle=\ttfamily\scriptsize,
  numbers=left,numberstyle=\tiny\color{gray},numbersep=6pt,frame=single,
  backgroundcolor=\color{codebg},rulecolor=\color{gray!60},breaklines=true,
  columns=fullflexible,keepspaces=true,showstringspaces=false,
  keywordstyle=\color{blue!60!black},commentstyle=\color{green!40!black},
  stringstyle=\color{red!50!black},captionpos=t,xleftmargin=1.5em,framexleftmargin=1.5em}
\newcommand{\listing}[4]{% file, first, last, title
  \Needspace{8\baselineskip}%
  \lstinputlisting[style=py,firstline=#2,lastline=#3,firstnumber=#2,
    title={\scriptsize\texttt{\detokenize{#1}}\,:\,#2--#3\quad #4}]{\srcroot/#1}}
\newcommand{\implnote}[1]{\par\smallskip\noindent\colorbox{notebg}{\parbox{\dimexpr\linewidth-2\fboxsep}{\small\textbf{Implementation note.} #1}}\par\smallskip}

% Release switch: the recursive-backstepping extension (Sec. bstrue) is kept out of the shared release until reviewed.
%   full version:     pdflatex ACC2027_controller_families
%   release version:  pdflatex -jobname=ACC2027_controller_families_release "\def\releasebuild{}% =============================================================================
% ACC 2027 CMG--PES paper: the six nominal controller families of Sec. VI-B
% equations, derivations, code, and architecture of each family
% =============================================================================
% Build from this folder (no Perl needed):
%   ..\..\.venv\Scripts\python.exe tools\build_family_tables.py     (from the repository root, see the script)
%   pdflatex ACC2027_controller_families
%   bibtex   ACC2027_controller_families
%   pdflatex ACC2027_controller_families   (x2)
% Paper numbers "P-(n)" refer to the manuscript PDF of 2026-09-30 (equations 1-30).
% Every listing is read from the repository file at the quoted lines (\lstinputlisting); every number in the
% generated/ tables is computed by tools/build_family_tables.py from the code or the A5 run outputs.
% =============================================================================
\documentclass[11pt]{article}
\usepackage[letterpaper,margin=0.85in]{geometry}
\usepackage{amsmath,amssymb,amsthm,bm,mathtools}
\usepackage{booktabs,array,longtable}
\usepackage[T1]{fontenc}
\usepackage{lmodern}
\usepackage{xcolor}
\usepackage{listings}
\usepackage{enumitem}
\usepackage{float}
\usepackage{needspace}
\usepackage{caption}
\captionsetup{font=small}
\captionsetup[table]{position=top,skip=5pt}
\usepackage{graphicx}
\usepackage{tikz}
\usetikzlibrary{arrows.meta,positioning,fit,backgrounds,calc}
\usepackage[obeyspaces,spaces]{url}
\urlstyle{tt}
\makeatletter\g@addto@macro\UrlBreaks{\do\_\do\.\do\-}\makeatother
\usepackage[colorlinks=true,linkcolor=blue!50!black,citecolor=green!35!black,urlcolor=blue!50!black]{hyperref}

\newcommand{\srcroot}{../..}
\newcommand{\T}{^{\mathsf T}}
\newcommand{\sat}{\operatorname{sat}}
\newcommand{\clip}{\operatorname{clip}}
\newcommand{\sgnp}{\operatorname{sgn}_{+}}
\newcommand{\vecpart}{\operatorname{vec}}
\newcommand{\Peq}[1]{P-(#1)}
\newcommand{\code}[1]{\path{#1}}
\newcommand{\ccode}[1]{\texttt{\detokenize{#1}}}% for captions (moving arguments)
\makeatletter\newcommand{\tabinput}[1]{\@@input #1 }\makeatother

\theoremstyle{remark}
\newtheorem*{status}{Status}

\definecolor{codebg}{gray}{0.97}
\definecolor{notebg}{RGB}{255,248,225}
\lstdefinestyle{py}{language=Python,basicstyle=\ttfamily\scriptsize,
  numbers=left,numberstyle=\tiny\color{gray},numbersep=6pt,frame=single,
  backgroundcolor=\color{codebg},rulecolor=\color{gray!60},breaklines=true,
  columns=fullflexible,keepspaces=true,showstringspaces=false,
  keywordstyle=\color{blue!60!black},commentstyle=\color{green!40!black},
  stringstyle=\color{red!50!black},captionpos=t,xleftmargin=1.5em,framexleftmargin=1.5em}
\newcommand{\listing}[4]{% file, first, last, title
  \Needspace{8\baselineskip}%
  \lstinputlisting[style=py,firstline=#2,lastline=#3,firstnumber=#2,
    title={\scriptsize\texttt{\detokenize{#1}}\,:\,#2--#3\quad #4}]{\srcroot/#1}}
\newcommand{\implnote}[1]{\par\smallskip\noindent\colorbox{notebg}{\parbox{\dimexpr\linewidth-2\fboxsep}{\small\textbf{Implementation note.} #1}}\par\smallskip}

% Release switch: the recursive-backstepping extension (Sec. bstrue) is kept out of the shared release until reviewed.
%   full version:     pdflatex ACC2027_controller_families
%   release version:  pdflatex -jobname=ACC2027_controller_families_release "\def\releasebuild{}\input{ACC2027_controller_families}"
\newif\ifwithbs \withbstrue
\ifdefined\releasebuild \withbsfalse \fi
% the manuscript's own figure style, plus the family extensions
\input{../v15_2026-09-29_corrected/architectures/tikz/arch_style.tex}
\input{figures/family_style.tex}
\input{generated/facts.tex}

\newcommand{\BL}{\code{experiments/nengo_nef_replication_2026-08-29/run_mujoco_nef_baselines.py}}
\newcommand{\RG}{\code{experiments/final_robustness_gating_2026-09-20/run_controllers_gated.py}}
\newcommand{\archfig}[3]{% file, caption, label
  \begin{figure}[H]\centering\resizebox{\textwidth}{!}{\input{figures/#1}}\caption{#2}\label{#3}\end{figure}}

\title{Online PES Decoder Learning in Neuromorphic Residual UAV Attitude Control\\[4pt]
\large The nominal controller families of Sec.~VI-B\ifwithbs{} and a recursive backstepping design\fi:\\ equations, derivations, code, and architectures\\
\normalsize companion to the ACC~2027 manuscript (version of September 30, 2026)}
\author{Gandhi Contreras, Steven Sandoval, and Luis Rodolfo Garcia Carrillo}
\date{October 1, 2026}

\begin{document}
\maketitle

\begin{abstract}
Section~VI-B of the manuscript attaches the same gated CMG--PES residual to six nominal attitude controllers: PD with
desired-rate feedforward in backstepping form (the primary baseline), PID with rate feedforward, LQR, sliding-mode control
(SMC), finite-horizon model predictive control (MPC), and grid-search nonlinear MPC (NMPC). This document gives, for each
family, its architecture, its control law, the derivation of its gains, the exact lines of code that implement it, and the
settings used in the portability study (Amendment A5). All six families share one torque interface and are analyzed in one
coordinate, the rotation-vector error $\bm e=2\bm e_q$, so their small-error behavior can be compared directly
(Table~\ref{tab:gains}): every family reduces near the reference to a second-order loop at the same bandwidth scale, with
damping ratios from $0.71$ (LQR, MPC) to $1$ (SMC). Each gain is read from the class the runs used, and each closed-form
derivation is checked numerically.\ifwithbs{} The family the code calls ``backstepping'' is computed torque; Section~\ref{sec:bstrue}
therefore adds a recursive Lyapunov backstepping design (plain and adaptive) in the style of the EHEKATL BWB controllers,
with its stability derivation, numerical verification, and a 40-run campaign with and without the residual. This new
family is not part of the manuscript.\fi{} The source files cited here are byte-identical to the ones recorded by every A5 run
(\fhashMatch/\fhashChecked{} hashes).
\end{abstract}

\tableofcontents
\newpage

% =============================================================================
\section{How to read this document}
% =============================================================================
\begin{itemize}[leftmargin=1.2em]
\item \textbf{Paper references.} ``\Peq{15}'' is equation (15) of the manuscript; Sec.~VI-B and Proposition~1 refer to the
  manuscript. The equation-to-code companion of September 30 (\code{paper/companion_2026-09-30}) covers the quaternion
  conventions, the residual and the gate in detail; this document does not repeat them.
\item \textbf{Code.} Every listing is read from the repository file at the quoted lines, with the file's own line numbers.
  Four files carry all six families: \BL{} (class \code{AdvancedBaseline}: backstepping, LQR, SMC, MPC, NMPC), \RG{}
  (the A5 plant loop and the PID path), and the campaign drivers \code{run_campaign_a5.py} and
  \code{run_campaign_a5_nmpc_corrected.py}. Appendix~\ref{app:hashes} gives their hashes.
\item \textbf{Status tags.} \emph{Standard}: textbook result, cited. \emph{Derived here}: the steps are written out.
  \emph{Verified}: recomputed by \code{tools/build_family_tables.py} from the code itself. \emph{Implementation note}: a
  detail of the code that the paper leaves implicit, with its effect on reported numbers.
\item \textbf{Figures.} Every architecture uses the manuscript's figure style. The left part (reference and attitude error)
  and the right part (torque box, gated residual, allocation, plant) are identical in all six figures; the dashed blue group is
  the only part that changes from family to family.
\end{itemize}

% =============================================================================
\section{What all six families share}\label{sec:common}
% =============================================================================
\subsection{The residual attaches to any nominal torque}
The residual reads only the nominal request $\bm u_0$ and adds a bounded, gated correction to it (Fig.~\ref{fig:interface}).
Proposition~1 assumes only $\bm u_{0,k}\in\mathcal U=\{\bm u:|u_i|\le\ell_i\}$, so it holds for every family whose output
is saturated to the torque box. All six are: LQR, backstepping, SMC and MPC through the shared clip (line~241 of \BL),
NMPC by clipping every candidate (lines 265 and 273), and PID at line~353 of \RG. Exact recovery at $\beta=0$ was certified
by running the augmentor with $\beta=0$ against the classical run: relative difference $0$ in all five cases (LQR, SMC, MPC
and two NMPC missions; \code{certify_beta0_recovery.py}). Backstepping is covered by the 20/20 bit-exact match with the
sealed primary runs; PID has no separate $\beta=0$ run, and for it the recovery rests on Proposition~1 alone.

\begin{figure}[H]\centering
\resizebox{\textwidth}{!}{\input{figures/fig_interface.tex}}
\caption{One residual, six nominal controllers. The nominal torque comes from one family per run; the residual path is the
manuscript's Fig.~1 unchanged ($\beta=.5$ on X2, $.2$ on CF2; $\eta=8\times10^{-4}$; $k_e=8$). The gate and the outer saturation
see only $\bm u_0$, so nothing in the residual depends on how $\bm u_0$ was computed.}
\label{fig:interface}
\end{figure}

\subsection{Error coordinates and the design model}\label{sec:design}
All families act on the shortest-path error $\bm e_q=\sgnp(q_{e,0})\bm q_{e,v}$ (\Peq{4}--\Peq{6}) and the rate error
$\bm\omega_e=\bm\omega-\bm\omega_d$, with $\bm\omega_d$ the central-difference reference rate of the companion, Sec.~2.5.
\code{AdvancedBaseline} works in the rotation-vector coordinate
\begin{equation}
\bm e=2\bm e_q,\qquad \|\bm e\|=2\sin\tfrac{\theta_q}{2}=\theta_q+O(\theta_q^3),
\label{eq:e}
\end{equation}
and designs every law for the per-axis double integrator obtained near the reference. From the error kinematics
\Peq{7}--\Peq{8}, $\dot{\bm e}=q_{e,0}\bm\omega_e+\bm e_q\times\bm\omega_e\approx\bm\omega_e$; from Euler's equation \Peq{2} with
the principal moments $\bm J=\bm j_{\rm tune}$ (Table~I of the paper),
\begin{equation}
\dot e_i\approx\omega_{e,i},\qquad \dot\omega_{e,i}=a_i-\dot\omega_{d,i}-\big[(\bm\omega\times\bm J\bm\omega)\oslash\bm J\big]_i,\qquad
a_i=\tau_i/J_i .
\label{eq:design}
\end{equation}
Every family outputs a desired angular acceleration $\bm a$; the torque is $\bm\tau=\bm J\odot\bm a$ plus, optionally, the
gyroscopic term, saturated to the box:
\begin{equation}
\bm u_0=\sat_{\bm\ell}\big(\bm J\odot\bm a+\chi_{\rm ff}\,\bm\omega\times(\bm J\odot\bm\omega)\big),\qquad
\bm a=\chi_{\rm ff}\dot{\bm\omega}_d+\bm a_{\rm fam}(\bm e,\bm\omega_e).
\label{eq:interface}
\end{equation}
\emph{Derived here.} In every run of the paper $\chi_{\rm ff}=0$ (\code{--no-model-ff}): the acceleration feedforward and the
gyroscopic term are off, so each family is a feedback law on $(\bm e,\bm\omega_e)$ scaled by inertia, and the residual is the
only adaptive element. The rate feedforward $\bm\omega_d$ stays on (\code{--rate-ff left-neg}) for all families.

\listing{experiments/nengo_nef_replication_2026-08-29/run_mujoco_nef_baselines.py}{189}{199}{$\bm\omega_d$ and $\dot{\bm\omega}_d$ (the latter is zero when model feedforward is off)}
\listing{experiments/nengo_nef_replication_2026-08-29/run_mujoco_nef_baselines.py}{216}{245}{\code{AdvancedBaseline.torque}: one interface, five families, shared clip at line 241}

\subsection{Torque box and platform constants}
$\ell_i=J_i\nu_i$ with $\bm\nu=(12,12,10)$\,rad/s$^2$, the same for both platforms, so the box limits every family to the
same angular acceleration $|a_i|\le\nu_i$:

\begin{table}[H]\centering\small
\caption{Constants read from the A5 run manifests (\ccode{derived} field).}\label{tab:platform}
\begin{tabular}{@{}lll@{}}\toprule
& X2 & CF2 (1.5$\times$ rotor design)\\\midrule
$\bm J$ (kg\,m$^2$) & \fJXtwo & \fJCFtwo\\
$\bm\ell$ (N\,m) & \fellXtwo & \fellCFtwo\\
$\bm\ell\oslash\bm J$ (rad/s$^2$) & \fnuXtwo & \fnuCFtwo\\
$\rho\bm\omega_n$ (rad/s) & 12, 12, 10.5 & 12, 12, 10.5\\\bottomrule
\end{tabular}\end{table}

Because the box is defined as an acceleration, every gain in this document is the same for X2 and CF2 (the script asserts
it), and the torque differs only by the inertia factor.

\subsection{The A5 settings}
All six families ran with the same plant, missions, residual, gate, $\rho=3$ bandwidth scale and seed; only the family flags
differ:
\listing{experiments/final_robustness_gating_2026-09-20/run_campaign_a5.py}{27}{40}{per-family flags}
\listing{experiments/final_robustness_gating_2026-09-20/run_campaign_a5.py}{49}{56}{common flags: $\rho=3$, rate feedforward on, model feedforward off, gate on in the CMG arm}

\implnote{The NMPC row of the paper comes from \code{run_campaign_a5_nmpc_corrected.py}, which drops the
\code{--baseline-rate-hz 100} hold of the first A5 run so that NMPC, like every other family, updates at every $0.5$-ms step
(Sec.~\ref{sec:nmpc}). The held run is kept on disk and not used.}

\subsection{Small-error equivalence of the six families}\label{sec:equiv}
Near the reference, inside the SMC boundary layer and away from the box, every family reduces to
$a_i=-k_{1,i}e_i-k_{2,i}\omega_{e,i}\,(-k_{I,i}\!\int\! e_i)$, i.e.\ the characteristic polynomial $s^2+k_2s+k_1$ (or
$s^3+k_2s^2+k_1s+k_I$ for PID). Table~\ref{tab:gains} lists the gains each family actually uses; Table~\ref{tab:poles} gives
the resulting poles of the linearized design model \eqref{eq:design}.

\begin{table}[H]\centering\small
\caption{Equivalent small-error gains, read from \ccode{AdvancedBaseline} with the A5 settings (roll\,/\,yaw; pitch equals roll).
$\omega_{\rm eq}=\sqrt{k_1}$, $\zeta_{\rm eq}=k_2/(2\sqrt{k_1})$. ``Box at'' is the rotation angle at which the
proportional term alone reaches the torque box, $\nu_i/k_{1,i}$. \emph{Verified.}}\label{tab:gains}
\begin{tabular}{@{}lllllll@{}}\toprule
Family & $k_1$ (s$^{-2}$) & $k_2$ (s$^{-1}$) & $\omega_{\rm eq}$ (rad/s) & $\zeta_{\rm eq}$ & box at ($^\circ$) & $k_I$ (s$^{-3}$)\\\midrule
\tabinput{generated/family_gains.tex}
\bottomrule\end{tabular}\end{table}

\begin{table}[H]\centering\small
\caption{Closed-loop poles of the linearized design model \eqref{eq:design} with each family's gains. \emph{Verified.}}\label{tab:poles}
\begin{tabular}{@{}lll@{}}\toprule
Family & axis & poles (s$^{-1}$)\\\midrule
\tabinput{generated/family_poles.tex}
\bottomrule\end{tabular}\end{table}

The families therefore differ in damping and in what happens away from the origin: how each one behaves at the box (SMC
uses a bounded reaching term, PID freezes its integrator, NMPC searches around the nominal law), not in bandwidth. The
proportional term of every linear law reaches the box at a rotation error of $5$--$6^\circ$, so Hardstep's steps of tens of
degrees saturate every family, which is where the gate matters.

% =============================================================================
\section{PD with desired-rate feedforward in computed-torque form (primary baseline)}\label{sec:pdff}
In the code this family is called \code{backstepping} (\code{--baseline backstepping}); it is the computed-torque form.\ifwithbs{}
It is not a recursive backstepping design; the recursive design is Sec.~\ref{sec:bstrue}.\fi
% =============================================================================
\archfig{fig_pdff.tex}{PD+ff in computed-torque (backstepping) form, the paper's $\rho=3$ primary baseline. The model
feedforward path (dotted) exists in the code and is switched off in every run, which makes this law identical to \Peq{15}.}{fig:pdff}

\paragraph{Law.} Computed-torque backstepping \cite{krstic1995nonlinear,slotine1991applied}: choose the error dynamics
$\dot{\bm\omega}_e=-k_1\bm e-k_2\bm\omega_e$ and invert \eqref{eq:design},
\begin{equation}
\bm u_0=\sat_{\bm\ell}\Big(\bm J\odot\big(\dot{\bm\omega}_d-k_1\bm e-k_2\bm\omega_e\big)+\bm\omega\times(\bm J\odot\bm\omega)\Big),\qquad
k_1=(\rho\omega_n)^2,\quad k_2=2\zeta\rho\omega_n,\ \zeta=0.95 .
\label{eq:bs}
\end{equation}
With the feedforward on, no saturation and a diagonal inertia, the gyroscopic term and $\dot{\bm\omega}_d$ cancel and the
rate error obeys $\dot{\bm\omega}_e=-k_1\bm e-k_2\bm\omega_e$ exactly; with $\dot{\bm e}\approx\bm\omega_e$ this is
$\ddot e_i+k_2\dot e_i+k_1e_i\approx0$ near the reference, the textbook reason for this form. \emph{Standard.}

\paragraph{Reduction to the paper's PD+ff.} With $\chi_{\rm ff}=0$ and $\bm e=2\bm e_q$,
$\bm u_0=\sat_{\bm\ell}(-2\bm J(\rho\bm\omega_n)^2\odot\bm e_q-2\zeta\bm J\rho\bm\omega_n\odot\bm\omega_e)$, which is \Peq{13}--\Peq{15}
with $\bm\omega_n\mapsto\rho\bm\omega_n$. \emph{Derived here; verified bit-exact in 20/20 cells}
(\code{run_campaign_a5_backstepping.py}, relative deviation $0$ against the sealed primary runs).

\listing{experiments/nengo_nef_replication_2026-08-29/run_mujoco_nef_baselines.py}{138}{141}{gains $k_1=\omega^2$, $k_2=2\zeta\omega$, $\omega=\rho\omega_n$}
\listing{experiments/final_robustness_gating_2026-09-20/run_controllers_gated.py}{79}{86}{the frozen paper baseline: backstepping, $\rho=3$, $\zeta=.95$, \code{left-neg}, no model feedforward}

\ifwithbs
% =============================================================================
\section{Recursive Lyapunov backstepping (new family)}\label{sec:bstrue}
% =============================================================================
The family of Sec.~\ref{sec:pdff} is computed torque: it chooses linear error dynamics and inverts the model, and with
model feedforward off it is exactly PD+ff. This section adds a backstepping controller in the strict sense: a recursive
design with a virtual control on the kinematic subsystem, a cross-coupling term that cancels the interconnection in the
Lyapunov derivative, and a final Lyapunov function that certifies the closed loop
\cite{krstic1995nonlinear,khalil2002nonlinear}. It follows step by step the design used for the longitudinal and adaptive
controllers of the EHEKATL BWB (strict-feedback form, errors $z_1,z_2$, virtual control with the term $-g_1\T z_1$,
$V=\sum\frac12z_i\T z_i$, La Salle--Yoshizawa, adaptive estimates with $\frac12\tilde f\T\Gamma\tilde f$). The plain and the
adaptive versions are implemented in \code{experiments/backstepping_true_2026-10-01/true_backstepping.py} and run through the
unchanged A5 runner. \textbf{They are not part of the manuscript.}

\archfig{fig_bstrue.tex}{Recursive backstepping. Step~1 stabilizes the quaternion kinematics with a virtual rate; step~2 drives
the rate to it with a torque that cancels the gyroscopic term, feeds forward the transported desired-rate derivative, and
contains the coupling term $-c\,z_1$. The adaptive variant adds a projected estimate of an unknown body torque.}{fig:bstrue}

\subsection{Strict-feedback form}
With the passive convention of the paper, $q_e=q_d\otimes q^{-1}$ and its body-frame rate
$\bm\omega_e=\bm\omega-\bar{\bm\omega}_d$, $\bar{\bm\omega}_d=R(q_e)\T\bm\omega_d$ (\Peq{7}--\Peq{8}, with the transport
$R(q_e)\T$ kept), the attitude system is in strict-feedback form:
\begin{equation}
\dot z_1=\tfrac12T(q_e)\,(\bm\omega-\bar{\bm\omega}_d),\qquad
\bm J\dot{\bm\omega}=-\bm\omega\times\bm J\bm\omega+\bm\tau+\bm d,\qquad
z_1=\bm e_q=\sgnp(q_{e,0})\bm q_{e,v},\quad T(q_e)=|q_{e,0}|I+[z_1]_\times ,
\label{eq:sf}
\end{equation}
i.e.\ $x_1=z_1$, $x_2=\bm\omega$, $f_1=-\frac12T\bar{\bm\omega}_d$, $g_1=\frac12T$, $f_2=-\bm J^{-1}(\bm\omega\times\bm J\bm\omega)$,
$g_2=\bm J^{-1}$ in the thesis notation, with $\bm d$ an unknown body torque (zero in the plain design). There is no actuator
state: the rotor forces of the MuJoCo models act without lag, so the thesis's $z_3$ step is not needed. The first identity
follows from \Peq{8}: $\frac{d}{dt}(\sgnp(q_{e,0})\bm q_{e,v})=\frac12(|q_{e,0}|\bm\omega_e+z_1\times\bm\omega_e)$.

\subsection{Step 1: kinematic subsystem}
Take $V_1=2(1-|q_{e,0}|)\ge0$, zero exactly at $q_e=\pm1$. With $\dot q_{e,0}=-\frac12\bm q_{e,v}\T\bm\omega_e$ (\Peq{8}),
\[
\dot V_1=-2\,\sgnp(q_{e,0})\,\dot q_{e,0}=z_1\T\bm\omega_e .
\]
Treat $\bm\omega_e$ as the control of this subsystem and choose the virtual rate $\bm\omega_e^d=-K_1z_1$, $K_1=\operatorname{diag}(k_{1,i})>0$:
if $\bm\omega_e=\bm\omega_e^d$, then $\dot V_1=-z_1\T K_1z_1$.

\subsection{Step 2: rate subsystem}
Define the second error $z_2=\bm\omega_e-\bm\omega_e^d=\bm\omega_e+K_1z_1$ and $V=c\,V_1+\frac12z_2\T\bm Jz_2$, $c>0$. Using
$\bm\omega_e=z_2-K_1z_1$ and $\dot{\bm\omega}_e=\dot{\bm\omega}-\dot{\bar{\bm\omega}}_d$,
\[
\dot V=c\,z_1\T z_2-c\,z_1\T K_1z_1+z_2\T\big(-\bm\omega\times\bm J\bm\omega+\bm\tau+\bm d-\bm J\dot{\bar{\bm\omega}}_d+\bm JK_1\dot z_1\big).
\]
The control
\begin{equation}
\boxed{\ \bm\tau=\bm\omega\times\bm J\bm\omega+\bm J\big(\dot{\bar{\bm\omega}}_d-K_1\dot z_1\big)-c\,z_1-K_2z_2-\hat{\bm d}\ },\qquad
\bm u_0=\sat_{\bm\ell}(\bm\tau),
\label{eq:bstrue}
\end{equation}
with $K_2=\operatorname{diag}(k_{2,i})>0$ and $\dot z_1=\frac12T(q_e)\bm\omega_e$ computed from measured signals, gives
\begin{equation}
\dot V=-c\,z_1\T K_1z_1-z_2\T K_2z_2+z_2\T\tilde{\bm d},\qquad \tilde{\bm d}=\bm d-\hat{\bm d}.
\label{eq:vdot}
\end{equation}
The term $-c\,z_1$ is the backstepping coupling term: it cancels $c\,z_1\T z_2$, the interconnection between the two
subsystems ($-g_1\T z_1$ in the thesis). Without it the design is computed torque.

\paragraph{Plain law} ($\bm d=\hat{\bm d}=\bm0$). $\dot V=-c\,z_1\T K_1z_1-z_2\T K_2z_2<0$ for $(z_1,z_2)\ne0$; by La Salle--Yoshizawa
$z_1,z_2\to0$, hence $q_e\to\pm1$ and $\bm\omega\to\bar{\bm\omega}_d$, and $z_1,z_2\in\mathcal L_2$. Because
$c\,V_1\ge c\|z_1\|^2/2$ near the origin, the convergence is exponential there.

\paragraph{Adaptive law.} For a constant unknown $\bm d$, add $\frac12\tilde{\bm d}\T\Gamma^{-1}\tilde{\bm d}$ to $V$. Since
$\dot{\tilde{\bm d}}=-\dot{\hat{\bm d}}$, the update
\begin{equation}
\dot{\hat{\bm d}}=\operatorname{Proj}_{|\hat d_i|\le D_i}\big(\Gamma z_2\big)
\label{eq:adapt}
\end{equation}
cancels $z_2\T\tilde{\bm d}$ in \eqref{eq:vdot} and leaves $\dot V_a\le-c\,z_1\T K_1z_1-z_2\T K_2z_2$; the projection keeps
$\hat{\bm d}$ bounded without increasing $\dot V_a$ when $|d_i|\le D_i$ \cite{ioannou1996robust,krstic1995nonlinear}. This is the
thesis's $\dot{\tilde f}=-\Gamma^{-1}z$ law with a bound added. Here $\Gamma_i=0.3\cdot2J_i\omega_i^2/k_{1,i}$, which makes the
integral action on $z_1$ ($\Gamma k_1\!\int\!z_1$) equal to the PID family's, and $D_i=0.25\,\ell_i$.

\paragraph{Rate transport.} $\bar{\bm\omega}_d=R(q_e)\T\bm\omega_d$ and, since $q_e$ rotates with body rate $\bm\omega_e$
($\dot R_e=R_e[\bm\omega_e]_\times$),
\[
\dot{\bar{\bm\omega}}_d=R(q_e)\T\dot{\bm\omega}_d-\bm\omega_e\times\bar{\bm\omega}_d ,
\]
with $\bm\omega_d$, $\dot{\bm\omega}_d$ the same central differences as Sec.~\ref{sec:design}. Unlike the A5 families, this law
uses the transport, the gyroscopic term and $\dot{\bm\omega}_d$ by construction; the \code{--no-model-ff} flag does not apply.

\subsection{Gain design: same linearization as PD+ff}
About $z=0$, \eqref{eq:sf}--\eqref{eq:bstrue} give, per axis, $\ddot z_{1,i}+(k_{1,i}/2+k_{2,i}/J_i)\dot z_{1,i}+(c+k_{1,i}k_{2,i})/(2J_i)\,z_{1,i}=0$.
Writing $k_{1,i}=2\lambda_{1,i}$, $k_{2,i}=J_i\lambda_{2,i}$ and imposing
\[
\lambda_{1,i}+\lambda_{2,i}=2\zeta\omega_i,\qquad \lambda_{1,i}\lambda_{2,i}+\frac{c}{2J_i}=\omega_i^2,\qquad \omega_i=\rho\omega_{n,i},\ \zeta=0.95,
\]
makes the small-error poles exactly those of the paper's PD+ff ($-11.40\pm3.75j$ on roll). Any difference between the two
families is then due to the backstepping structure, not to bandwidth or damping. $V_1$ is a scalar, so one $c$ serves all
axes; real $\lambda$'s need $c\ge2J_i\omega_i^2(1-\zeta^2)$ on every axis, and the smallest such value is used
($\lambda_1\le\lambda_2$: the outer kinematic loop is the slower one). For X2: $c=\fbsC$\,N\,m,
$\bm k_1=(\fbsKone)$\,s$^{-1}$, $\bm k_2=(\fbsKtwo)$\,N\,m\,s.

\subsection{Verification}
\code{verify_true_backstepping.py} checks the derivation on the rigid-body model (\code{results/verification.json}). \emph{Verified}:
\begin{itemize}[leftmargin=1.2em]
\item $\dot z_1$ and $\dot{\bar{\bm\omega}}_d$ in closed form against central finite differences along rigid-body trajectories on Full3,
  Hard45 and Yaw360: worst absolute error \fbsKinErr{} (signals of order $0.1$--$1$);
\item closed loop, plain law, no torque limit, Full3 and Hard45 from a $\approx20^\circ$ initial error: $V$ decreases monotonically to
  $10^{-12}$, and the finite-difference $\dot V$ equals $-c\,z_1\T K_1z_1-z_2\T K_2z_2$ within \fbsVdotErr{} (relative, step-size limited);
\item adaptive law with a constant unknown torque of $20\%$ of the box: $\hat{\bm d}$ reaches $\bm d$ within \fbsAdaptErr{} after $20$\,s;
\item linearized poles equal the PD+ff poles to \fbsPoleErr.
\end{itemize}

\listing{experiments/backstepping_true_2026-10-01/true_backstepping.py}{44}{55}{gain design: smallest common $c$, $\lambda_1\le\lambda_2$}
\listing{experiments/backstepping_true_2026-10-01/true_backstepping.py}{84}{114}{$z_1$, $z_2$, $\dot z_1$, transport, law \eqref{eq:bstrue}, projected adaptation \eqref{eq:adapt}}
\listing{experiments/backstepping_true_2026-10-01/run_true_backstepping.py}{30}{42}{the family enters the unchanged A5 runner}

\implnote{The Lyapunov certificate holds for the model \eqref{eq:sf} without saturation. In the plant the torque box, rotor
clipping, the non-diagonal X2 inertia (the law uses the principal moments, as every family does) and the residual add terms
that the certificate does not cover; the sign selection makes $V_1$ non-smooth at $q_{e,0}=0$, the same local scope as the
paper. The output is saturated to the box, so Proposition~1 applies unchanged.}

\input{generated/bs_results_section.tex}

\paragraph{What the two tables show.}
\begin{itemize}[leftmargin=1.2em]
\item \textbf{The recursive law is a much better nominal controller} (Table~\ref{tab:bsvsct}): with the same linearization it
  reduces the computed-torque PD+ff error on four of five missions on both platforms, most strongly on CF2. Its law contains
  $\bm J\dot{\bar{\bm\omega}}_d$ and $\bm\omega\times\bm J\bm\omega$, which every A5 family ran without (\code{--no-model-ff}). The
  computed-torque law \emph{with} model feedforward was not run, so this campaign does not separate the share of the gain due to the
  acceleration feedforward from the share due to the recursive structure (cross term and rate transport). That run is the
  natural next control.
\item \textbf{The unchanged residual does not help it} (Table~\ref{tab:bsres}). It still reduces the error on X2 Full3, Hard45
  and Yaw360, but where the recursive law's error is already very small (Roll; CF2 Roll and Yaw360, a few hundredths of a mrad) it
  raises the error to tenths of a mrad, worse in absolute terms than the residual on the paper's baseline. The residual's
  $\beta$, $\eta$ and $k_e$ were set for a lagging PD+ff; on a nominal law that already cancels the dominant error, the residual's
  own correction becomes the error. Retuning ($\beta$ in particular) was not tested. This is a limit of the unchanged residual,
  not of Proposition~1, which still holds: the output stays in the box and $\beta=0$ recovers the law.
\item \textbf{The adaptive variant fails on CF2 Hardstep} under persistent rotor saturation (79\% of samples): the projected estimate
  keeps integrating $z_2$ while the torque it requests cannot be delivered, a windup through the adaptation law that the
  Lyapunov argument (which assumes no saturation) does not exclude. Freezing the adaptation when the box or a rotor saturates, the
  rule the PID family uses, is the standard fix and was not run. On the other missions the adaptive and plain laws are close.
\end{itemize}
\fi % withbs

% =============================================================================
\section{PID with rate feedforward and rotor-level anti-windup}\label{sec:pid}
% =============================================================================
\archfig{fig_pid.tex}{PID+ff. The integrator runs on $\bm e_q$ and is frozen per axis whenever the torque box clipped that axis,
and on all axes whenever any rotor saturated at the previous step (dashed paths).}{fig:pid}

\paragraph{Law.} The PID path bypasses \code{AdvancedBaseline} (\code{--baseline pd}) and adds an integral term to \Peq{15}:
\begin{equation}
\bm u_0=\sat_{\bm\ell}\big(-\bm k_q\odot\bm e_q-\bm k_\omega\odot\bm\omega_e-\bm k_i\odot\bm z\big),\qquad
\bm k_q=2\bm J\odot(\rho\bm\omega_n)^2,\ \bm k_\omega=2\zeta\bm J\odot\rho\bm\omega_n,\ \bm k_i=0.3\,\bm k_q .
\label{eq:pid}
\end{equation}
\paragraph{Conditional integration.} With $h=5\times10^{-4}$\,s and the clip test of line 357,
\begin{equation}
z_{i,k+1}=z_{i,k}+h\,e_{q,i,k}\,m_{i,k},\qquad
m_{i,k}=\mathbf 1\big[u_{0,i,k}=\tau_{{\rm u},i,k}\big]\cdot\mathbf 1\big[\text{no rotor saturated at }k-1\big].
\label{eq:aw}
\end{equation}
This is conditional integration (integrator clamping) \cite{astrom2006advancedpid}, applied at the rotor level as well as at the
torque box: a rotor at its force limit means the realized torque differs from the command even if the box is inactive.
\emph{Standard; the rotor-level rule is this study's choice.}

\paragraph{Loop.} In the coordinate $\bm e$ the gains per unit inertia are $k_1=(\rho\omega_n)^2$, $k_2=2\zeta\rho\omega_n$ and
$k_I=0.3(\rho\omega_n)^2$ (Table~\ref{tab:gains}). The Routh condition for $s^3+k_2s^2+k_1s+k_I$ is $k_2k_1>k_I$, met with
a ratio of at least \fpidRouthMargin. The integral adds a slow real pole at $-0.32$\,s$^{-1}$ (time constant $\approx3$\,s,
Table~\ref{tab:poles}): the integrator is weak, which the A5 report discloses as ``undertuned''. \emph{Derived here; verified.}

\listing{experiments/final_robustness_gating_2026-09-20/run_controllers_gated.py}{253}{260}{gains; PID uses the plain PD path}
\listing{experiments/final_robustness_gating_2026-09-20/run_controllers_gated.py}{347}{361}{\eqref{eq:pid} and the conditional integrator \eqref{eq:aw}}
\listing{experiments/final_robustness_gating_2026-09-20/run_controllers_gated.py}{475}{475}{rotor-saturation flag for the next step}

\implnote{The integrator acts on $\bm e_q$ (not $\bm e=2\bm e_q$), so $\bm k_i=0.3\bm k_q$ means
$k_I=0.3(\rho\omega_n)^2$ in Table~\ref{tab:gains}'s coordinate. With \code{--ki-factor 0} this path is \Peq{15} exactly.}

% =============================================================================
\section{LQR}\label{sec:lqr}
% =============================================================================
\archfig{fig_lqr.tex}{LQR. The gain is computed once from the closed-form solution of the continuous algebraic Riccati
equation (CARE) for the per-axis double integrator; online it is a fixed state feedback.}{fig:lqr}

\paragraph{Design problem.} For each axis, $\bm x=[e,\omega_e]\T$, $\dot{\bm x}=A\bm x+Ba$ with
$A=\left[\begin{smallmatrix}0&1\\0&0\end{smallmatrix}\right]$, $B=[0,1]\T$, minimize
$\int_0^\infty(q_e e^2+ra^2)\,dt$ \cite{andersonmoore1990optimal}.

\paragraph{Closed-form CARE.} With $P=\left[\begin{smallmatrix}p_1&p_2\\p_2&p_3\end{smallmatrix}\right]$, the CARE
$A\T P+PA-PBr^{-1}B\T P+Q=0$ reads, entry by entry,
\[
(1,1):\ q_e-\frac{p_2^2}{r}=0,\qquad (1,2):\ p_1-\frac{p_2p_3}{r}=0,\qquad (2,2):\ 2p_2-\frac{p_3^2}{r}=0 .
\]
The positive-definite solution has $p_2=\sqrt{q_er}$, $p_3=\sqrt{2p_2r}$, $p_1=p_2p_3/r$, so
\begin{equation}
K=r^{-1}B\T P=\Big[\sqrt{q_e/r},\ \sqrt{2\sqrt{q_e/r}}\Big]\quad\xrightarrow{\ q_e/r=\omega^4\ }\quad K=[\omega^2,\ \sqrt2\,\omega],
\label{eq:lqr}
\end{equation}
a closed loop $s^2+\sqrt2\omega s+\omega^2$ with $\zeta=1/\sqrt2$ at the same $\omega=\rho\omega_n$ as the other families.
\emph{Derived here. Verified}: $P$ from these formulas leaves a CARE residual of \flqrCareRelErr{} (relative), its smallest
eigenvalue is \flqrPminEig{} $>0$, and $B\T P$ equals the code's gain.

\listing{experiments/nengo_nef_replication_2026-08-29/run_mujoco_nef_baselines.py}{132}{137}{$K=[\omega^2,\sqrt2\omega]$}
\listing{experiments/nengo_nef_replication_2026-08-29/run_mujoco_nef_baselines.py}{232}{233}{the law: $\bm a=\dot{\bm\omega}_d-k_1\bm e-k_2\bm\omega_e$}

% =============================================================================
\section{Sliding-mode control}\label{sec:smc}
% =============================================================================
\archfig{fig_smc.tex}{SMC. A linear sliding variable and a boundary-layer reaching law; inside the layer the law is a critically
damped PD, outside it the reaching term has a fixed magnitude of half the torque box.}{fig:smc}

\paragraph{Law.} Sliding variable and reaching law \cite{utkin1992sliding,slotine1991applied}, per axis:
\begin{equation}
\sigma_i=\omega_{e,i}+\lambda_ie_i,\qquad
a_i=\dot\omega_{d,i}-\lambda_i\omega_{e,i}-\eta_i\,\sat\!\big(\sigma_i/\phi_i\big),\qquad
\lambda_i=\rho\omega_{n,i},\ \ \eta_i=0.5\,\ell_i/J_i,\ \ \phi_i=\eta_i/\lambda_i .
\label{eq:smc}
\end{equation}
With the A5 setting, $\bm\eta=(\fsmcEtaXtwo)$\,rad/s$^2$ and $\bm\phi=(\fsmcPhiXtwo)$\,rad/s on both platforms.

\paragraph{Inside the boundary layer} ($|\sigma_i|<\phi_i$): $a_i=-\lambda_i\omega_{e,i}-(\eta_i/\phi_i)(\omega_{e,i}+\lambda_ie_i)$, and
$\eta_i/\phi_i=\lambda_i$, so
\[
a_i=-\lambda_i^2e_i-2\lambda_i\omega_{e,i}:\qquad k_1=\lambda^2,\quad k_2=2\lambda,\quad \zeta=1,
\]
a double pole at $-\rho\omega_n$ (Table~\ref{tab:poles}). \emph{Derived here; verified from the code's
\code{lam}, \code{eta}, \code{phi}.}

\paragraph{Outside the layer.} On the design model \eqref{eq:design} with model feedforward off,
$\dot\sigma_i=\dot\omega_{e,i}+\lambda_i\dot e_i=a_i-\dot\omega_{d,i}+\lambda_i\omega_{e,i}=-\eta_i\operatorname{sgn}(\sigma_i)-\dot\omega_{d,i}$,
hence
\[
\sigma_i\dot\sigma_i\le-(\eta_i-|\dot\omega_{d,i}|)\,|\sigma_i| ,
\]
the reaching condition, whenever the reference acceleration stays below $\eta_i$ ($6$\,rad/s$^2$ on roll and pitch). The layer
is then reached in finite time and the error decays inside it; the boundary layer trades exact sliding for a continuous,
chattering-free command. Because $J_i\eta_i=\ell_i/2$, the reaching term alone uses half the torque box, leaving the
other half for $-\lambda\bm\omega_e$. \emph{Standard argument, specialized here.}

\listing{experiments/nengo_nef_replication_2026-08-29/run_mujoco_nef_baselines.py}{142}{149}{$\lambda=\omega$, $\eta=0.5\,\bm\ell\oslash\bm J$, $\phi=\eta/\omega$}
\listing{experiments/nengo_nef_replication_2026-08-29/run_mujoco_nef_baselines.py}{234}{236}{\eqref{eq:smc}}

% =============================================================================
\section{Model predictive control (unconstrained, finite horizon)}\label{sec:mpc}
% =============================================================================
\archfig{fig_mpc.tex}{MPC. The finite-horizon quadratic program has no constraints, so its solution is linear in the state; the
first move of the optimal sequence is a fixed gain $K_0$, computed once and applied at every plant step.}{fig:mpc}

\paragraph{Prediction model.} Per axis, the exact zero-order-hold discretization of the double integrator with sample time
$T=20$\,ms,
\[
\bm x_{k+1}=A_d\bm x_k+B_da_k,\qquad A_d=\begin{bmatrix}1&T\\0&1\end{bmatrix},\quad B_d=\begin{bmatrix}T^2/2\\T\end{bmatrix}.
\]
\paragraph{Condensed problem.} Stacking $N=120$ predictions, $X=S_x\bm x_0+S_uU$ with
$S_x=[A_d;A_d^2;\dots;A_d^N]$ and the block-Toeplitz $S_u$ with blocks $A_d^{k-j}B_d$ ($j\le k$). The cost
$\sum_{k=1}^{N}\bm x_k\T Q\bm x_k+\sum_{k=0}^{N-1}Ra_k^2$ with $Q=\operatorname{diag}(\omega^4,q_\omega\omega^2)$, $q_\omega=0$, $R=1$ becomes
\[
U\T HU+2\bm x_0\T F\T U+\text{const},\qquad H=S_u\T\bar QS_u+\bar R,\quad F=S_u\T\bar QS_x,
\]
with $\bar Q=I_N\otimes Q$, $\bar R=I_N\otimes R$. Without constraints the minimizer is $U^\star=-H^{-1}F\bm x_0$, and the
receding-horizon law applies its first element, $a_0=-K_0\bm x_0$ with $K_0=[H^{-1}F]_{1,:}$
\cite{borrelli2017predictive,rawlings2017mpc}. \emph{Standard.}

\paragraph{What the horizon does.} Table~\ref{tab:mpc} recomputes $K_0$ with the code's own function for growing $N$. From
$N\approx30$ ($0.6$\,s, several closed-loop time constants) the first-move gain no longer changes: at the A5 horizon it equals
the infinite-horizon discrete LQR gain obtained by iterating the Riccati recursion to its fixed point to a relative
\fmpcVsDare{} ($K_0=[\fmpcKone,\ \fmpcKtwo]$ on roll; DARE residual \fdareResidual). The MPC row is therefore a discrete LQR
designed at $T=20$\,ms. That sample time, not the horizon, is why its proportional gain is $16\%$ below the continuous LQR
of Sec.~\ref{sec:lqr} ($121.5$ vs.\ $144$\,s$^{-2}$). \emph{Verified.}

\begin{table}[H]\centering\small
\caption{First-move MPC gain on the roll axis ($\omega=12$\,rad/s) versus horizon, from \ccode{AdvancedBaseline._mpc_first_gain};
the last two columns divide by the continuous LQR gain $[\omega^2,\sqrt2\omega]$. A5 uses $N=120$.}\label{tab:mpc}
\begin{tabular}{@{}rrrrrr@{}}\toprule
$N$ & $NT$ (s) & $K_{0,1}$ & $K_{0,2}$ & $K_{0,1}/\omega^2$ & $K_{0,2}/(\sqrt2\omega)$\\\midrule
\tabinput{generated/mpc_horizon.tex}
\bottomrule\end{tabular}\end{table}

\listing{experiments/nengo_nef_replication_2026-08-29/run_mujoco_nef_baselines.py}{167}{187}{condensed matrices $S_x,S_u,H,F$ and the first-move gain}
\listing{experiments/nengo_nef_replication_2026-08-29/run_mujoco_nef_baselines.py}{237}{238}{the online law $\bm a=\dot{\bm\omega}_d-K_0[\bm e;\bm\omega_e]$}

\implnote{The gain is designed for a 20-ms move but applied as a continuous-time state feedback at every 0.5-ms step (A5 ran
every family at full rate). This is a common simplification; it is the reason the MPC row behaves like a softer LQR
($\zeta_{\rm eq}=0.71$, $\omega_{\rm eq}=11.0$\,rad/s) and not like a sampled controller. ``MPC'' here means an unconstrained
receding-horizon gain; the torque box is enforced by the shared clip, not inside the optimization.}

% =============================================================================
\section{Nonlinear MPC (grid search around the backstepping law)}\label{sec:nmpc}
% =============================================================================
\archfig{fig_nmpc.tex}{NMPC. At every plant step, 27 candidate offsets are added to the nominal backstepping torque; each
candidate is scored by an RK4 rollout of the closed loop on the rigid-body model over $0.3$\,s, and the first torque of the
cheapest candidate is applied.}{fig:nmpc}

\paragraph{Candidates.} The nominal law is \eqref{eq:bs} ($\zeta=0.95$). A grid of $3^3=\fnmpcCandidates$ constant offsets
$\bm d_c\in\{-1,0,1\}^3\odot0.1\bm\ell$ is added to it; $\bm d_c=\bm0$ is one of them (\fnmpcHasZero), so the nominal law is always
a candidate.

\paragraph{Prediction (``policy'' mode).} From the measured state, each candidate is rolled out for $N=5$ steps of
$\Delta=60$\,ms with the rigid-body model $\dot Q=\frac12Q\otimes\tilde{\bm\omega}$,
$\dot{\bm\omega}=(\bm\tau-\bm\omega\times\bm J\bm\omega)\oslash\bm J$, integrated by classical RK4 \cite{hairer1993ode}. At
$k=0$ the candidate's first torque $\bm\tau^{(c)}_0=\sat_{\bm\ell}(\bm\tau_{\rm nom}+\bm d_c)$ is used; at $k\ge1$ the nominal law is
re-evaluated on the predicted state and the same offset added,
$\bm\tau^{(c)}_k=\sat_{\bm\ell}(\bm J\odot\bm a(\hat q_k,\hat{\bm\omega}_k,t+k\Delta)+\bm d_c)$. The prediction thus models the
closed loop that will actually run, not a constant torque.

\paragraph{Cost and selection.} With $\bm e_k,\bm\omega_{e,k}$ the predicted errors against the reference at $t+(k+1)\Delta$,
\begin{equation}
J_c=\sum_{k=0}^{N-1}\Big[w_k\big(\textstyle\sum_i\omega_i^4e_{k,i}^2+2\sum_i\omega_i^2\omega_{e,k,i}^2\big)+\big\|\bm\tau^{(c)}_k\oslash\bm J\big\|^2\Big],
\qquad w_k=\begin{cases}1.5,&k=N-1\\1,&\text{otherwise}\end{cases}
\label{eq:nmpc}
\end{equation}
and $\bm u_0=\bm\tau^{(c^\star)}_0$, $c^\star=\arg\min_cJ_c$. The angle weight $\omega^4$ is the LQR weight $q_e/r$ of
Sec.~\ref{sec:lqr}; the effort term is the squared commanded acceleration. This is a sampling-based (grid) receding-horizon
controller in the sense of \cite{rawlings2017mpc}, not a gradient-based NMPC solver. \emph{Derived here from the code.}

\listing{experiments/nengo_nef_replication_2026-08-29/run_mujoco_nef_baselines.py}{152}{159}{candidate grid and cost weights}
\listing{experiments/nengo_nef_replication_2026-08-29/run_mujoco_nef_baselines.py}{201}{214}{rigid-body model and RK4 step (MuJoCo's body-to-world convention)}
\listing{experiments/nengo_nef_replication_2026-08-29/run_mujoco_nef_baselines.py}{247}{289}{\code{_nmpc}: rollout, cost \eqref{eq:nmpc}, selection, one-period self-check}
\listing{experiments/final_robustness_gating_2026-09-20/run_campaign_a5_nmpc_corrected.py}{25}{26}{corrected A5 settings: no update hold}

\implnote{(i) The rollout predicts the nominal family only; it does not know the residual that will be added. (ii) The first A5
NMPC run updated at 100\,Hz and held the torque between updates, a leftover from an earlier exploratory setting. The audit
(\code{AUDIT_NMPC_HOLD_RATE_2026-09-25.md}) showed that the hold, not the predictive machinery, caused that row's anomaly:
removing it restored NMPC, and imposing the same hold on MPC broke MPC identically. The paper's NMPC row is the full-rate re-run.
(iii) The class logs a one-period-ahead self-check of its own prediction (\code{nmpc_pred_err_mean_deg}); it was used to
diagnose (ii), not as a tuning signal.}

% =============================================================================
\section{Results across families (Sec.~VI-B)}\label{sec:results}
% =============================================================================
\begin{table}[H]\centering\small
\caption{Reduction of mean quaternion-angle error by the gated CMG--PES residual on each nominal family, noise-free, $\rho=3$,
45-s missions (\% , positive is better; ``5-mission'' is the reduction of geometric means). Read from
\ccode{results/a5_numbers.json}, which \ccode{analyze_a5.py} computes from the run summaries; NMPC from the corrected
full-rate run. Last column: baseline geometric mean (mrad).}\label{tab:results}
\begin{tabular}{@{}llrrrrrrr@{}}\toprule
Family & Platform & Roll & Full3 & Hard45 & Hardstep & Yaw360 & 5-mission & base (mrad)\\\midrule
\tabinput{generated/family_results.tex}
\bottomrule\end{tabular}\end{table}

These are the numbers the manuscript reports (X2/CF2: backstepping $37.3/22.4$, PID+ff $36.7/22.4$, LQR $31.1/22.4$, SMC
$43.8/26.6$, MPC $33.4/22.7$, NMPC $31.9/21.7$), and the claim verifier checks them (\code{reproduce/verify_acc2027_claims.py}).
What the table shows, without further interpretation:
\begin{itemize}[leftmargin=1.2em]
\item The five-mission reduction stays within $31$--$44\%$ on X2 and $22$--$27\%$ on CF2 across families whose baseline errors
  differ by up to $1.7\times$ (X2 geometric means $3.7$--$6.3$\,mrad), with the residual unchanged.
\item The reduction is largest on Roll and Yaw360 and smallest on Hardstep, whose step changes of tens of degrees drive every
  linear family far past the $5$--$6^\circ$ at which its proportional term reaches the box (Table~\ref{tab:gains}).
\item SMC is the only family with a clearly positive Hardstep result ($+34.2\%$ X2, $+12.2\%$ CF2). SMC is also the only
  family whose law is itself bounded below the box away from the origin (Sec.~\ref{sec:smc}); A5 did not isolate whether that
  is the cause, so it is recorded here as an open observation.
\end{itemize}

% =============================================================================
\section{Scope of the family comparison}
% =============================================================================
\begin{itemize}[leftmargin=1.2em]
\item The families are \emph{attachment tests}, not tuned competitors: all use the same bandwidth scale $\rho=3$, model
  feedforward off, and the A5 settings above. No family was retuned per mission or per platform, and the residual
  ($\beta$, $\eta$, $k_e$) is the same for all.
\item MPC is unconstrained and realized as its first-move gain; NMPC is a 27-point grid search around backstepping. Neither is a
  constrained optimal controller; the torque box acts through the shared saturation.
\item PID's integrator is weak (slow pole at $-0.32$\,s$^{-1}$) and is disclosed as undertuned in the A5 report.
\item All gains are designed on the per-axis double integrator \eqref{eq:design} with principal moments; the X2 inertia tensor is
  not diagonal, so the poles of Table~\ref{tab:poles} are design values, not measured closed-loop poles.
\end{itemize}

\appendix
\section{Source files and hashes}\label{app:hashes}
SHA-256 (first 16 hex digits) of each cited file, compared with the \code{source_sha256} field of one A5 run per family and
platform (Roll, classical arm; NMPC from the corrected run): \fhashMatch/\fhashChecked{} recorded hashes are identical.
\begin{table}[H]\centering\small
\begin{tabular}{@{}p{0.55\linewidth}ll@{}}\toprule
File & SHA-256 & vs.\ A5 manifests\\\midrule
\tabinput{generated/family_hashes.tex}
\bottomrule\end{tabular}\end{table}

\bibliographystyle{IEEEtran}
\bibliography{../v15_2026-09-29_corrected/references_ACC2027,family_refs}
\end{document}
"
\newif\ifwithbs \withbstrue
\ifdefined\releasebuild \withbsfalse \fi
% the manuscript's own figure style, plus the family extensions
% Shared TikZ style for all architecture figures. Grayscale-safe.
\tikzset{
  blk/.style   ={draw,rounded corners=2pt,align=center,font=\footnotesize,
                 minimum height=8mm,inner sep=3pt,fill=white},
  clas/.style  ={blk,fill=black!4},
  neuro/.style ={blk,fill=black!12,very thick},
  plant/.style ={blk,fill=black!20},
  meas/.style  ={blk,dashed},
  sum/.style   ={draw,circle,inner sep=0pt,minimum size=4.5mm,font=\scriptsize},
  sig/.style   ={-{Latex[length=2mm]},thick},
  sigd/.style  ={-{Latex[length=2mm]},thick,dashed},
  lbl/.style   ={font=\scriptsize,inner sep=1.5pt,fill=white},
  grp/.style   ={draw,dashed,rounded corners=3pt,inner sep=6pt},
  note/.style   ={font=\scriptsize\itshape,align=left},
}

% Extra styles for the controller-family figures; the base styles (blk, clas, neuro, plant, sig, ...) are the
% manuscript's own architectures/tikz/arch_style.tex, so every figure matches the paper's Fig. 1.
\tikzset{
  law/.style    ={blk,fill=blue!7,draw=blue!45!black,thick},
  offl/.style   ={blk,dashed,fill=black!2},
  fam/.style    ={draw=blue!45!black,dashed,rounded corners=3pt,inner sep=7pt},
  famlbl/.style ={font=\scriptsize\itshape,text=blue!45!black,inner sep=1pt},
  off/.style    ={-{Latex[length=2mm]},thick,dotted,gray},
  every node/.append style={font=\footnotesize},
}
% Common left half: reference generator and attitude error (identical for every family).
\newcommand{\commonleft}{%
  \node[clas] (ref) at (0,0) {\textbf{Mission reference}\\$q_d(t)$\\[1pt]
        $\bm\omega_d$: central difference};
  \node[clas,right=7mm of ref] (err) {\textbf{Attitude error}\\$q_e=q_d\otimes q^{-1}$\\
        $\bm e=2\,\sgnp(q_{e,0})\,\bm q_{e,v}$\\$\bm\omega_e=\bm\omega-\bm\omega_d$};
  \draw[sig] (ref) -- node[lbl,above]{$q_d,\bm\omega_d$} (err);
}
% Common right half and feedback of every family figure. Needs nodes (err) and (law-out) placed by the figure,
% and the coordinate (satx) giving the x position of the torque box.
\newcommand{\commonloop}[2]{% #1: label on u0 arrow source node, #2: text inside the torque-box node
  \node[clas] (sat) at (satx) {#2};
  \node[neuro,below=13mm of sat] (res) {\textbf{Gated CMG--PES residual}\\(Fig.~\ref{fig:interface})\\
        $\bm u=\sat_{\bm\ell}(\bm u_0+\beta\bm\ell\odot\tilde{\bm s})$};
  \node[plant,anchor=east] (plant) at ($(res.west)+(-12mm,0)$) {\textbf{Allocation and MuJoCo plant}\\
        $\bm f=\clip_{[\bm0,\bm f_{\max}]}(B^{-1}[T,\bm u\T]\T)$\\X2 / CF2 rigid body};
  \draw[sig] (#1) -- (sat);
  \draw[sig] (sat) -- node[lbl,right]{$\bm u_0\in\mathcal U$} (res);
  \draw[sig] (res) -- node[lbl,above]{$\bm u$} (plant);
  \draw[sig] (plant.west) -- ++(-4mm,0) coordinate (fb1) -| node[lbl,pos=.75,left]{measured $q,\ \bm\omega$} (err.south);
}

\newcommand{\flqrCareRelErr}{2.7e-18}
\newcommand{\flqrPminEig}{7.39}
\newcommand{\fsmcEtaXtwo}{6.00, 6.00, 5.00}
\newcommand{\fsmcPhiXtwo}{0.500, 0.500, 0.476}
\newcommand{\fsmcEtaCFtwo}{6.00, 6.00, 5.00}
\newcommand{\fsmcPhiCFtwo}{0.500, 0.500, 0.476}
\newcommand{\fpidRouthMargin}{66.5}
\newcommand{\fdareResidual}{9.5e-14}
\newcommand{\fmpcDareKone}{121.55}
\newcommand{\fmpcDareKtwo}{15.592}
\newcommand{\fmpcKone}{121.55}
\newcommand{\fmpcKtwo}{15.592}
\newcommand{\fmpcVsDare}{4.6e-10}
\newcommand{\fnmpcCandidates}{27}
\newcommand{\fnmpcHasZero}{yes}
\newcommand{\fhashMatch}{60}
\newcommand{\fhashChecked}{60}
\newcommand{\fbsC}{1.705}
\newcommand{\fbsKone}{22.80, 16.69, 10.40}
\newcommand{\fbsKtwo}{0.692, 0.527, 0.375}
\newcommand{\fbsKinErr}{$5.6\times10^{-7}$}
\newcommand{\fbsVdotErr}{0.22\%}
\newcommand{\fbsAdaptErr}{0.25\%}
\newcommand{\fbsPoleErr}{$2.7\times10^{-15}$}
\newcommand{\fbsGeotrueXtwo}{-0.4}
\newcommand{\fbsGeotrueCFtwo}{-230.0}
\newcommand{\fbsGeoadaptiveXtwo}{-4.3}
\newcommand{\fbsGeoadaptiveCFtwo}{-204.7}
\newcommand{\fbsCampaignComplete}{yes}
\newcommand{\fJXtwo}{0.06071, 0.03647, 0.02541}
\newcommand{\fellXtwo}{0.7285, 0.4376, 0.2541}
\newcommand{\fnuXtwo}{12.00, 12.00, 10.00}
\newcommand{\fJCFtwo}{2.395e-05, 2.395e-05, 3.235e-05}
\newcommand{\fellCFtwo}{0.0002874, 0.0002874, 0.0003235}
\newcommand{\fnuCFtwo}{12.00, 12.00, 10.00}


\newcommand{\BL}{\code{experiments/nengo_nef_replication_2026-08-29/run_mujoco_nef_baselines.py}}
\newcommand{\RG}{\code{experiments/final_robustness_gating_2026-09-20/run_controllers_gated.py}}
\newcommand{\archfig}[3]{% file, caption, label
  \begin{figure}[H]\centering\resizebox{\textwidth}{!}{\input{figures/#1}}\caption{#2}\label{#3}\end{figure}}

\title{Online PES Decoder Learning in Neuromorphic Residual UAV Attitude Control\\[4pt]
\large The nominal controller families of Sec.~VI-B\ifwithbs{} and a recursive backstepping design\fi:\\ equations, derivations, code, and architectures\\
\normalsize companion to the ACC~2027 manuscript (version of September 30, 2026)}
\author{Gandhi Contreras, Steven Sandoval, and Luis Rodolfo Garcia Carrillo}
\date{October 1, 2026}

\begin{document}
\maketitle

\begin{abstract}
Section~VI-B of the manuscript attaches the same gated CMG--PES residual to six nominal attitude controllers: PD with
desired-rate feedforward in backstepping form (the primary baseline), PID with rate feedforward, LQR, sliding-mode control
(SMC), finite-horizon model predictive control (MPC), and grid-search nonlinear MPC (NMPC). This document gives, for each
family, its architecture, its control law, the derivation of its gains, the exact lines of code that implement it, and the
settings used in the portability study (Amendment A5). All six families share one torque interface and are analyzed in one
coordinate, the rotation-vector error $\bm e=2\bm e_q$, so their small-error behavior can be compared directly
(Table~\ref{tab:gains}): every family reduces near the reference to a second-order loop at the same bandwidth scale, with
damping ratios from $0.71$ (LQR, MPC) to $1$ (SMC). Each gain is read from the class the runs used, and each closed-form
derivation is checked numerically.\ifwithbs{} The family the code calls ``backstepping'' is computed torque; Section~\ref{sec:bstrue}
therefore adds a recursive Lyapunov backstepping design (plain and adaptive) in the style of the EHEKATL BWB controllers,
with its stability derivation, numerical verification, and a 40-run campaign with and without the residual. This new
family is not part of the manuscript.\fi{} The source files cited here are byte-identical to the ones recorded by every A5 run
(\fhashMatch/\fhashChecked{} hashes).
\end{abstract}

\tableofcontents
\newpage

% =============================================================================
\section{How to read this document}
% =============================================================================
\begin{itemize}[leftmargin=1.2em]
\item \textbf{Paper references.} ``\Peq{15}'' is equation (15) of the manuscript; Sec.~VI-B and Proposition~1 refer to the
  manuscript. The equation-to-code companion of September 30 (\code{paper/companion_2026-09-30}) covers the quaternion
  conventions, the residual and the gate in detail; this document does not repeat them.
\item \textbf{Code.} Every listing is read from the repository file at the quoted lines, with the file's own line numbers.
  Four files carry all six families: \BL{} (class \code{AdvancedBaseline}: backstepping, LQR, SMC, MPC, NMPC), \RG{}
  (the A5 plant loop and the PID path), and the campaign drivers \code{run_campaign_a5.py} and
  \code{run_campaign_a5_nmpc_corrected.py}. Appendix~\ref{app:hashes} gives their hashes.
\item \textbf{Status tags.} \emph{Standard}: textbook result, cited. \emph{Derived here}: the steps are written out.
  \emph{Verified}: recomputed by \code{tools/build_family_tables.py} from the code itself. \emph{Implementation note}: a
  detail of the code that the paper leaves implicit, with its effect on reported numbers.
\item \textbf{Figures.} Every architecture uses the manuscript's figure style. The left part (reference and attitude error)
  and the right part (torque box, gated residual, allocation, plant) are identical in all six figures; the dashed blue group is
  the only part that changes from family to family.
\end{itemize}

% =============================================================================
\section{What all six families share}\label{sec:common}
% =============================================================================
\subsection{The residual attaches to any nominal torque}
The residual reads only the nominal request $\bm u_0$ and adds a bounded, gated correction to it (Fig.~\ref{fig:interface}).
Proposition~1 assumes only $\bm u_{0,k}\in\mathcal U=\{\bm u:|u_i|\le\ell_i\}$, so it holds for every family whose output
is saturated to the torque box. All six are: LQR, backstepping, SMC and MPC through the shared clip (line~241 of \BL),
NMPC by clipping every candidate (lines 265 and 273), and PID at line~353 of \RG. Exact recovery at $\beta=0$ was certified
by running the augmentor with $\beta=0$ against the classical run: relative difference $0$ in all five cases (LQR, SMC, MPC
and two NMPC missions; \code{certify_beta0_recovery.py}). Backstepping is covered by the 20/20 bit-exact match with the
sealed primary runs; PID has no separate $\beta=0$ run, and for it the recovery rests on Proposition~1 alone.

\begin{figure}[H]\centering
\resizebox{\textwidth}{!}{% How every family attaches to the same gated CMG--PES residual (Sec. 2). Row 1: the six u0 sources and the
% plant path; rows 2-3: the residual path of the manuscript's Fig. 1, unchanged.
\begin{tikzpicture}
  % ---- six u0 sources -----------------------------------------------------------------------------------------
  \foreach \n/\t [count=\i] in {f1/PD+ff (backstepping),f2/PID+ff,f3/LQR,f4/SMC,f5/MPC,f6/NMPC}
    \node[law,minimum height=6.5mm,minimum width=30mm] (\n) at (0,{1.75-0.9*\i}) {\t};
  \ifwithbs
  \node[law,dashed,minimum height=6.5mm,minimum width=30mm,align=center] (f7) at (0,{1.75-0.9*7-0.25}) {Recursive backstepping\\(new, not in paper)};
  \node[fam,fit=(f1)(f7)] (fams) {};
  \else
  \node[fam,fit=(f1)(f6)] (fams) {};
  \fi
  \node[famlbl,anchor=south] at (fams.north) {one family per run (A5)};
  \ifwithbs\def\nfam{7}\else\def\nfam{6}\fi
  \foreach \i in {1,...,\nfam} \draw[thick] (f\i.east) -- (f\i.east-|fams.east);

  % ---- residual path (manuscript Fig. 1) -----------------------------------------------------------------------
  \node[neuro] (enc)  at (4.9,-3.0) {\textbf{Paired currents}\\$\bar{\bm u}=\bm u_0\oslash\bm\ell$\\
        $\bm I^\pm=I_b\bm1+I_g\max(\pm\bar{\bm u},\bm0)$};
  \node[neuro] (cmg)  [right=7mm of enc] {\textbf{Six CMG units}\\$(v_j,r_j)$, no spikes\\$\bm a=\clip_{[0,1]}(G\bm r^+)$};
  \node[neuro] (dec)  [right=7mm of cmg] {\textbf{Decoder}\\$\bm s=\clip_{[-1,1]}(W\bm a)$};
  \node[clas]  (gate) [right=8mm of dec] {\textbf{Saturation-aware gate}\\$\tilde s_i=0$ if $|u_{0,i}|\ge\ell_i$ and $s_iu_{0,i}<0$\\
        else $\tilde s_i=s_i$};
  \node[neuro] (pes)  [below=11mm of cmg.south west,anchor=north west] {\textbf{PES} (online, every step)\\
        $W^+=\Pi_{\mathcal W}\big[(1-\lambda_W)W+\eta\bm\varepsilon\bm a\T\big]$,\ \ $\bm\varepsilon=\bm y^*-\bm s$};
  \node[neuro] (tgt)  [left=7mm of pes] {\textbf{PES teaching target}\\$\bm y^*=\clip_{[-1,1]}(-k_e\bm e_q)$\\from the measured attitude};

  % ---- plant path ----------------------------------------------------------------------------------------------
  \node[clas]  (sel)   at (enc|-0,0) {\textbf{Nominal torque}\\$\bm u_0=\pi_{\rm fam}(q,\bm\omega,t)\in\mathcal U$};
  \node[clas]  (box)   at (gate|-0,0) {\textbf{Bounded torque addition}\\$\bm u=\sat_{\bm\ell}(\bm u_0+\beta\bm\ell\odot\tilde{\bm s})$};
  \node[plant] (plant) [right=8mm of box] {\textbf{Allocation and plant}\\$\bm f=\clip(B^{-1}[T,\bm u\T]\T)$\\MuJoCo X2 / CF2};
  \draw[sig] (fams.east|-sel) -- (sel);
  \draw[sig] (sel) -- node[lbl,above]{$\bm u_0$} (box);
  \draw[sig] (box) -- node[lbl,above]{$\bm u$} (plant);

  % ---- signals ---------------------------------------------------------------------------------------------------
  \draw[sig] (sel) -- node[lbl,right]{$\bm u_0$} (enc);
  \draw[sig] (enc) -- node[lbl,above]{$\bm I$} (cmg);
  \draw[sig] (cmg) -- node[lbl,above]{$\bm a$} (dec);
  \draw[sig] (dec) -- node[lbl,above]{$\bm s$} (gate);
  \draw[sig] (gate) -- node[lbl,right]{$\tilde{\bm s}$} (box);
  \draw[sig] (cmg.south) -- node[lbl,left]{$\bm a$} (cmg.south|-pes.north);
  \draw[sig] ($(pes.north-|dec)+(-3mm,0)$) -- node[lbl,left]{$W$} ($(dec.south)+(-3mm,0)$);
  \draw[sig] ($(dec.south)+(3mm,0)$) -- node[lbl,right]{$\bm s$} ($(pes.north-|dec)+(3mm,0)$);
  \draw[sig] (tgt) -- node[lbl,above]{$\bm y^*$} (pes);
  \coordinate (u0tap) at ($(sel.east)+(6mm,0)$);
  \fill (u0tap) circle (1.1pt);
  \draw[sigd] (u0tap) -- ++(0,-1.25) -| node[lbl,pos=.25,above]{$\bm u_0$ (load test of the gate)} ($(gate.north)+(-12mm,0)$);

  \begin{scope}[on background layer]
    \node[fam,draw=black!55,fit=(enc)(cmg)(dec)(gate)(pes)(tgt),inner sep=9pt] (rp) {};
  \end{scope}
  \node[note,anchor=north east] at (rp.south east) {CMG--PES residual path of the manuscript, unchanged across families};
\end{tikzpicture}
}
\caption{One residual, six nominal controllers. The nominal torque comes from one family per run; the residual path is the
manuscript's Fig.~1 unchanged ($\beta=.5$ on X2, $.2$ on CF2; $\eta=8\times10^{-4}$; $k_e=8$). The gate and the outer saturation
see only $\bm u_0$, so nothing in the residual depends on how $\bm u_0$ was computed.}
\label{fig:interface}
\end{figure}

\subsection{Error coordinates and the design model}\label{sec:design}
All families act on the shortest-path error $\bm e_q=\sgnp(q_{e,0})\bm q_{e,v}$ (\Peq{4}--\Peq{6}) and the rate error
$\bm\omega_e=\bm\omega-\bm\omega_d$, with $\bm\omega_d$ the central-difference reference rate of the companion, Sec.~2.5.
\code{AdvancedBaseline} works in the rotation-vector coordinate
\begin{equation}
\bm e=2\bm e_q,\qquad \|\bm e\|=2\sin\tfrac{\theta_q}{2}=\theta_q+O(\theta_q^3),
\label{eq:e}
\end{equation}
and designs every law for the per-axis double integrator obtained near the reference. From the error kinematics
\Peq{7}--\Peq{8}, $\dot{\bm e}=q_{e,0}\bm\omega_e+\bm e_q\times\bm\omega_e\approx\bm\omega_e$; from Euler's equation \Peq{2} with
the principal moments $\bm J=\bm j_{\rm tune}$ (Table~I of the paper),
\begin{equation}
\dot e_i\approx\omega_{e,i},\qquad \dot\omega_{e,i}=a_i-\dot\omega_{d,i}-\big[(\bm\omega\times\bm J\bm\omega)\oslash\bm J\big]_i,\qquad
a_i=\tau_i/J_i .
\label{eq:design}
\end{equation}
Every family outputs a desired angular acceleration $\bm a$; the torque is $\bm\tau=\bm J\odot\bm a$ plus, optionally, the
gyroscopic term, saturated to the box:
\begin{equation}
\bm u_0=\sat_{\bm\ell}\big(\bm J\odot\bm a+\chi_{\rm ff}\,\bm\omega\times(\bm J\odot\bm\omega)\big),\qquad
\bm a=\chi_{\rm ff}\dot{\bm\omega}_d+\bm a_{\rm fam}(\bm e,\bm\omega_e).
\label{eq:interface}
\end{equation}
\emph{Derived here.} In every run of the paper $\chi_{\rm ff}=0$ (\code{--no-model-ff}): the acceleration feedforward and the
gyroscopic term are off, so each family is a feedback law on $(\bm e,\bm\omega_e)$ scaled by inertia, and the residual is the
only adaptive element. The rate feedforward $\bm\omega_d$ stays on (\code{--rate-ff left-neg}) for all families.

\listing{experiments/nengo_nef_replication_2026-08-29/run_mujoco_nef_baselines.py}{189}{199}{$\bm\omega_d$ and $\dot{\bm\omega}_d$ (the latter is zero when model feedforward is off)}
\listing{experiments/nengo_nef_replication_2026-08-29/run_mujoco_nef_baselines.py}{216}{245}{\code{AdvancedBaseline.torque}: one interface, five families, shared clip at line 241}

\subsection{Torque box and platform constants}
$\ell_i=J_i\nu_i$ with $\bm\nu=(12,12,10)$\,rad/s$^2$, the same for both platforms, so the box limits every family to the
same angular acceleration $|a_i|\le\nu_i$:

\begin{table}[H]\centering\small
\caption{Constants read from the A5 run manifests (\ccode{derived} field).}\label{tab:platform}
\begin{tabular}{@{}lll@{}}\toprule
& X2 & CF2 (1.5$\times$ rotor design)\\\midrule
$\bm J$ (kg\,m$^2$) & \fJXtwo & \fJCFtwo\\
$\bm\ell$ (N\,m) & \fellXtwo & \fellCFtwo\\
$\bm\ell\oslash\bm J$ (rad/s$^2$) & \fnuXtwo & \fnuCFtwo\\
$\rho\bm\omega_n$ (rad/s) & 12, 12, 10.5 & 12, 12, 10.5\\\bottomrule
\end{tabular}\end{table}

Because the box is defined as an acceleration, every gain in this document is the same for X2 and CF2 (the script asserts
it), and the torque differs only by the inertia factor.

\subsection{The A5 settings}
All six families ran with the same plant, missions, residual, gate, $\rho=3$ bandwidth scale and seed; only the family flags
differ:
\listing{experiments/final_robustness_gating_2026-09-20/run_campaign_a5.py}{27}{40}{per-family flags}
\listing{experiments/final_robustness_gating_2026-09-20/run_campaign_a5.py}{49}{56}{common flags: $\rho=3$, rate feedforward on, model feedforward off, gate on in the CMG arm}

\implnote{The NMPC row of the paper comes from \code{run_campaign_a5_nmpc_corrected.py}, which drops the
\code{--baseline-rate-hz 100} hold of the first A5 run so that NMPC, like every other family, updates at every $0.5$-ms step
(Sec.~\ref{sec:nmpc}). The held run is kept on disk and not used.}

\subsection{Small-error equivalence of the six families}\label{sec:equiv}
Near the reference, inside the SMC boundary layer and away from the box, every family reduces to
$a_i=-k_{1,i}e_i-k_{2,i}\omega_{e,i}\,(-k_{I,i}\!\int\! e_i)$, i.e.\ the characteristic polynomial $s^2+k_2s+k_1$ (or
$s^3+k_2s^2+k_1s+k_I$ for PID). Table~\ref{tab:gains} lists the gains each family actually uses; Table~\ref{tab:poles} gives
the resulting poles of the linearized design model \eqref{eq:design}.

\begin{table}[H]\centering\small
\caption{Equivalent small-error gains, read from \ccode{AdvancedBaseline} with the A5 settings (roll\,/\,yaw; pitch equals roll).
$\omega_{\rm eq}=\sqrt{k_1}$, $\zeta_{\rm eq}=k_2/(2\sqrt{k_1})$. ``Box at'' is the rotation angle at which the
proportional term alone reaches the torque box, $\nu_i/k_{1,i}$. \emph{Verified.}}\label{tab:gains}
\begin{tabular}{@{}lllllll@{}}\toprule
Family & $k_1$ (s$^{-2}$) & $k_2$ (s$^{-1}$) & $\omega_{\rm eq}$ (rad/s) & $\zeta_{\rm eq}$ & box at ($^\circ$) & $k_I$ (s$^{-3}$)\\\midrule
\tabinput{generated/family_gains.tex}
\bottomrule\end{tabular}\end{table}

\begin{table}[H]\centering\small
\caption{Closed-loop poles of the linearized design model \eqref{eq:design} with each family's gains. \emph{Verified.}}\label{tab:poles}
\begin{tabular}{@{}lll@{}}\toprule
Family & axis & poles (s$^{-1}$)\\\midrule
\tabinput{generated/family_poles.tex}
\bottomrule\end{tabular}\end{table}

The families therefore differ in damping and in what happens away from the origin: how each one behaves at the box (SMC
uses a bounded reaching term, PID freezes its integrator, NMPC searches around the nominal law), not in bandwidth. The
proportional term of every linear law reaches the box at a rotation error of $5$--$6^\circ$, so Hardstep's steps of tens of
degrees saturate every family, which is where the gate matters.

% =============================================================================
\section{PD with desired-rate feedforward in computed-torque form (primary baseline)}\label{sec:pdff}
In the code this family is called \code{backstepping} (\code{--baseline backstepping}); it is the computed-torque form.\ifwithbs{}
It is not a recursive backstepping design; the recursive design is Sec.~\ref{sec:bstrue}.\fi
% =============================================================================
\archfig{fig_pdff.tex}{PD+ff in computed-torque (backstepping) form, the paper's $\rho=3$ primary baseline. The model
feedforward path (dotted) exists in the code and is switched off in every run, which makes this law identical to \Peq{15}.}{fig:pdff}

\paragraph{Law.} Computed-torque backstepping \cite{krstic1995nonlinear,slotine1991applied}: choose the error dynamics
$\dot{\bm\omega}_e=-k_1\bm e-k_2\bm\omega_e$ and invert \eqref{eq:design},
\begin{equation}
\bm u_0=\sat_{\bm\ell}\Big(\bm J\odot\big(\dot{\bm\omega}_d-k_1\bm e-k_2\bm\omega_e\big)+\bm\omega\times(\bm J\odot\bm\omega)\Big),\qquad
k_1=(\rho\omega_n)^2,\quad k_2=2\zeta\rho\omega_n,\ \zeta=0.95 .
\label{eq:bs}
\end{equation}
With the feedforward on, no saturation and a diagonal inertia, the gyroscopic term and $\dot{\bm\omega}_d$ cancel and the
rate error obeys $\dot{\bm\omega}_e=-k_1\bm e-k_2\bm\omega_e$ exactly; with $\dot{\bm e}\approx\bm\omega_e$ this is
$\ddot e_i+k_2\dot e_i+k_1e_i\approx0$ near the reference, the textbook reason for this form. \emph{Standard.}

\paragraph{Reduction to the paper's PD+ff.} With $\chi_{\rm ff}=0$ and $\bm e=2\bm e_q$,
$\bm u_0=\sat_{\bm\ell}(-2\bm J(\rho\bm\omega_n)^2\odot\bm e_q-2\zeta\bm J\rho\bm\omega_n\odot\bm\omega_e)$, which is \Peq{13}--\Peq{15}
with $\bm\omega_n\mapsto\rho\bm\omega_n$. \emph{Derived here; verified bit-exact in 20/20 cells}
(\code{run_campaign_a5_backstepping.py}, relative deviation $0$ against the sealed primary runs).

\listing{experiments/nengo_nef_replication_2026-08-29/run_mujoco_nef_baselines.py}{138}{141}{gains $k_1=\omega^2$, $k_2=2\zeta\omega$, $\omega=\rho\omega_n$}
\listing{experiments/final_robustness_gating_2026-09-20/run_controllers_gated.py}{79}{86}{the frozen paper baseline: backstepping, $\rho=3$, $\zeta=.95$, \code{left-neg}, no model feedforward}

\ifwithbs
% =============================================================================
\section{Recursive Lyapunov backstepping (new family)}\label{sec:bstrue}
% =============================================================================
The family of Sec.~\ref{sec:pdff} is computed torque: it chooses linear error dynamics and inverts the model, and with
model feedforward off it is exactly PD+ff. This section adds a backstepping controller in the strict sense: a recursive
design with a virtual control on the kinematic subsystem, a cross-coupling term that cancels the interconnection in the
Lyapunov derivative, and a final Lyapunov function that certifies the closed loop
\cite{krstic1995nonlinear,khalil2002nonlinear}. It follows step by step the design used for the longitudinal and adaptive
controllers of the EHEKATL BWB (strict-feedback form, errors $z_1,z_2$, virtual control with the term $-g_1\T z_1$,
$V=\sum\frac12z_i\T z_i$, La Salle--Yoshizawa, adaptive estimates with $\frac12\tilde f\T\Gamma\tilde f$). The plain and the
adaptive versions are implemented in \code{experiments/backstepping_true_2026-10-01/true_backstepping.py} and run through the
unchanged A5 runner. \textbf{They are not part of the manuscript.}

\archfig{fig_bstrue.tex}{Recursive backstepping. Step~1 stabilizes the quaternion kinematics with a virtual rate; step~2 drives
the rate to it with a torque that cancels the gyroscopic term, feeds forward the transported desired-rate derivative, and
contains the coupling term $-c\,z_1$. The adaptive variant adds a projected estimate of an unknown body torque.}{fig:bstrue}

\subsection{Strict-feedback form}
With the passive convention of the paper, $q_e=q_d\otimes q^{-1}$ and its body-frame rate
$\bm\omega_e=\bm\omega-\bar{\bm\omega}_d$, $\bar{\bm\omega}_d=R(q_e)\T\bm\omega_d$ (\Peq{7}--\Peq{8}, with the transport
$R(q_e)\T$ kept), the attitude system is in strict-feedback form:
\begin{equation}
\dot z_1=\tfrac12T(q_e)\,(\bm\omega-\bar{\bm\omega}_d),\qquad
\bm J\dot{\bm\omega}=-\bm\omega\times\bm J\bm\omega+\bm\tau+\bm d,\qquad
z_1=\bm e_q=\sgnp(q_{e,0})\bm q_{e,v},\quad T(q_e)=|q_{e,0}|I+[z_1]_\times ,
\label{eq:sf}
\end{equation}
i.e.\ $x_1=z_1$, $x_2=\bm\omega$, $f_1=-\frac12T\bar{\bm\omega}_d$, $g_1=\frac12T$, $f_2=-\bm J^{-1}(\bm\omega\times\bm J\bm\omega)$,
$g_2=\bm J^{-1}$ in the thesis notation, with $\bm d$ an unknown body torque (zero in the plain design). There is no actuator
state: the rotor forces of the MuJoCo models act without lag, so the thesis's $z_3$ step is not needed. The first identity
follows from \Peq{8}: $\frac{d}{dt}(\sgnp(q_{e,0})\bm q_{e,v})=\frac12(|q_{e,0}|\bm\omega_e+z_1\times\bm\omega_e)$.

\subsection{Step 1: kinematic subsystem}
Take $V_1=2(1-|q_{e,0}|)\ge0$, zero exactly at $q_e=\pm1$. With $\dot q_{e,0}=-\frac12\bm q_{e,v}\T\bm\omega_e$ (\Peq{8}),
\[
\dot V_1=-2\,\sgnp(q_{e,0})\,\dot q_{e,0}=z_1\T\bm\omega_e .
\]
Treat $\bm\omega_e$ as the control of this subsystem and choose the virtual rate $\bm\omega_e^d=-K_1z_1$, $K_1=\operatorname{diag}(k_{1,i})>0$:
if $\bm\omega_e=\bm\omega_e^d$, then $\dot V_1=-z_1\T K_1z_1$.

\subsection{Step 2: rate subsystem}
Define the second error $z_2=\bm\omega_e-\bm\omega_e^d=\bm\omega_e+K_1z_1$ and $V=c\,V_1+\frac12z_2\T\bm Jz_2$, $c>0$. Using
$\bm\omega_e=z_2-K_1z_1$ and $\dot{\bm\omega}_e=\dot{\bm\omega}-\dot{\bar{\bm\omega}}_d$,
\[
\dot V=c\,z_1\T z_2-c\,z_1\T K_1z_1+z_2\T\big(-\bm\omega\times\bm J\bm\omega+\bm\tau+\bm d-\bm J\dot{\bar{\bm\omega}}_d+\bm JK_1\dot z_1\big).
\]
The control
\begin{equation}
\boxed{\ \bm\tau=\bm\omega\times\bm J\bm\omega+\bm J\big(\dot{\bar{\bm\omega}}_d-K_1\dot z_1\big)-c\,z_1-K_2z_2-\hat{\bm d}\ },\qquad
\bm u_0=\sat_{\bm\ell}(\bm\tau),
\label{eq:bstrue}
\end{equation}
with $K_2=\operatorname{diag}(k_{2,i})>0$ and $\dot z_1=\frac12T(q_e)\bm\omega_e$ computed from measured signals, gives
\begin{equation}
\dot V=-c\,z_1\T K_1z_1-z_2\T K_2z_2+z_2\T\tilde{\bm d},\qquad \tilde{\bm d}=\bm d-\hat{\bm d}.
\label{eq:vdot}
\end{equation}
The term $-c\,z_1$ is the backstepping coupling term: it cancels $c\,z_1\T z_2$, the interconnection between the two
subsystems ($-g_1\T z_1$ in the thesis). Without it the design is computed torque.

\paragraph{Plain law} ($\bm d=\hat{\bm d}=\bm0$). $\dot V=-c\,z_1\T K_1z_1-z_2\T K_2z_2<0$ for $(z_1,z_2)\ne0$; by La Salle--Yoshizawa
$z_1,z_2\to0$, hence $q_e\to\pm1$ and $\bm\omega\to\bar{\bm\omega}_d$, and $z_1,z_2\in\mathcal L_2$. Because
$c\,V_1\ge c\|z_1\|^2/2$ near the origin, the convergence is exponential there.

\paragraph{Adaptive law.} For a constant unknown $\bm d$, add $\frac12\tilde{\bm d}\T\Gamma^{-1}\tilde{\bm d}$ to $V$. Since
$\dot{\tilde{\bm d}}=-\dot{\hat{\bm d}}$, the update
\begin{equation}
\dot{\hat{\bm d}}=\operatorname{Proj}_{|\hat d_i|\le D_i}\big(\Gamma z_2\big)
\label{eq:adapt}
\end{equation}
cancels $z_2\T\tilde{\bm d}$ in \eqref{eq:vdot} and leaves $\dot V_a\le-c\,z_1\T K_1z_1-z_2\T K_2z_2$; the projection keeps
$\hat{\bm d}$ bounded without increasing $\dot V_a$ when $|d_i|\le D_i$ \cite{ioannou1996robust,krstic1995nonlinear}. This is the
thesis's $\dot{\tilde f}=-\Gamma^{-1}z$ law with a bound added. Here $\Gamma_i=0.3\cdot2J_i\omega_i^2/k_{1,i}$, which makes the
integral action on $z_1$ ($\Gamma k_1\!\int\!z_1$) equal to the PID family's, and $D_i=0.25\,\ell_i$.

\paragraph{Rate transport.} $\bar{\bm\omega}_d=R(q_e)\T\bm\omega_d$ and, since $q_e$ rotates with body rate $\bm\omega_e$
($\dot R_e=R_e[\bm\omega_e]_\times$),
\[
\dot{\bar{\bm\omega}}_d=R(q_e)\T\dot{\bm\omega}_d-\bm\omega_e\times\bar{\bm\omega}_d ,
\]
with $\bm\omega_d$, $\dot{\bm\omega}_d$ the same central differences as Sec.~\ref{sec:design}. Unlike the A5 families, this law
uses the transport, the gyroscopic term and $\dot{\bm\omega}_d$ by construction; the \code{--no-model-ff} flag does not apply.

\subsection{Gain design: same linearization as PD+ff}
About $z=0$, \eqref{eq:sf}--\eqref{eq:bstrue} give, per axis, $\ddot z_{1,i}+(k_{1,i}/2+k_{2,i}/J_i)\dot z_{1,i}+(c+k_{1,i}k_{2,i})/(2J_i)\,z_{1,i}=0$.
Writing $k_{1,i}=2\lambda_{1,i}$, $k_{2,i}=J_i\lambda_{2,i}$ and imposing
\[
\lambda_{1,i}+\lambda_{2,i}=2\zeta\omega_i,\qquad \lambda_{1,i}\lambda_{2,i}+\frac{c}{2J_i}=\omega_i^2,\qquad \omega_i=\rho\omega_{n,i},\ \zeta=0.95,
\]
makes the small-error poles exactly those of the paper's PD+ff ($-11.40\pm3.75j$ on roll). Any difference between the two
families is then due to the backstepping structure, not to bandwidth or damping. $V_1$ is a scalar, so one $c$ serves all
axes; real $\lambda$'s need $c\ge2J_i\omega_i^2(1-\zeta^2)$ on every axis, and the smallest such value is used
($\lambda_1\le\lambda_2$: the outer kinematic loop is the slower one). For X2: $c=\fbsC$\,N\,m,
$\bm k_1=(\fbsKone)$\,s$^{-1}$, $\bm k_2=(\fbsKtwo)$\,N\,m\,s.

\subsection{Verification}
\code{verify_true_backstepping.py} checks the derivation on the rigid-body model (\code{results/verification.json}). \emph{Verified}:
\begin{itemize}[leftmargin=1.2em]
\item $\dot z_1$ and $\dot{\bar{\bm\omega}}_d$ in closed form against central finite differences along rigid-body trajectories on Full3,
  Hard45 and Yaw360: worst absolute error \fbsKinErr{} (signals of order $0.1$--$1$);
\item closed loop, plain law, no torque limit, Full3 and Hard45 from a $\approx20^\circ$ initial error: $V$ decreases monotonically to
  $10^{-12}$, and the finite-difference $\dot V$ equals $-c\,z_1\T K_1z_1-z_2\T K_2z_2$ within \fbsVdotErr{} (relative, step-size limited);
\item adaptive law with a constant unknown torque of $20\%$ of the box: $\hat{\bm d}$ reaches $\bm d$ within \fbsAdaptErr{} after $20$\,s;
\item linearized poles equal the PD+ff poles to \fbsPoleErr.
\end{itemize}

\listing{experiments/backstepping_true_2026-10-01/true_backstepping.py}{44}{55}{gain design: smallest common $c$, $\lambda_1\le\lambda_2$}
\listing{experiments/backstepping_true_2026-10-01/true_backstepping.py}{84}{114}{$z_1$, $z_2$, $\dot z_1$, transport, law \eqref{eq:bstrue}, projected adaptation \eqref{eq:adapt}}
\listing{experiments/backstepping_true_2026-10-01/run_true_backstepping.py}{30}{42}{the family enters the unchanged A5 runner}

\implnote{The Lyapunov certificate holds for the model \eqref{eq:sf} without saturation. In the plant the torque box, rotor
clipping, the non-diagonal X2 inertia (the law uses the principal moments, as every family does) and the residual add terms
that the certificate does not cover; the sign selection makes $V_1$ non-smooth at $q_{e,0}=0$, the same local scope as the
paper. The output is saturated to the box, so Proposition~1 applies unchanged.}

\subsection{Results: recursive backstepping with and without the residual}
Same protocol as A5 (Table~\ref{tab:results}): 45-s missions, $\rho=3$, noise-free, seed 1, gate on, the residual unchanged. Run by \code{experiments/backstepping_true_2026-10-01/run_campaign_true_bs.py}; numbers computed from the run summaries by \code{tools/build_family_tables.py}.
\begin{table}[H]\centering\small
\caption{Reduction of mean quaternion-angle error by the gated CMG--PES residual on the recursive backstepping laws (\%, positive is better; 5-mission = reduction of geometric means; last column: baseline geometric mean, mrad).}\label{tab:bsres}
\begin{tabular}{@{}llrrrrrrr@{}}\toprule
Family & Platform & Roll & Full3 & Hard45 & Hardstep & Yaw360 & 5-mission & base (mrad)\\\midrule
Recursive backstepping & X2 & -114.8 & +7.2 & +11.8 & -2.9 & +43.5 & \textbf{-0.4} & 3.30\\
 & CF2 & -1416.5 & -18.2 & -38.1 & -0.8 & -1467.9 & \textbf{-230.0} & 0.59\\\addlinespace
Adaptive recursive backstepping & X2 & -114.7 & +6.5 & +11.7 & -0.8 & +30.9 & \textbf{-4.3} & 2.63\\
 & CF2 & -669.0 & -9.6 & -19.6 & -0.4 & -2494.5 & \textbf{-204.7} & 0.88\\\addlinespace
\bottomrule\end{tabular}\end{table}
\begin{table}[H]\centering\small
\caption{Nominal laws alone (no residual): error reduction of the recursive designs relative to the computed-torque PD+ff (the paper's primary baseline), \% (positive: the recursive law tracks better). The linearizations are identical, so these differences come from the backstepping structure.}\label{tab:bsvsct}
\begin{tabular}{@{}llrrrrrr@{}}\toprule
Platform & law & Roll & Full3 & Hard45 & Hardstep & Yaw360 & 5-mission\\\midrule
X2 & recursive & +29.6 & +30.5 & +31.0 & -2.4 & +41.5 & \textbf{+27.4}\\
X2 & adaptive & +36.5 & +30.8 & +29.1 & -7.9 & +80.7 & \textbf{+42.2}\\
CF2 & recursive & +95.0 & +64.8 & +73.3 & -0.5 & +97.9 & \textbf{+84.2}\\
CF2 & adaptive & +91.2 & +47.5 & +58.9 & -695.8 & +99.6 & \textbf{+76.7}\\
\bottomrule\end{tabular}\end{table}


\paragraph{What the two tables show.}
\begin{itemize}[leftmargin=1.2em]
\item \textbf{The recursive law is a much better nominal controller} (Table~\ref{tab:bsvsct}): with the same linearization it
  reduces the computed-torque PD+ff error on four of five missions on both platforms, most strongly on CF2. Its law contains
  $\bm J\dot{\bar{\bm\omega}}_d$ and $\bm\omega\times\bm J\bm\omega$, which every A5 family ran without (\code{--no-model-ff}). The
  computed-torque law \emph{with} model feedforward was not run, so this campaign does not separate the share of the gain due to the
  acceleration feedforward from the share due to the recursive structure (cross term and rate transport). That run is the
  natural next control.
\item \textbf{The unchanged residual does not help it} (Table~\ref{tab:bsres}). It still reduces the error on X2 Full3, Hard45
  and Yaw360, but where the recursive law's error is already very small (Roll; CF2 Roll and Yaw360, a few hundredths of a mrad) it
  raises the error to tenths of a mrad, worse in absolute terms than the residual on the paper's baseline. The residual's
  $\beta$, $\eta$ and $k_e$ were set for a lagging PD+ff; on a nominal law that already cancels the dominant error, the residual's
  own correction becomes the error. Retuning ($\beta$ in particular) was not tested. This is a limit of the unchanged residual,
  not of Proposition~1, which still holds: the output stays in the box and $\beta=0$ recovers the law.
\item \textbf{The adaptive variant fails on CF2 Hardstep} under persistent rotor saturation (79\% of samples): the projected estimate
  keeps integrating $z_2$ while the torque it requests cannot be delivered, a windup through the adaptation law that the
  Lyapunov argument (which assumes no saturation) does not exclude. Freezing the adaptation when the box or a rotor saturates, the
  rule the PID family uses, is the standard fix and was not run. On the other missions the adaptive and plain laws are close.
\end{itemize}
\fi % withbs

% =============================================================================
\section{PID with rate feedforward and rotor-level anti-windup}\label{sec:pid}
% =============================================================================
\archfig{fig_pid.tex}{PID+ff. The integrator runs on $\bm e_q$ and is frozen per axis whenever the torque box clipped that axis,
and on all axes whenever any rotor saturated at the previous step (dashed paths).}{fig:pid}

\paragraph{Law.} The PID path bypasses \code{AdvancedBaseline} (\code{--baseline pd}) and adds an integral term to \Peq{15}:
\begin{equation}
\bm u_0=\sat_{\bm\ell}\big(-\bm k_q\odot\bm e_q-\bm k_\omega\odot\bm\omega_e-\bm k_i\odot\bm z\big),\qquad
\bm k_q=2\bm J\odot(\rho\bm\omega_n)^2,\ \bm k_\omega=2\zeta\bm J\odot\rho\bm\omega_n,\ \bm k_i=0.3\,\bm k_q .
\label{eq:pid}
\end{equation}
\paragraph{Conditional integration.} With $h=5\times10^{-4}$\,s and the clip test of line 357,
\begin{equation}
z_{i,k+1}=z_{i,k}+h\,e_{q,i,k}\,m_{i,k},\qquad
m_{i,k}=\mathbf 1\big[u_{0,i,k}=\tau_{{\rm u},i,k}\big]\cdot\mathbf 1\big[\text{no rotor saturated at }k-1\big].
\label{eq:aw}
\end{equation}
This is conditional integration (integrator clamping) \cite{astrom2006advancedpid}, applied at the rotor level as well as at the
torque box: a rotor at its force limit means the realized torque differs from the command even if the box is inactive.
\emph{Standard; the rotor-level rule is this study's choice.}

\paragraph{Loop.} In the coordinate $\bm e$ the gains per unit inertia are $k_1=(\rho\omega_n)^2$, $k_2=2\zeta\rho\omega_n$ and
$k_I=0.3(\rho\omega_n)^2$ (Table~\ref{tab:gains}). The Routh condition for $s^3+k_2s^2+k_1s+k_I$ is $k_2k_1>k_I$, met with
a ratio of at least \fpidRouthMargin. The integral adds a slow real pole at $-0.32$\,s$^{-1}$ (time constant $\approx3$\,s,
Table~\ref{tab:poles}): the integrator is weak, which the A5 report discloses as ``undertuned''. \emph{Derived here; verified.}

\listing{experiments/final_robustness_gating_2026-09-20/run_controllers_gated.py}{253}{260}{gains; PID uses the plain PD path}
\listing{experiments/final_robustness_gating_2026-09-20/run_controllers_gated.py}{347}{361}{\eqref{eq:pid} and the conditional integrator \eqref{eq:aw}}
\listing{experiments/final_robustness_gating_2026-09-20/run_controllers_gated.py}{475}{475}{rotor-saturation flag for the next step}

\implnote{The integrator acts on $\bm e_q$ (not $\bm e=2\bm e_q$), so $\bm k_i=0.3\bm k_q$ means
$k_I=0.3(\rho\omega_n)^2$ in Table~\ref{tab:gains}'s coordinate. With \code{--ki-factor 0} this path is \Peq{15} exactly.}

% =============================================================================
\section{LQR}\label{sec:lqr}
% =============================================================================
\archfig{fig_lqr.tex}{LQR. The gain is computed once from the closed-form solution of the continuous algebraic Riccati
equation (CARE) for the per-axis double integrator; online it is a fixed state feedback.}{fig:lqr}

\paragraph{Design problem.} For each axis, $\bm x=[e,\omega_e]\T$, $\dot{\bm x}=A\bm x+Ba$ with
$A=\left[\begin{smallmatrix}0&1\\0&0\end{smallmatrix}\right]$, $B=[0,1]\T$, minimize
$\int_0^\infty(q_e e^2+ra^2)\,dt$ \cite{andersonmoore1990optimal}.

\paragraph{Closed-form CARE.} With $P=\left[\begin{smallmatrix}p_1&p_2\\p_2&p_3\end{smallmatrix}\right]$, the CARE
$A\T P+PA-PBr^{-1}B\T P+Q=0$ reads, entry by entry,
\[
(1,1):\ q_e-\frac{p_2^2}{r}=0,\qquad (1,2):\ p_1-\frac{p_2p_3}{r}=0,\qquad (2,2):\ 2p_2-\frac{p_3^2}{r}=0 .
\]
The positive-definite solution has $p_2=\sqrt{q_er}$, $p_3=\sqrt{2p_2r}$, $p_1=p_2p_3/r$, so
\begin{equation}
K=r^{-1}B\T P=\Big[\sqrt{q_e/r},\ \sqrt{2\sqrt{q_e/r}}\Big]\quad\xrightarrow{\ q_e/r=\omega^4\ }\quad K=[\omega^2,\ \sqrt2\,\omega],
\label{eq:lqr}
\end{equation}
a closed loop $s^2+\sqrt2\omega s+\omega^2$ with $\zeta=1/\sqrt2$ at the same $\omega=\rho\omega_n$ as the other families.
\emph{Derived here. Verified}: $P$ from these formulas leaves a CARE residual of \flqrCareRelErr{} (relative), its smallest
eigenvalue is \flqrPminEig{} $>0$, and $B\T P$ equals the code's gain.

\listing{experiments/nengo_nef_replication_2026-08-29/run_mujoco_nef_baselines.py}{132}{137}{$K=[\omega^2,\sqrt2\omega]$}
\listing{experiments/nengo_nef_replication_2026-08-29/run_mujoco_nef_baselines.py}{232}{233}{the law: $\bm a=\dot{\bm\omega}_d-k_1\bm e-k_2\bm\omega_e$}

% =============================================================================
\section{Sliding-mode control}\label{sec:smc}
% =============================================================================
\archfig{fig_smc.tex}{SMC. A linear sliding variable and a boundary-layer reaching law; inside the layer the law is a critically
damped PD, outside it the reaching term has a fixed magnitude of half the torque box.}{fig:smc}

\paragraph{Law.} Sliding variable and reaching law \cite{utkin1992sliding,slotine1991applied}, per axis:
\begin{equation}
\sigma_i=\omega_{e,i}+\lambda_ie_i,\qquad
a_i=\dot\omega_{d,i}-\lambda_i\omega_{e,i}-\eta_i\,\sat\!\big(\sigma_i/\phi_i\big),\qquad
\lambda_i=\rho\omega_{n,i},\ \ \eta_i=0.5\,\ell_i/J_i,\ \ \phi_i=\eta_i/\lambda_i .
\label{eq:smc}
\end{equation}
With the A5 setting, $\bm\eta=(\fsmcEtaXtwo)$\,rad/s$^2$ and $\bm\phi=(\fsmcPhiXtwo)$\,rad/s on both platforms.

\paragraph{Inside the boundary layer} ($|\sigma_i|<\phi_i$): $a_i=-\lambda_i\omega_{e,i}-(\eta_i/\phi_i)(\omega_{e,i}+\lambda_ie_i)$, and
$\eta_i/\phi_i=\lambda_i$, so
\[
a_i=-\lambda_i^2e_i-2\lambda_i\omega_{e,i}:\qquad k_1=\lambda^2,\quad k_2=2\lambda,\quad \zeta=1,
\]
a double pole at $-\rho\omega_n$ (Table~\ref{tab:poles}). \emph{Derived here; verified from the code's
\code{lam}, \code{eta}, \code{phi}.}

\paragraph{Outside the layer.} On the design model \eqref{eq:design} with model feedforward off,
$\dot\sigma_i=\dot\omega_{e,i}+\lambda_i\dot e_i=a_i-\dot\omega_{d,i}+\lambda_i\omega_{e,i}=-\eta_i\operatorname{sgn}(\sigma_i)-\dot\omega_{d,i}$,
hence
\[
\sigma_i\dot\sigma_i\le-(\eta_i-|\dot\omega_{d,i}|)\,|\sigma_i| ,
\]
the reaching condition, whenever the reference acceleration stays below $\eta_i$ ($6$\,rad/s$^2$ on roll and pitch). The layer
is then reached in finite time and the error decays inside it; the boundary layer trades exact sliding for a continuous,
chattering-free command. Because $J_i\eta_i=\ell_i/2$, the reaching term alone uses half the torque box, leaving the
other half for $-\lambda\bm\omega_e$. \emph{Standard argument, specialized here.}

\listing{experiments/nengo_nef_replication_2026-08-29/run_mujoco_nef_baselines.py}{142}{149}{$\lambda=\omega$, $\eta=0.5\,\bm\ell\oslash\bm J$, $\phi=\eta/\omega$}
\listing{experiments/nengo_nef_replication_2026-08-29/run_mujoco_nef_baselines.py}{234}{236}{\eqref{eq:smc}}

% =============================================================================
\section{Model predictive control (unconstrained, finite horizon)}\label{sec:mpc}
% =============================================================================
\archfig{fig_mpc.tex}{MPC. The finite-horizon quadratic program has no constraints, so its solution is linear in the state; the
first move of the optimal sequence is a fixed gain $K_0$, computed once and applied at every plant step.}{fig:mpc}

\paragraph{Prediction model.} Per axis, the exact zero-order-hold discretization of the double integrator with sample time
$T=20$\,ms,
\[
\bm x_{k+1}=A_d\bm x_k+B_da_k,\qquad A_d=\begin{bmatrix}1&T\\0&1\end{bmatrix},\quad B_d=\begin{bmatrix}T^2/2\\T\end{bmatrix}.
\]
\paragraph{Condensed problem.} Stacking $N=120$ predictions, $X=S_x\bm x_0+S_uU$ with
$S_x=[A_d;A_d^2;\dots;A_d^N]$ and the block-Toeplitz $S_u$ with blocks $A_d^{k-j}B_d$ ($j\le k$). The cost
$\sum_{k=1}^{N}\bm x_k\T Q\bm x_k+\sum_{k=0}^{N-1}Ra_k^2$ with $Q=\operatorname{diag}(\omega^4,q_\omega\omega^2)$, $q_\omega=0$, $R=1$ becomes
\[
U\T HU+2\bm x_0\T F\T U+\text{const},\qquad H=S_u\T\bar QS_u+\bar R,\quad F=S_u\T\bar QS_x,
\]
with $\bar Q=I_N\otimes Q$, $\bar R=I_N\otimes R$. Without constraints the minimizer is $U^\star=-H^{-1}F\bm x_0$, and the
receding-horizon law applies its first element, $a_0=-K_0\bm x_0$ with $K_0=[H^{-1}F]_{1,:}$
\cite{borrelli2017predictive,rawlings2017mpc}. \emph{Standard.}

\paragraph{What the horizon does.} Table~\ref{tab:mpc} recomputes $K_0$ with the code's own function for growing $N$. From
$N\approx30$ ($0.6$\,s, several closed-loop time constants) the first-move gain no longer changes: at the A5 horizon it equals
the infinite-horizon discrete LQR gain obtained by iterating the Riccati recursion to its fixed point to a relative
\fmpcVsDare{} ($K_0=[\fmpcKone,\ \fmpcKtwo]$ on roll; DARE residual \fdareResidual). The MPC row is therefore a discrete LQR
designed at $T=20$\,ms. That sample time, not the horizon, is why its proportional gain is $16\%$ below the continuous LQR
of Sec.~\ref{sec:lqr} ($121.5$ vs.\ $144$\,s$^{-2}$). \emph{Verified.}

\begin{table}[H]\centering\small
\caption{First-move MPC gain on the roll axis ($\omega=12$\,rad/s) versus horizon, from \ccode{AdvancedBaseline._mpc_first_gain};
the last two columns divide by the continuous LQR gain $[\omega^2,\sqrt2\omega]$. A5 uses $N=120$.}\label{tab:mpc}
\begin{tabular}{@{}rrrrrr@{}}\toprule
$N$ & $NT$ (s) & $K_{0,1}$ & $K_{0,2}$ & $K_{0,1}/\omega^2$ & $K_{0,2}/(\sqrt2\omega)$\\\midrule
\tabinput{generated/mpc_horizon.tex}
\bottomrule\end{tabular}\end{table}

\listing{experiments/nengo_nef_replication_2026-08-29/run_mujoco_nef_baselines.py}{167}{187}{condensed matrices $S_x,S_u,H,F$ and the first-move gain}
\listing{experiments/nengo_nef_replication_2026-08-29/run_mujoco_nef_baselines.py}{237}{238}{the online law $\bm a=\dot{\bm\omega}_d-K_0[\bm e;\bm\omega_e]$}

\implnote{The gain is designed for a 20-ms move but applied as a continuous-time state feedback at every 0.5-ms step (A5 ran
every family at full rate). This is a common simplification; it is the reason the MPC row behaves like a softer LQR
($\zeta_{\rm eq}=0.71$, $\omega_{\rm eq}=11.0$\,rad/s) and not like a sampled controller. ``MPC'' here means an unconstrained
receding-horizon gain; the torque box is enforced by the shared clip, not inside the optimization.}

% =============================================================================
\section{Nonlinear MPC (grid search around the backstepping law)}\label{sec:nmpc}
% =============================================================================
\archfig{fig_nmpc.tex}{NMPC. At every plant step, 27 candidate offsets are added to the nominal backstepping torque; each
candidate is scored by an RK4 rollout of the closed loop on the rigid-body model over $0.3$\,s, and the first torque of the
cheapest candidate is applied.}{fig:nmpc}

\paragraph{Candidates.} The nominal law is \eqref{eq:bs} ($\zeta=0.95$). A grid of $3^3=\fnmpcCandidates$ constant offsets
$\bm d_c\in\{-1,0,1\}^3\odot0.1\bm\ell$ is added to it; $\bm d_c=\bm0$ is one of them (\fnmpcHasZero), so the nominal law is always
a candidate.

\paragraph{Prediction (``policy'' mode).} From the measured state, each candidate is rolled out for $N=5$ steps of
$\Delta=60$\,ms with the rigid-body model $\dot Q=\frac12Q\otimes\tilde{\bm\omega}$,
$\dot{\bm\omega}=(\bm\tau-\bm\omega\times\bm J\bm\omega)\oslash\bm J$, integrated by classical RK4 \cite{hairer1993ode}. At
$k=0$ the candidate's first torque $\bm\tau^{(c)}_0=\sat_{\bm\ell}(\bm\tau_{\rm nom}+\bm d_c)$ is used; at $k\ge1$ the nominal law is
re-evaluated on the predicted state and the same offset added,
$\bm\tau^{(c)}_k=\sat_{\bm\ell}(\bm J\odot\bm a(\hat q_k,\hat{\bm\omega}_k,t+k\Delta)+\bm d_c)$. The prediction thus models the
closed loop that will actually run, not a constant torque.

\paragraph{Cost and selection.} With $\bm e_k,\bm\omega_{e,k}$ the predicted errors against the reference at $t+(k+1)\Delta$,
\begin{equation}
J_c=\sum_{k=0}^{N-1}\Big[w_k\big(\textstyle\sum_i\omega_i^4e_{k,i}^2+2\sum_i\omega_i^2\omega_{e,k,i}^2\big)+\big\|\bm\tau^{(c)}_k\oslash\bm J\big\|^2\Big],
\qquad w_k=\begin{cases}1.5,&k=N-1\\1,&\text{otherwise}\end{cases}
\label{eq:nmpc}
\end{equation}
and $\bm u_0=\bm\tau^{(c^\star)}_0$, $c^\star=\arg\min_cJ_c$. The angle weight $\omega^4$ is the LQR weight $q_e/r$ of
Sec.~\ref{sec:lqr}; the effort term is the squared commanded acceleration. This is a sampling-based (grid) receding-horizon
controller in the sense of \cite{rawlings2017mpc}, not a gradient-based NMPC solver. \emph{Derived here from the code.}

\listing{experiments/nengo_nef_replication_2026-08-29/run_mujoco_nef_baselines.py}{152}{159}{candidate grid and cost weights}
\listing{experiments/nengo_nef_replication_2026-08-29/run_mujoco_nef_baselines.py}{201}{214}{rigid-body model and RK4 step (MuJoCo's body-to-world convention)}
\listing{experiments/nengo_nef_replication_2026-08-29/run_mujoco_nef_baselines.py}{247}{289}{\code{_nmpc}: rollout, cost \eqref{eq:nmpc}, selection, one-period self-check}
\listing{experiments/final_robustness_gating_2026-09-20/run_campaign_a5_nmpc_corrected.py}{25}{26}{corrected A5 settings: no update hold}

\implnote{(i) The rollout predicts the nominal family only; it does not know the residual that will be added. (ii) The first A5
NMPC run updated at 100\,Hz and held the torque between updates, a leftover from an earlier exploratory setting. The audit
(\code{AUDIT_NMPC_HOLD_RATE_2026-09-25.md}) showed that the hold, not the predictive machinery, caused that row's anomaly:
removing it restored NMPC, and imposing the same hold on MPC broke MPC identically. The paper's NMPC row is the full-rate re-run.
(iii) The class logs a one-period-ahead self-check of its own prediction (\code{nmpc_pred_err_mean_deg}); it was used to
diagnose (ii), not as a tuning signal.}

% =============================================================================
\section{Results across families (Sec.~VI-B)}\label{sec:results}
% =============================================================================
\begin{table}[H]\centering\small
\caption{Reduction of mean quaternion-angle error by the gated CMG--PES residual on each nominal family, noise-free, $\rho=3$,
45-s missions (\% , positive is better; ``5-mission'' is the reduction of geometric means). Read from
\ccode{results/a5_numbers.json}, which \ccode{analyze_a5.py} computes from the run summaries; NMPC from the corrected
full-rate run. Last column: baseline geometric mean (mrad).}\label{tab:results}
\begin{tabular}{@{}llrrrrrrr@{}}\toprule
Family & Platform & Roll & Full3 & Hard45 & Hardstep & Yaw360 & 5-mission & base (mrad)\\\midrule
\tabinput{generated/family_results.tex}
\bottomrule\end{tabular}\end{table}

These are the numbers the manuscript reports (X2/CF2: backstepping $37.3/22.4$, PID+ff $36.7/22.4$, LQR $31.1/22.4$, SMC
$43.8/26.6$, MPC $33.4/22.7$, NMPC $31.9/21.7$), and the claim verifier checks them (\code{reproduce/verify_acc2027_claims.py}).
What the table shows, without further interpretation:
\begin{itemize}[leftmargin=1.2em]
\item The five-mission reduction stays within $31$--$44\%$ on X2 and $22$--$27\%$ on CF2 across families whose baseline errors
  differ by up to $1.7\times$ (X2 geometric means $3.7$--$6.3$\,mrad), with the residual unchanged.
\item The reduction is largest on Roll and Yaw360 and smallest on Hardstep, whose step changes of tens of degrees drive every
  linear family far past the $5$--$6^\circ$ at which its proportional term reaches the box (Table~\ref{tab:gains}).
\item SMC is the only family with a clearly positive Hardstep result ($+34.2\%$ X2, $+12.2\%$ CF2). SMC is also the only
  family whose law is itself bounded below the box away from the origin (Sec.~\ref{sec:smc}); A5 did not isolate whether that
  is the cause, so it is recorded here as an open observation.
\end{itemize}

% =============================================================================
\section{Scope of the family comparison}
% =============================================================================
\begin{itemize}[leftmargin=1.2em]
\item The families are \emph{attachment tests}, not tuned competitors: all use the same bandwidth scale $\rho=3$, model
  feedforward off, and the A5 settings above. No family was retuned per mission or per platform, and the residual
  ($\beta$, $\eta$, $k_e$) is the same for all.
\item MPC is unconstrained and realized as its first-move gain; NMPC is a 27-point grid search around backstepping. Neither is a
  constrained optimal controller; the torque box acts through the shared saturation.
\item PID's integrator is weak (slow pole at $-0.32$\,s$^{-1}$) and is disclosed as undertuned in the A5 report.
\item All gains are designed on the per-axis double integrator \eqref{eq:design} with principal moments; the X2 inertia tensor is
  not diagonal, so the poles of Table~\ref{tab:poles} are design values, not measured closed-loop poles.
\end{itemize}

\appendix
\section{Source files and hashes}\label{app:hashes}
SHA-256 (first 16 hex digits) of each cited file, compared with the \code{source_sha256} field of one A5 run per family and
platform (Roll, classical arm; NMPC from the corrected run): \fhashMatch/\fhashChecked{} recorded hashes are identical.
\begin{table}[H]\centering\small
\begin{tabular}{@{}p{0.55\linewidth}ll@{}}\toprule
File & SHA-256 & vs.\ A5 manifests\\\midrule
\tabinput{generated/family_hashes.tex}
\bottomrule\end{tabular}\end{table}

\bibliographystyle{IEEEtran}
\bibliography{../v15_2026-09-29_corrected/references_ACC2027,family_refs}
\end{document}
"
\newif\ifwithbs \withbstrue
\ifdefined\releasebuild \withbsfalse \fi
% the manuscript's own figure style, plus the family extensions
% Shared TikZ style for all architecture figures. Grayscale-safe.
\tikzset{
  blk/.style   ={draw,rounded corners=2pt,align=center,font=\footnotesize,
                 minimum height=8mm,inner sep=3pt,fill=white},
  clas/.style  ={blk,fill=black!4},
  neuro/.style ={blk,fill=black!12,very thick},
  plant/.style ={blk,fill=black!20},
  meas/.style  ={blk,dashed},
  sum/.style   ={draw,circle,inner sep=0pt,minimum size=4.5mm,font=\scriptsize},
  sig/.style   ={-{Latex[length=2mm]},thick},
  sigd/.style  ={-{Latex[length=2mm]},thick,dashed},
  lbl/.style   ={font=\scriptsize,inner sep=1.5pt,fill=white},
  grp/.style   ={draw,dashed,rounded corners=3pt,inner sep=6pt},
  note/.style   ={font=\scriptsize\itshape,align=left},
}

% Extra styles for the controller-family figures; the base styles (blk, clas, neuro, plant, sig, ...) are the
% manuscript's own architectures/tikz/arch_style.tex, so every figure matches the paper's Fig. 1.
\tikzset{
  law/.style    ={blk,fill=blue!7,draw=blue!45!black,thick},
  offl/.style   ={blk,dashed,fill=black!2},
  fam/.style    ={draw=blue!45!black,dashed,rounded corners=3pt,inner sep=7pt},
  famlbl/.style ={font=\scriptsize\itshape,text=blue!45!black,inner sep=1pt},
  off/.style    ={-{Latex[length=2mm]},thick,dotted,gray},
  every node/.append style={font=\footnotesize},
}
% Common left half: reference generator and attitude error (identical for every family).
\newcommand{\commonleft}{%
  \node[clas] (ref) at (0,0) {\textbf{Mission reference}\\$q_d(t)$\\[1pt]
        $\bm\omega_d$: central difference};
  \node[clas,right=7mm of ref] (err) {\textbf{Attitude error}\\$q_e=q_d\otimes q^{-1}$\\
        $\bm e=2\,\sgnp(q_{e,0})\,\bm q_{e,v}$\\$\bm\omega_e=\bm\omega-\bm\omega_d$};
  \draw[sig] (ref) -- node[lbl,above]{$q_d,\bm\omega_d$} (err);
}
% Common right half and feedback of every family figure. Needs nodes (err) and (law-out) placed by the figure,
% and the coordinate (satx) giving the x position of the torque box.
\newcommand{\commonloop}[2]{% #1: label on u0 arrow source node, #2: text inside the torque-box node
  \node[clas] (sat) at (satx) {#2};
  \node[neuro,below=13mm of sat] (res) {\textbf{Gated CMG--PES residual}\\(Fig.~\ref{fig:interface})\\
        $\bm u=\sat_{\bm\ell}(\bm u_0+\beta\bm\ell\odot\tilde{\bm s})$};
  \node[plant,anchor=east] (plant) at ($(res.west)+(-12mm,0)$) {\textbf{Allocation and MuJoCo plant}\\
        $\bm f=\clip_{[\bm0,\bm f_{\max}]}(B^{-1}[T,\bm u\T]\T)$\\X2 / CF2 rigid body};
  \draw[sig] (#1) -- (sat);
  \draw[sig] (sat) -- node[lbl,right]{$\bm u_0\in\mathcal U$} (res);
  \draw[sig] (res) -- node[lbl,above]{$\bm u$} (plant);
  \draw[sig] (plant.west) -- ++(-4mm,0) coordinate (fb1) -| node[lbl,pos=.75,left]{measured $q,\ \bm\omega$} (err.south);
}

\newcommand{\flqrCareRelErr}{2.7e-18}
\newcommand{\flqrPminEig}{7.39}
\newcommand{\fsmcEtaXtwo}{6.00, 6.00, 5.00}
\newcommand{\fsmcPhiXtwo}{0.500, 0.500, 0.476}
\newcommand{\fsmcEtaCFtwo}{6.00, 6.00, 5.00}
\newcommand{\fsmcPhiCFtwo}{0.500, 0.500, 0.476}
\newcommand{\fpidRouthMargin}{66.5}
\newcommand{\fdareResidual}{9.5e-14}
\newcommand{\fmpcDareKone}{121.55}
\newcommand{\fmpcDareKtwo}{15.592}
\newcommand{\fmpcKone}{121.55}
\newcommand{\fmpcKtwo}{15.592}
\newcommand{\fmpcVsDare}{4.6e-10}
\newcommand{\fnmpcCandidates}{27}
\newcommand{\fnmpcHasZero}{yes}
\newcommand{\fhashMatch}{60}
\newcommand{\fhashChecked}{60}
\newcommand{\fbsC}{1.705}
\newcommand{\fbsKone}{22.80, 16.69, 10.40}
\newcommand{\fbsKtwo}{0.692, 0.527, 0.375}
\newcommand{\fbsKinErr}{$5.6\times10^{-7}$}
\newcommand{\fbsVdotErr}{0.22\%}
\newcommand{\fbsAdaptErr}{0.25\%}
\newcommand{\fbsPoleErr}{$2.7\times10^{-15}$}
\newcommand{\fbsGeotrueXtwo}{-0.4}
\newcommand{\fbsGeotrueCFtwo}{-230.0}
\newcommand{\fbsGeoadaptiveXtwo}{-4.3}
\newcommand{\fbsGeoadaptiveCFtwo}{-204.7}
\newcommand{\fbsCampaignComplete}{yes}
\newcommand{\fJXtwo}{0.06071, 0.03647, 0.02541}
\newcommand{\fellXtwo}{0.7285, 0.4376, 0.2541}
\newcommand{\fnuXtwo}{12.00, 12.00, 10.00}
\newcommand{\fJCFtwo}{2.395e-05, 2.395e-05, 3.235e-05}
\newcommand{\fellCFtwo}{0.0002874, 0.0002874, 0.0003235}
\newcommand{\fnuCFtwo}{12.00, 12.00, 10.00}


\newcommand{\BL}{\code{experiments/nengo_nef_replication_2026-08-29/run_mujoco_nef_baselines.py}}
\newcommand{\RG}{\code{experiments/final_robustness_gating_2026-09-20/run_controllers_gated.py}}
\newcommand{\archfig}[3]{% file, caption, label
  \begin{figure}[H]\centering\resizebox{\textwidth}{!}{\input{figures/#1}}\caption{#2}\label{#3}\end{figure}}

\title{Online PES Decoder Learning in Neuromorphic Residual UAV Attitude Control\\[4pt]
\large The nominal controller families of Sec.~VI-B\ifwithbs{} and a recursive backstepping design\fi:\\ equations, derivations, code, and architectures\\
\normalsize companion to the ACC~2027 manuscript (version of September 30, 2026)}
\author{Gandhi Contreras, Steven Sandoval, and Luis Rodolfo Garcia Carrillo}
\date{October 1, 2026}

\begin{document}
\maketitle

\begin{abstract}
Section~VI-B of the manuscript attaches the same gated CMG--PES residual to six nominal attitude controllers: PD with
desired-rate feedforward in backstepping form (the primary baseline), PID with rate feedforward, LQR, sliding-mode control
(SMC), finite-horizon model predictive control (MPC), and grid-search nonlinear MPC (NMPC). This document gives, for each
family, its architecture, its control law, the derivation of its gains, the exact lines of code that implement it, and the
settings used in the portability study (Amendment A5). All six families share one torque interface and are analyzed in one
coordinate, the rotation-vector error $\bm e=2\bm e_q$, so their small-error behavior can be compared directly
(Table~\ref{tab:gains}): every family reduces near the reference to a second-order loop at the same bandwidth scale, with
damping ratios from $0.71$ (LQR, MPC) to $1$ (SMC). Each gain is read from the class the runs used, and each closed-form
derivation is checked numerically.\ifwithbs{} The family the code calls ``backstepping'' is computed torque; Section~\ref{sec:bstrue}
therefore adds a recursive Lyapunov backstepping design (plain and adaptive) in the style of the EHEKATL BWB controllers,
with its stability derivation, numerical verification, and a 40-run campaign with and without the residual. This new
family is not part of the manuscript.\fi{} The source files cited here are byte-identical to the ones recorded by every A5 run
(\fhashMatch/\fhashChecked{} hashes).
\end{abstract}

\tableofcontents
\newpage

% =============================================================================
\section{How to read this document}
% =============================================================================
\begin{itemize}[leftmargin=1.2em]
\item \textbf{Paper references.} ``\Peq{15}'' is equation (15) of the manuscript; Sec.~VI-B and Proposition~1 refer to the
  manuscript. The equation-to-code companion of September 30 (\code{paper/companion_2026-09-30}) covers the quaternion
  conventions, the residual and the gate in detail; this document does not repeat them.
\item \textbf{Code.} Every listing is read from the repository file at the quoted lines, with the file's own line numbers.
  Four files carry all six families: \BL{} (class \code{AdvancedBaseline}: backstepping, LQR, SMC, MPC, NMPC), \RG{}
  (the A5 plant loop and the PID path), and the campaign drivers \code{run_campaign_a5.py} and
  \code{run_campaign_a5_nmpc_corrected.py}. Appendix~\ref{app:hashes} gives their hashes.
\item \textbf{Status tags.} \emph{Standard}: textbook result, cited. \emph{Derived here}: the steps are written out.
  \emph{Verified}: recomputed by \code{tools/build_family_tables.py} from the code itself. \emph{Implementation note}: a
  detail of the code that the paper leaves implicit, with its effect on reported numbers.
\item \textbf{Figures.} Every architecture uses the manuscript's figure style. The left part (reference and attitude error)
  and the right part (torque box, gated residual, allocation, plant) are identical in all six figures; the dashed blue group is
  the only part that changes from family to family.
\end{itemize}

% =============================================================================
\section{What all six families share}\label{sec:common}
% =============================================================================
\subsection{The residual attaches to any nominal torque}
The residual reads only the nominal request $\bm u_0$ and adds a bounded, gated correction to it (Fig.~\ref{fig:interface}).
Proposition~1 assumes only $\bm u_{0,k}\in\mathcal U=\{\bm u:|u_i|\le\ell_i\}$, so it holds for every family whose output
is saturated to the torque box. All six are: LQR, backstepping, SMC and MPC through the shared clip (line~241 of \BL),
NMPC by clipping every candidate (lines 265 and 273), and PID at line~353 of \RG. Exact recovery at $\beta=0$ was certified
by running the augmentor with $\beta=0$ against the classical run: relative difference $0$ in all five cases (LQR, SMC, MPC
and two NMPC missions; \code{certify_beta0_recovery.py}). Backstepping is covered by the 20/20 bit-exact match with the
sealed primary runs; PID has no separate $\beta=0$ run, and for it the recovery rests on Proposition~1 alone.

\begin{figure}[H]\centering
\resizebox{\textwidth}{!}{% How every family attaches to the same gated CMG--PES residual (Sec. 2). Row 1: the six u0 sources and the
% plant path; rows 2-3: the residual path of the manuscript's Fig. 1, unchanged.
\begin{tikzpicture}
  % ---- six u0 sources -----------------------------------------------------------------------------------------
  \foreach \n/\t [count=\i] in {f1/PD+ff (backstepping),f2/PID+ff,f3/LQR,f4/SMC,f5/MPC,f6/NMPC}
    \node[law,minimum height=6.5mm,minimum width=30mm] (\n) at (0,{1.75-0.9*\i}) {\t};
  \ifwithbs
  \node[law,dashed,minimum height=6.5mm,minimum width=30mm,align=center] (f7) at (0,{1.75-0.9*7-0.25}) {Recursive backstepping\\(new, not in paper)};
  \node[fam,fit=(f1)(f7)] (fams) {};
  \else
  \node[fam,fit=(f1)(f6)] (fams) {};
  \fi
  \node[famlbl,anchor=south] at (fams.north) {one family per run (A5)};
  \ifwithbs\def\nfam{7}\else\def\nfam{6}\fi
  \foreach \i in {1,...,\nfam} \draw[thick] (f\i.east) -- (f\i.east-|fams.east);

  % ---- residual path (manuscript Fig. 1) -----------------------------------------------------------------------
  \node[neuro] (enc)  at (4.9,-3.0) {\textbf{Paired currents}\\$\bar{\bm u}=\bm u_0\oslash\bm\ell$\\
        $\bm I^\pm=I_b\bm1+I_g\max(\pm\bar{\bm u},\bm0)$};
  \node[neuro] (cmg)  [right=7mm of enc] {\textbf{Six CMG units}\\$(v_j,r_j)$, no spikes\\$\bm a=\clip_{[0,1]}(G\bm r^+)$};
  \node[neuro] (dec)  [right=7mm of cmg] {\textbf{Decoder}\\$\bm s=\clip_{[-1,1]}(W\bm a)$};
  \node[clas]  (gate) [right=8mm of dec] {\textbf{Saturation-aware gate}\\$\tilde s_i=0$ if $|u_{0,i}|\ge\ell_i$ and $s_iu_{0,i}<0$\\
        else $\tilde s_i=s_i$};
  \node[neuro] (pes)  [below=11mm of cmg.south west,anchor=north west] {\textbf{PES} (online, every step)\\
        $W^+=\Pi_{\mathcal W}\big[(1-\lambda_W)W+\eta\bm\varepsilon\bm a\T\big]$,\ \ $\bm\varepsilon=\bm y^*-\bm s$};
  \node[neuro] (tgt)  [left=7mm of pes] {\textbf{PES teaching target}\\$\bm y^*=\clip_{[-1,1]}(-k_e\bm e_q)$\\from the measured attitude};

  % ---- plant path ----------------------------------------------------------------------------------------------
  \node[clas]  (sel)   at (enc|-0,0) {\textbf{Nominal torque}\\$\bm u_0=\pi_{\rm fam}(q,\bm\omega,t)\in\mathcal U$};
  \node[clas]  (box)   at (gate|-0,0) {\textbf{Bounded torque addition}\\$\bm u=\sat_{\bm\ell}(\bm u_0+\beta\bm\ell\odot\tilde{\bm s})$};
  \node[plant] (plant) [right=8mm of box] {\textbf{Allocation and plant}\\$\bm f=\clip(B^{-1}[T,\bm u\T]\T)$\\MuJoCo X2 / CF2};
  \draw[sig] (fams.east|-sel) -- (sel);
  \draw[sig] (sel) -- node[lbl,above]{$\bm u_0$} (box);
  \draw[sig] (box) -- node[lbl,above]{$\bm u$} (plant);

  % ---- signals ---------------------------------------------------------------------------------------------------
  \draw[sig] (sel) -- node[lbl,right]{$\bm u_0$} (enc);
  \draw[sig] (enc) -- node[lbl,above]{$\bm I$} (cmg);
  \draw[sig] (cmg) -- node[lbl,above]{$\bm a$} (dec);
  \draw[sig] (dec) -- node[lbl,above]{$\bm s$} (gate);
  \draw[sig] (gate) -- node[lbl,right]{$\tilde{\bm s}$} (box);
  \draw[sig] (cmg.south) -- node[lbl,left]{$\bm a$} (cmg.south|-pes.north);
  \draw[sig] ($(pes.north-|dec)+(-3mm,0)$) -- node[lbl,left]{$W$} ($(dec.south)+(-3mm,0)$);
  \draw[sig] ($(dec.south)+(3mm,0)$) -- node[lbl,right]{$\bm s$} ($(pes.north-|dec)+(3mm,0)$);
  \draw[sig] (tgt) -- node[lbl,above]{$\bm y^*$} (pes);
  \coordinate (u0tap) at ($(sel.east)+(6mm,0)$);
  \fill (u0tap) circle (1.1pt);
  \draw[sigd] (u0tap) -- ++(0,-1.25) -| node[lbl,pos=.25,above]{$\bm u_0$ (load test of the gate)} ($(gate.north)+(-12mm,0)$);

  \begin{scope}[on background layer]
    \node[fam,draw=black!55,fit=(enc)(cmg)(dec)(gate)(pes)(tgt),inner sep=9pt] (rp) {};
  \end{scope}
  \node[note,anchor=north east] at (rp.south east) {CMG--PES residual path of the manuscript, unchanged across families};
\end{tikzpicture}
}
\caption{One residual, six nominal controllers. The nominal torque comes from one family per run; the residual path is the
manuscript's Fig.~1 unchanged ($\beta=.5$ on X2, $.2$ on CF2; $\eta=8\times10^{-4}$; $k_e=8$). The gate and the outer saturation
see only $\bm u_0$, so nothing in the residual depends on how $\bm u_0$ was computed.}
\label{fig:interface}
\end{figure}

\subsection{Error coordinates and the design model}\label{sec:design}
All families act on the shortest-path error $\bm e_q=\sgnp(q_{e,0})\bm q_{e,v}$ (\Peq{4}--\Peq{6}) and the rate error
$\bm\omega_e=\bm\omega-\bm\omega_d$, with $\bm\omega_d$ the central-difference reference rate of the companion, Sec.~2.5.
\code{AdvancedBaseline} works in the rotation-vector coordinate
\begin{equation}
\bm e=2\bm e_q,\qquad \|\bm e\|=2\sin\tfrac{\theta_q}{2}=\theta_q+O(\theta_q^3),
\label{eq:e}
\end{equation}
and designs every law for the per-axis double integrator obtained near the reference. From the error kinematics
\Peq{7}--\Peq{8}, $\dot{\bm e}=q_{e,0}\bm\omega_e+\bm e_q\times\bm\omega_e\approx\bm\omega_e$; from Euler's equation \Peq{2} with
the principal moments $\bm J=\bm j_{\rm tune}$ (Table~I of the paper),
\begin{equation}
\dot e_i\approx\omega_{e,i},\qquad \dot\omega_{e,i}=a_i-\dot\omega_{d,i}-\big[(\bm\omega\times\bm J\bm\omega)\oslash\bm J\big]_i,\qquad
a_i=\tau_i/J_i .
\label{eq:design}
\end{equation}
Every family outputs a desired angular acceleration $\bm a$; the torque is $\bm\tau=\bm J\odot\bm a$ plus, optionally, the
gyroscopic term, saturated to the box:
\begin{equation}
\bm u_0=\sat_{\bm\ell}\big(\bm J\odot\bm a+\chi_{\rm ff}\,\bm\omega\times(\bm J\odot\bm\omega)\big),\qquad
\bm a=\chi_{\rm ff}\dot{\bm\omega}_d+\bm a_{\rm fam}(\bm e,\bm\omega_e).
\label{eq:interface}
\end{equation}
\emph{Derived here.} In every run of the paper $\chi_{\rm ff}=0$ (\code{--no-model-ff}): the acceleration feedforward and the
gyroscopic term are off, so each family is a feedback law on $(\bm e,\bm\omega_e)$ scaled by inertia, and the residual is the
only adaptive element. The rate feedforward $\bm\omega_d$ stays on (\code{--rate-ff left-neg}) for all families.

\listing{experiments/nengo_nef_replication_2026-08-29/run_mujoco_nef_baselines.py}{189}{199}{$\bm\omega_d$ and $\dot{\bm\omega}_d$ (the latter is zero when model feedforward is off)}
\listing{experiments/nengo_nef_replication_2026-08-29/run_mujoco_nef_baselines.py}{216}{245}{\code{AdvancedBaseline.torque}: one interface, five families, shared clip at line 241}

\subsection{Torque box and platform constants}
$\ell_i=J_i\nu_i$ with $\bm\nu=(12,12,10)$\,rad/s$^2$, the same for both platforms, so the box limits every family to the
same angular acceleration $|a_i|\le\nu_i$:

\begin{table}[H]\centering\small
\caption{Constants read from the A5 run manifests (\ccode{derived} field).}\label{tab:platform}
\begin{tabular}{@{}lll@{}}\toprule
& X2 & CF2 (1.5$\times$ rotor design)\\\midrule
$\bm J$ (kg\,m$^2$) & \fJXtwo & \fJCFtwo\\
$\bm\ell$ (N\,m) & \fellXtwo & \fellCFtwo\\
$\bm\ell\oslash\bm J$ (rad/s$^2$) & \fnuXtwo & \fnuCFtwo\\
$\rho\bm\omega_n$ (rad/s) & 12, 12, 10.5 & 12, 12, 10.5\\\bottomrule
\end{tabular}\end{table}

Because the box is defined as an acceleration, every gain in this document is the same for X2 and CF2 (the script asserts
it), and the torque differs only by the inertia factor.

\subsection{The A5 settings}
All six families ran with the same plant, missions, residual, gate, $\rho=3$ bandwidth scale and seed; only the family flags
differ:
\listing{experiments/final_robustness_gating_2026-09-20/run_campaign_a5.py}{27}{40}{per-family flags}
\listing{experiments/final_robustness_gating_2026-09-20/run_campaign_a5.py}{49}{56}{common flags: $\rho=3$, rate feedforward on, model feedforward off, gate on in the CMG arm}

\implnote{The NMPC row of the paper comes from \code{run_campaign_a5_nmpc_corrected.py}, which drops the
\code{--baseline-rate-hz 100} hold of the first A5 run so that NMPC, like every other family, updates at every $0.5$-ms step
(Sec.~\ref{sec:nmpc}). The held run is kept on disk and not used.}

\subsection{Small-error equivalence of the six families}\label{sec:equiv}
Near the reference, inside the SMC boundary layer and away from the box, every family reduces to
$a_i=-k_{1,i}e_i-k_{2,i}\omega_{e,i}\,(-k_{I,i}\!\int\! e_i)$, i.e.\ the characteristic polynomial $s^2+k_2s+k_1$ (or
$s^3+k_2s^2+k_1s+k_I$ for PID). Table~\ref{tab:gains} lists the gains each family actually uses; Table~\ref{tab:poles} gives
the resulting poles of the linearized design model \eqref{eq:design}.

\begin{table}[H]\centering\small
\caption{Equivalent small-error gains, read from \ccode{AdvancedBaseline} with the A5 settings (roll\,/\,yaw; pitch equals roll).
$\omega_{\rm eq}=\sqrt{k_1}$, $\zeta_{\rm eq}=k_2/(2\sqrt{k_1})$. ``Box at'' is the rotation angle at which the
proportional term alone reaches the torque box, $\nu_i/k_{1,i}$. \emph{Verified.}}\label{tab:gains}
\begin{tabular}{@{}lllllll@{}}\toprule
Family & $k_1$ (s$^{-2}$) & $k_2$ (s$^{-1}$) & $\omega_{\rm eq}$ (rad/s) & $\zeta_{\rm eq}$ & box at ($^\circ$) & $k_I$ (s$^{-3}$)\\\midrule
\tabinput{generated/family_gains.tex}
\bottomrule\end{tabular}\end{table}

\begin{table}[H]\centering\small
\caption{Closed-loop poles of the linearized design model \eqref{eq:design} with each family's gains. \emph{Verified.}}\label{tab:poles}
\begin{tabular}{@{}lll@{}}\toprule
Family & axis & poles (s$^{-1}$)\\\midrule
\tabinput{generated/family_poles.tex}
\bottomrule\end{tabular}\end{table}

The families therefore differ in damping and in what happens away from the origin: how each one behaves at the box (SMC
uses a bounded reaching term, PID freezes its integrator, NMPC searches around the nominal law), not in bandwidth. The
proportional term of every linear law reaches the box at a rotation error of $5$--$6^\circ$, so Hardstep's steps of tens of
degrees saturate every family, which is where the gate matters.

% =============================================================================
\section{PD with desired-rate feedforward in computed-torque form (primary baseline)}\label{sec:pdff}
In the code this family is called \code{backstepping} (\code{--baseline backstepping}); it is the computed-torque form.\ifwithbs{}
It is not a recursive backstepping design; the recursive design is Sec.~\ref{sec:bstrue}.\fi
% =============================================================================
\archfig{fig_pdff.tex}{PD+ff in computed-torque (backstepping) form, the paper's $\rho=3$ primary baseline. The model
feedforward path (dotted) exists in the code and is switched off in every run, which makes this law identical to \Peq{15}.}{fig:pdff}

\paragraph{Law.} Computed-torque backstepping \cite{krstic1995nonlinear,slotine1991applied}: choose the error dynamics
$\dot{\bm\omega}_e=-k_1\bm e-k_2\bm\omega_e$ and invert \eqref{eq:design},
\begin{equation}
\bm u_0=\sat_{\bm\ell}\Big(\bm J\odot\big(\dot{\bm\omega}_d-k_1\bm e-k_2\bm\omega_e\big)+\bm\omega\times(\bm J\odot\bm\omega)\Big),\qquad
k_1=(\rho\omega_n)^2,\quad k_2=2\zeta\rho\omega_n,\ \zeta=0.95 .
\label{eq:bs}
\end{equation}
With the feedforward on, no saturation and a diagonal inertia, the gyroscopic term and $\dot{\bm\omega}_d$ cancel and the
rate error obeys $\dot{\bm\omega}_e=-k_1\bm e-k_2\bm\omega_e$ exactly; with $\dot{\bm e}\approx\bm\omega_e$ this is
$\ddot e_i+k_2\dot e_i+k_1e_i\approx0$ near the reference, the textbook reason for this form. \emph{Standard.}

\paragraph{Reduction to the paper's PD+ff.} With $\chi_{\rm ff}=0$ and $\bm e=2\bm e_q$,
$\bm u_0=\sat_{\bm\ell}(-2\bm J(\rho\bm\omega_n)^2\odot\bm e_q-2\zeta\bm J\rho\bm\omega_n\odot\bm\omega_e)$, which is \Peq{13}--\Peq{15}
with $\bm\omega_n\mapsto\rho\bm\omega_n$. \emph{Derived here; verified bit-exact in 20/20 cells}
(\code{run_campaign_a5_backstepping.py}, relative deviation $0$ against the sealed primary runs).

\listing{experiments/nengo_nef_replication_2026-08-29/run_mujoco_nef_baselines.py}{138}{141}{gains $k_1=\omega^2$, $k_2=2\zeta\omega$, $\omega=\rho\omega_n$}
\listing{experiments/final_robustness_gating_2026-09-20/run_controllers_gated.py}{79}{86}{the frozen paper baseline: backstepping, $\rho=3$, $\zeta=.95$, \code{left-neg}, no model feedforward}

\ifwithbs
% =============================================================================
\section{Recursive Lyapunov backstepping (new family)}\label{sec:bstrue}
% =============================================================================
The family of Sec.~\ref{sec:pdff} is computed torque: it chooses linear error dynamics and inverts the model, and with
model feedforward off it is exactly PD+ff. This section adds a backstepping controller in the strict sense: a recursive
design with a virtual control on the kinematic subsystem, a cross-coupling term that cancels the interconnection in the
Lyapunov derivative, and a final Lyapunov function that certifies the closed loop
\cite{krstic1995nonlinear,khalil2002nonlinear}. It follows step by step the design used for the longitudinal and adaptive
controllers of the EHEKATL BWB (strict-feedback form, errors $z_1,z_2$, virtual control with the term $-g_1\T z_1$,
$V=\sum\frac12z_i\T z_i$, La Salle--Yoshizawa, adaptive estimates with $\frac12\tilde f\T\Gamma\tilde f$). The plain and the
adaptive versions are implemented in \code{experiments/backstepping_true_2026-10-01/true_backstepping.py} and run through the
unchanged A5 runner. \textbf{They are not part of the manuscript.}

\archfig{fig_bstrue.tex}{Recursive backstepping. Step~1 stabilizes the quaternion kinematics with a virtual rate; step~2 drives
the rate to it with a torque that cancels the gyroscopic term, feeds forward the transported desired-rate derivative, and
contains the coupling term $-c\,z_1$. The adaptive variant adds a projected estimate of an unknown body torque.}{fig:bstrue}

\subsection{Strict-feedback form}
With the passive convention of the paper, $q_e=q_d\otimes q^{-1}$ and its body-frame rate
$\bm\omega_e=\bm\omega-\bar{\bm\omega}_d$, $\bar{\bm\omega}_d=R(q_e)\T\bm\omega_d$ (\Peq{7}--\Peq{8}, with the transport
$R(q_e)\T$ kept), the attitude system is in strict-feedback form:
\begin{equation}
\dot z_1=\tfrac12T(q_e)\,(\bm\omega-\bar{\bm\omega}_d),\qquad
\bm J\dot{\bm\omega}=-\bm\omega\times\bm J\bm\omega+\bm\tau+\bm d,\qquad
z_1=\bm e_q=\sgnp(q_{e,0})\bm q_{e,v},\quad T(q_e)=|q_{e,0}|I+[z_1]_\times ,
\label{eq:sf}
\end{equation}
i.e.\ $x_1=z_1$, $x_2=\bm\omega$, $f_1=-\frac12T\bar{\bm\omega}_d$, $g_1=\frac12T$, $f_2=-\bm J^{-1}(\bm\omega\times\bm J\bm\omega)$,
$g_2=\bm J^{-1}$ in the thesis notation, with $\bm d$ an unknown body torque (zero in the plain design). There is no actuator
state: the rotor forces of the MuJoCo models act without lag, so the thesis's $z_3$ step is not needed. The first identity
follows from \Peq{8}: $\frac{d}{dt}(\sgnp(q_{e,0})\bm q_{e,v})=\frac12(|q_{e,0}|\bm\omega_e+z_1\times\bm\omega_e)$.

\subsection{Step 1: kinematic subsystem}
Take $V_1=2(1-|q_{e,0}|)\ge0$, zero exactly at $q_e=\pm1$. With $\dot q_{e,0}=-\frac12\bm q_{e,v}\T\bm\omega_e$ (\Peq{8}),
\[
\dot V_1=-2\,\sgnp(q_{e,0})\,\dot q_{e,0}=z_1\T\bm\omega_e .
\]
Treat $\bm\omega_e$ as the control of this subsystem and choose the virtual rate $\bm\omega_e^d=-K_1z_1$, $K_1=\operatorname{diag}(k_{1,i})>0$:
if $\bm\omega_e=\bm\omega_e^d$, then $\dot V_1=-z_1\T K_1z_1$.

\subsection{Step 2: rate subsystem}
Define the second error $z_2=\bm\omega_e-\bm\omega_e^d=\bm\omega_e+K_1z_1$ and $V=c\,V_1+\frac12z_2\T\bm Jz_2$, $c>0$. Using
$\bm\omega_e=z_2-K_1z_1$ and $\dot{\bm\omega}_e=\dot{\bm\omega}-\dot{\bar{\bm\omega}}_d$,
\[
\dot V=c\,z_1\T z_2-c\,z_1\T K_1z_1+z_2\T\big(-\bm\omega\times\bm J\bm\omega+\bm\tau+\bm d-\bm J\dot{\bar{\bm\omega}}_d+\bm JK_1\dot z_1\big).
\]
The control
\begin{equation}
\boxed{\ \bm\tau=\bm\omega\times\bm J\bm\omega+\bm J\big(\dot{\bar{\bm\omega}}_d-K_1\dot z_1\big)-c\,z_1-K_2z_2-\hat{\bm d}\ },\qquad
\bm u_0=\sat_{\bm\ell}(\bm\tau),
\label{eq:bstrue}
\end{equation}
with $K_2=\operatorname{diag}(k_{2,i})>0$ and $\dot z_1=\frac12T(q_e)\bm\omega_e$ computed from measured signals, gives
\begin{equation}
\dot V=-c\,z_1\T K_1z_1-z_2\T K_2z_2+z_2\T\tilde{\bm d},\qquad \tilde{\bm d}=\bm d-\hat{\bm d}.
\label{eq:vdot}
\end{equation}
The term $-c\,z_1$ is the backstepping coupling term: it cancels $c\,z_1\T z_2$, the interconnection between the two
subsystems ($-g_1\T z_1$ in the thesis). Without it the design is computed torque.

\paragraph{Plain law} ($\bm d=\hat{\bm d}=\bm0$). $\dot V=-c\,z_1\T K_1z_1-z_2\T K_2z_2<0$ for $(z_1,z_2)\ne0$; by La Salle--Yoshizawa
$z_1,z_2\to0$, hence $q_e\to\pm1$ and $\bm\omega\to\bar{\bm\omega}_d$, and $z_1,z_2\in\mathcal L_2$. Because
$c\,V_1\ge c\|z_1\|^2/2$ near the origin, the convergence is exponential there.

\paragraph{Adaptive law.} For a constant unknown $\bm d$, add $\frac12\tilde{\bm d}\T\Gamma^{-1}\tilde{\bm d}$ to $V$. Since
$\dot{\tilde{\bm d}}=-\dot{\hat{\bm d}}$, the update
\begin{equation}
\dot{\hat{\bm d}}=\operatorname{Proj}_{|\hat d_i|\le D_i}\big(\Gamma z_2\big)
\label{eq:adapt}
\end{equation}
cancels $z_2\T\tilde{\bm d}$ in \eqref{eq:vdot} and leaves $\dot V_a\le-c\,z_1\T K_1z_1-z_2\T K_2z_2$; the projection keeps
$\hat{\bm d}$ bounded without increasing $\dot V_a$ when $|d_i|\le D_i$ \cite{ioannou1996robust,krstic1995nonlinear}. This is the
thesis's $\dot{\tilde f}=-\Gamma^{-1}z$ law with a bound added. Here $\Gamma_i=0.3\cdot2J_i\omega_i^2/k_{1,i}$, which makes the
integral action on $z_1$ ($\Gamma k_1\!\int\!z_1$) equal to the PID family's, and $D_i=0.25\,\ell_i$.

\paragraph{Rate transport.} $\bar{\bm\omega}_d=R(q_e)\T\bm\omega_d$ and, since $q_e$ rotates with body rate $\bm\omega_e$
($\dot R_e=R_e[\bm\omega_e]_\times$),
\[
\dot{\bar{\bm\omega}}_d=R(q_e)\T\dot{\bm\omega}_d-\bm\omega_e\times\bar{\bm\omega}_d ,
\]
with $\bm\omega_d$, $\dot{\bm\omega}_d$ the same central differences as Sec.~\ref{sec:design}. Unlike the A5 families, this law
uses the transport, the gyroscopic term and $\dot{\bm\omega}_d$ by construction; the \code{--no-model-ff} flag does not apply.

\subsection{Gain design: same linearization as PD+ff}
About $z=0$, \eqref{eq:sf}--\eqref{eq:bstrue} give, per axis, $\ddot z_{1,i}+(k_{1,i}/2+k_{2,i}/J_i)\dot z_{1,i}+(c+k_{1,i}k_{2,i})/(2J_i)\,z_{1,i}=0$.
Writing $k_{1,i}=2\lambda_{1,i}$, $k_{2,i}=J_i\lambda_{2,i}$ and imposing
\[
\lambda_{1,i}+\lambda_{2,i}=2\zeta\omega_i,\qquad \lambda_{1,i}\lambda_{2,i}+\frac{c}{2J_i}=\omega_i^2,\qquad \omega_i=\rho\omega_{n,i},\ \zeta=0.95,
\]
makes the small-error poles exactly those of the paper's PD+ff ($-11.40\pm3.75j$ on roll). Any difference between the two
families is then due to the backstepping structure, not to bandwidth or damping. $V_1$ is a scalar, so one $c$ serves all
axes; real $\lambda$'s need $c\ge2J_i\omega_i^2(1-\zeta^2)$ on every axis, and the smallest such value is used
($\lambda_1\le\lambda_2$: the outer kinematic loop is the slower one). For X2: $c=\fbsC$\,N\,m,
$\bm k_1=(\fbsKone)$\,s$^{-1}$, $\bm k_2=(\fbsKtwo)$\,N\,m\,s.

\subsection{Verification}
\code{verify_true_backstepping.py} checks the derivation on the rigid-body model (\code{results/verification.json}). \emph{Verified}:
\begin{itemize}[leftmargin=1.2em]
\item $\dot z_1$ and $\dot{\bar{\bm\omega}}_d$ in closed form against central finite differences along rigid-body trajectories on Full3,
  Hard45 and Yaw360: worst absolute error \fbsKinErr{} (signals of order $0.1$--$1$);
\item closed loop, plain law, no torque limit, Full3 and Hard45 from a $\approx20^\circ$ initial error: $V$ decreases monotonically to
  $10^{-12}$, and the finite-difference $\dot V$ equals $-c\,z_1\T K_1z_1-z_2\T K_2z_2$ within \fbsVdotErr{} (relative, step-size limited);
\item adaptive law with a constant unknown torque of $20\%$ of the box: $\hat{\bm d}$ reaches $\bm d$ within \fbsAdaptErr{} after $20$\,s;
\item linearized poles equal the PD+ff poles to \fbsPoleErr.
\end{itemize}

\listing{experiments/backstepping_true_2026-10-01/true_backstepping.py}{44}{55}{gain design: smallest common $c$, $\lambda_1\le\lambda_2$}
\listing{experiments/backstepping_true_2026-10-01/true_backstepping.py}{84}{114}{$z_1$, $z_2$, $\dot z_1$, transport, law \eqref{eq:bstrue}, projected adaptation \eqref{eq:adapt}}
\listing{experiments/backstepping_true_2026-10-01/run_true_backstepping.py}{30}{42}{the family enters the unchanged A5 runner}

\implnote{The Lyapunov certificate holds for the model \eqref{eq:sf} without saturation. In the plant the torque box, rotor
clipping, the non-diagonal X2 inertia (the law uses the principal moments, as every family does) and the residual add terms
that the certificate does not cover; the sign selection makes $V_1$ non-smooth at $q_{e,0}=0$, the same local scope as the
paper. The output is saturated to the box, so Proposition~1 applies unchanged.}

\subsection{Results: recursive backstepping with and without the residual}
Same protocol as A5 (Table~\ref{tab:results}): 45-s missions, $\rho=3$, noise-free, seed 1, gate on, the residual unchanged. Run by \code{experiments/backstepping_true_2026-10-01/run_campaign_true_bs.py}; numbers computed from the run summaries by \code{tools/build_family_tables.py}.
\begin{table}[H]\centering\small
\caption{Reduction of mean quaternion-angle error by the gated CMG--PES residual on the recursive backstepping laws (\%, positive is better; 5-mission = reduction of geometric means; last column: baseline geometric mean, mrad).}\label{tab:bsres}
\begin{tabular}{@{}llrrrrrrr@{}}\toprule
Family & Platform & Roll & Full3 & Hard45 & Hardstep & Yaw360 & 5-mission & base (mrad)\\\midrule
Recursive backstepping & X2 & -114.8 & +7.2 & +11.8 & -2.9 & +43.5 & \textbf{-0.4} & 3.30\\
 & CF2 & -1416.5 & -18.2 & -38.1 & -0.8 & -1467.9 & \textbf{-230.0} & 0.59\\\addlinespace
Adaptive recursive backstepping & X2 & -114.7 & +6.5 & +11.7 & -0.8 & +30.9 & \textbf{-4.3} & 2.63\\
 & CF2 & -669.0 & -9.6 & -19.6 & -0.4 & -2494.5 & \textbf{-204.7} & 0.88\\\addlinespace
\bottomrule\end{tabular}\end{table}
\begin{table}[H]\centering\small
\caption{Nominal laws alone (no residual): error reduction of the recursive designs relative to the computed-torque PD+ff (the paper's primary baseline), \% (positive: the recursive law tracks better). The linearizations are identical, so these differences come from the backstepping structure.}\label{tab:bsvsct}
\begin{tabular}{@{}llrrrrrr@{}}\toprule
Platform & law & Roll & Full3 & Hard45 & Hardstep & Yaw360 & 5-mission\\\midrule
X2 & recursive & +29.6 & +30.5 & +31.0 & -2.4 & +41.5 & \textbf{+27.4}\\
X2 & adaptive & +36.5 & +30.8 & +29.1 & -7.9 & +80.7 & \textbf{+42.2}\\
CF2 & recursive & +95.0 & +64.8 & +73.3 & -0.5 & +97.9 & \textbf{+84.2}\\
CF2 & adaptive & +91.2 & +47.5 & +58.9 & -695.8 & +99.6 & \textbf{+76.7}\\
\bottomrule\end{tabular}\end{table}


\paragraph{What the two tables show.}
\begin{itemize}[leftmargin=1.2em]
\item \textbf{The recursive law is a much better nominal controller} (Table~\ref{tab:bsvsct}): with the same linearization it
  reduces the computed-torque PD+ff error on four of five missions on both platforms, most strongly on CF2. Its law contains
  $\bm J\dot{\bar{\bm\omega}}_d$ and $\bm\omega\times\bm J\bm\omega$, which every A5 family ran without (\code{--no-model-ff}). The
  computed-torque law \emph{with} model feedforward was not run, so this campaign does not separate the share of the gain due to the
  acceleration feedforward from the share due to the recursive structure (cross term and rate transport). That run is the
  natural next control.
\item \textbf{The unchanged residual does not help it} (Table~\ref{tab:bsres}). It still reduces the error on X2 Full3, Hard45
  and Yaw360, but where the recursive law's error is already very small (Roll; CF2 Roll and Yaw360, a few hundredths of a mrad) it
  raises the error to tenths of a mrad, worse in absolute terms than the residual on the paper's baseline. The residual's
  $\beta$, $\eta$ and $k_e$ were set for a lagging PD+ff; on a nominal law that already cancels the dominant error, the residual's
  own correction becomes the error. Retuning ($\beta$ in particular) was not tested. This is a limit of the unchanged residual,
  not of Proposition~1, which still holds: the output stays in the box and $\beta=0$ recovers the law.
\item \textbf{The adaptive variant fails on CF2 Hardstep} under persistent rotor saturation (79\% of samples): the projected estimate
  keeps integrating $z_2$ while the torque it requests cannot be delivered, a windup through the adaptation law that the
  Lyapunov argument (which assumes no saturation) does not exclude. Freezing the adaptation when the box or a rotor saturates, the
  rule the PID family uses, is the standard fix and was not run. On the other missions the adaptive and plain laws are close.
\end{itemize}
\fi % withbs

% =============================================================================
\section{PID with rate feedforward and rotor-level anti-windup}\label{sec:pid}
% =============================================================================
\archfig{fig_pid.tex}{PID+ff. The integrator runs on $\bm e_q$ and is frozen per axis whenever the torque box clipped that axis,
and on all axes whenever any rotor saturated at the previous step (dashed paths).}{fig:pid}

\paragraph{Law.} The PID path bypasses \code{AdvancedBaseline} (\code{--baseline pd}) and adds an integral term to \Peq{15}:
\begin{equation}
\bm u_0=\sat_{\bm\ell}\big(-\bm k_q\odot\bm e_q-\bm k_\omega\odot\bm\omega_e-\bm k_i\odot\bm z\big),\qquad
\bm k_q=2\bm J\odot(\rho\bm\omega_n)^2,\ \bm k_\omega=2\zeta\bm J\odot\rho\bm\omega_n,\ \bm k_i=0.3\,\bm k_q .
\label{eq:pid}
\end{equation}
\paragraph{Conditional integration.} With $h=5\times10^{-4}$\,s and the clip test of line 357,
\begin{equation}
z_{i,k+1}=z_{i,k}+h\,e_{q,i,k}\,m_{i,k},\qquad
m_{i,k}=\mathbf 1\big[u_{0,i,k}=\tau_{{\rm u},i,k}\big]\cdot\mathbf 1\big[\text{no rotor saturated at }k-1\big].
\label{eq:aw}
\end{equation}
This is conditional integration (integrator clamping) \cite{astrom2006advancedpid}, applied at the rotor level as well as at the
torque box: a rotor at its force limit means the realized torque differs from the command even if the box is inactive.
\emph{Standard; the rotor-level rule is this study's choice.}

\paragraph{Loop.} In the coordinate $\bm e$ the gains per unit inertia are $k_1=(\rho\omega_n)^2$, $k_2=2\zeta\rho\omega_n$ and
$k_I=0.3(\rho\omega_n)^2$ (Table~\ref{tab:gains}). The Routh condition for $s^3+k_2s^2+k_1s+k_I$ is $k_2k_1>k_I$, met with
a ratio of at least \fpidRouthMargin. The integral adds a slow real pole at $-0.32$\,s$^{-1}$ (time constant $\approx3$\,s,
Table~\ref{tab:poles}): the integrator is weak, which the A5 report discloses as ``undertuned''. \emph{Derived here; verified.}

\listing{experiments/final_robustness_gating_2026-09-20/run_controllers_gated.py}{253}{260}{gains; PID uses the plain PD path}
\listing{experiments/final_robustness_gating_2026-09-20/run_controllers_gated.py}{347}{361}{\eqref{eq:pid} and the conditional integrator \eqref{eq:aw}}
\listing{experiments/final_robustness_gating_2026-09-20/run_controllers_gated.py}{475}{475}{rotor-saturation flag for the next step}

\implnote{The integrator acts on $\bm e_q$ (not $\bm e=2\bm e_q$), so $\bm k_i=0.3\bm k_q$ means
$k_I=0.3(\rho\omega_n)^2$ in Table~\ref{tab:gains}'s coordinate. With \code{--ki-factor 0} this path is \Peq{15} exactly.}

% =============================================================================
\section{LQR}\label{sec:lqr}
% =============================================================================
\archfig{fig_lqr.tex}{LQR. The gain is computed once from the closed-form solution of the continuous algebraic Riccati
equation (CARE) for the per-axis double integrator; online it is a fixed state feedback.}{fig:lqr}

\paragraph{Design problem.} For each axis, $\bm x=[e,\omega_e]\T$, $\dot{\bm x}=A\bm x+Ba$ with
$A=\left[\begin{smallmatrix}0&1\\0&0\end{smallmatrix}\right]$, $B=[0,1]\T$, minimize
$\int_0^\infty(q_e e^2+ra^2)\,dt$ \cite{andersonmoore1990optimal}.

\paragraph{Closed-form CARE.} With $P=\left[\begin{smallmatrix}p_1&p_2\\p_2&p_3\end{smallmatrix}\right]$, the CARE
$A\T P+PA-PBr^{-1}B\T P+Q=0$ reads, entry by entry,
\[
(1,1):\ q_e-\frac{p_2^2}{r}=0,\qquad (1,2):\ p_1-\frac{p_2p_3}{r}=0,\qquad (2,2):\ 2p_2-\frac{p_3^2}{r}=0 .
\]
The positive-definite solution has $p_2=\sqrt{q_er}$, $p_3=\sqrt{2p_2r}$, $p_1=p_2p_3/r$, so
\begin{equation}
K=r^{-1}B\T P=\Big[\sqrt{q_e/r},\ \sqrt{2\sqrt{q_e/r}}\Big]\quad\xrightarrow{\ q_e/r=\omega^4\ }\quad K=[\omega^2,\ \sqrt2\,\omega],
\label{eq:lqr}
\end{equation}
a closed loop $s^2+\sqrt2\omega s+\omega^2$ with $\zeta=1/\sqrt2$ at the same $\omega=\rho\omega_n$ as the other families.
\emph{Derived here. Verified}: $P$ from these formulas leaves a CARE residual of \flqrCareRelErr{} (relative), its smallest
eigenvalue is \flqrPminEig{} $>0$, and $B\T P$ equals the code's gain.

\listing{experiments/nengo_nef_replication_2026-08-29/run_mujoco_nef_baselines.py}{132}{137}{$K=[\omega^2,\sqrt2\omega]$}
\listing{experiments/nengo_nef_replication_2026-08-29/run_mujoco_nef_baselines.py}{232}{233}{the law: $\bm a=\dot{\bm\omega}_d-k_1\bm e-k_2\bm\omega_e$}

% =============================================================================
\section{Sliding-mode control}\label{sec:smc}
% =============================================================================
\archfig{fig_smc.tex}{SMC. A linear sliding variable and a boundary-layer reaching law; inside the layer the law is a critically
damped PD, outside it the reaching term has a fixed magnitude of half the torque box.}{fig:smc}

\paragraph{Law.} Sliding variable and reaching law \cite{utkin1992sliding,slotine1991applied}, per axis:
\begin{equation}
\sigma_i=\omega_{e,i}+\lambda_ie_i,\qquad
a_i=\dot\omega_{d,i}-\lambda_i\omega_{e,i}-\eta_i\,\sat\!\big(\sigma_i/\phi_i\big),\qquad
\lambda_i=\rho\omega_{n,i},\ \ \eta_i=0.5\,\ell_i/J_i,\ \ \phi_i=\eta_i/\lambda_i .
\label{eq:smc}
\end{equation}
With the A5 setting, $\bm\eta=(\fsmcEtaXtwo)$\,rad/s$^2$ and $\bm\phi=(\fsmcPhiXtwo)$\,rad/s on both platforms.

\paragraph{Inside the boundary layer} ($|\sigma_i|<\phi_i$): $a_i=-\lambda_i\omega_{e,i}-(\eta_i/\phi_i)(\omega_{e,i}+\lambda_ie_i)$, and
$\eta_i/\phi_i=\lambda_i$, so
\[
a_i=-\lambda_i^2e_i-2\lambda_i\omega_{e,i}:\qquad k_1=\lambda^2,\quad k_2=2\lambda,\quad \zeta=1,
\]
a double pole at $-\rho\omega_n$ (Table~\ref{tab:poles}). \emph{Derived here; verified from the code's
\code{lam}, \code{eta}, \code{phi}.}

\paragraph{Outside the layer.} On the design model \eqref{eq:design} with model feedforward off,
$\dot\sigma_i=\dot\omega_{e,i}+\lambda_i\dot e_i=a_i-\dot\omega_{d,i}+\lambda_i\omega_{e,i}=-\eta_i\operatorname{sgn}(\sigma_i)-\dot\omega_{d,i}$,
hence
\[
\sigma_i\dot\sigma_i\le-(\eta_i-|\dot\omega_{d,i}|)\,|\sigma_i| ,
\]
the reaching condition, whenever the reference acceleration stays below $\eta_i$ ($6$\,rad/s$^2$ on roll and pitch). The layer
is then reached in finite time and the error decays inside it; the boundary layer trades exact sliding for a continuous,
chattering-free command. Because $J_i\eta_i=\ell_i/2$, the reaching term alone uses half the torque box, leaving the
other half for $-\lambda\bm\omega_e$. \emph{Standard argument, specialized here.}

\listing{experiments/nengo_nef_replication_2026-08-29/run_mujoco_nef_baselines.py}{142}{149}{$\lambda=\omega$, $\eta=0.5\,\bm\ell\oslash\bm J$, $\phi=\eta/\omega$}
\listing{experiments/nengo_nef_replication_2026-08-29/run_mujoco_nef_baselines.py}{234}{236}{\eqref{eq:smc}}

% =============================================================================
\section{Model predictive control (unconstrained, finite horizon)}\label{sec:mpc}
% =============================================================================
\archfig{fig_mpc.tex}{MPC. The finite-horizon quadratic program has no constraints, so its solution is linear in the state; the
first move of the optimal sequence is a fixed gain $K_0$, computed once and applied at every plant step.}{fig:mpc}

\paragraph{Prediction model.} Per axis, the exact zero-order-hold discretization of the double integrator with sample time
$T=20$\,ms,
\[
\bm x_{k+1}=A_d\bm x_k+B_da_k,\qquad A_d=\begin{bmatrix}1&T\\0&1\end{bmatrix},\quad B_d=\begin{bmatrix}T^2/2\\T\end{bmatrix}.
\]
\paragraph{Condensed problem.} Stacking $N=120$ predictions, $X=S_x\bm x_0+S_uU$ with
$S_x=[A_d;A_d^2;\dots;A_d^N]$ and the block-Toeplitz $S_u$ with blocks $A_d^{k-j}B_d$ ($j\le k$). The cost
$\sum_{k=1}^{N}\bm x_k\T Q\bm x_k+\sum_{k=0}^{N-1}Ra_k^2$ with $Q=\operatorname{diag}(\omega^4,q_\omega\omega^2)$, $q_\omega=0$, $R=1$ becomes
\[
U\T HU+2\bm x_0\T F\T U+\text{const},\qquad H=S_u\T\bar QS_u+\bar R,\quad F=S_u\T\bar QS_x,
\]
with $\bar Q=I_N\otimes Q$, $\bar R=I_N\otimes R$. Without constraints the minimizer is $U^\star=-H^{-1}F\bm x_0$, and the
receding-horizon law applies its first element, $a_0=-K_0\bm x_0$ with $K_0=[H^{-1}F]_{1,:}$
\cite{borrelli2017predictive,rawlings2017mpc}. \emph{Standard.}

\paragraph{What the horizon does.} Table~\ref{tab:mpc} recomputes $K_0$ with the code's own function for growing $N$. From
$N\approx30$ ($0.6$\,s, several closed-loop time constants) the first-move gain no longer changes: at the A5 horizon it equals
the infinite-horizon discrete LQR gain obtained by iterating the Riccati recursion to its fixed point to a relative
\fmpcVsDare{} ($K_0=[\fmpcKone,\ \fmpcKtwo]$ on roll; DARE residual \fdareResidual). The MPC row is therefore a discrete LQR
designed at $T=20$\,ms. That sample time, not the horizon, is why its proportional gain is $16\%$ below the continuous LQR
of Sec.~\ref{sec:lqr} ($121.5$ vs.\ $144$\,s$^{-2}$). \emph{Verified.}

\begin{table}[H]\centering\small
\caption{First-move MPC gain on the roll axis ($\omega=12$\,rad/s) versus horizon, from \ccode{AdvancedBaseline._mpc_first_gain};
the last two columns divide by the continuous LQR gain $[\omega^2,\sqrt2\omega]$. A5 uses $N=120$.}\label{tab:mpc}
\begin{tabular}{@{}rrrrrr@{}}\toprule
$N$ & $NT$ (s) & $K_{0,1}$ & $K_{0,2}$ & $K_{0,1}/\omega^2$ & $K_{0,2}/(\sqrt2\omega)$\\\midrule
\tabinput{generated/mpc_horizon.tex}
\bottomrule\end{tabular}\end{table}

\listing{experiments/nengo_nef_replication_2026-08-29/run_mujoco_nef_baselines.py}{167}{187}{condensed matrices $S_x,S_u,H,F$ and the first-move gain}
\listing{experiments/nengo_nef_replication_2026-08-29/run_mujoco_nef_baselines.py}{237}{238}{the online law $\bm a=\dot{\bm\omega}_d-K_0[\bm e;\bm\omega_e]$}

\implnote{The gain is designed for a 20-ms move but applied as a continuous-time state feedback at every 0.5-ms step (A5 ran
every family at full rate). This is a common simplification; it is the reason the MPC row behaves like a softer LQR
($\zeta_{\rm eq}=0.71$, $\omega_{\rm eq}=11.0$\,rad/s) and not like a sampled controller. ``MPC'' here means an unconstrained
receding-horizon gain; the torque box is enforced by the shared clip, not inside the optimization.}

% =============================================================================
\section{Nonlinear MPC (grid search around the backstepping law)}\label{sec:nmpc}
% =============================================================================
\archfig{fig_nmpc.tex}{NMPC. At every plant step, 27 candidate offsets are added to the nominal backstepping torque; each
candidate is scored by an RK4 rollout of the closed loop on the rigid-body model over $0.3$\,s, and the first torque of the
cheapest candidate is applied.}{fig:nmpc}

\paragraph{Candidates.} The nominal law is \eqref{eq:bs} ($\zeta=0.95$). A grid of $3^3=\fnmpcCandidates$ constant offsets
$\bm d_c\in\{-1,0,1\}^3\odot0.1\bm\ell$ is added to it; $\bm d_c=\bm0$ is one of them (\fnmpcHasZero), so the nominal law is always
a candidate.

\paragraph{Prediction (``policy'' mode).} From the measured state, each candidate is rolled out for $N=5$ steps of
$\Delta=60$\,ms with the rigid-body model $\dot Q=\frac12Q\otimes\tilde{\bm\omega}$,
$\dot{\bm\omega}=(\bm\tau-\bm\omega\times\bm J\bm\omega)\oslash\bm J$, integrated by classical RK4 \cite{hairer1993ode}. At
$k=0$ the candidate's first torque $\bm\tau^{(c)}_0=\sat_{\bm\ell}(\bm\tau_{\rm nom}+\bm d_c)$ is used; at $k\ge1$ the nominal law is
re-evaluated on the predicted state and the same offset added,
$\bm\tau^{(c)}_k=\sat_{\bm\ell}(\bm J\odot\bm a(\hat q_k,\hat{\bm\omega}_k,t+k\Delta)+\bm d_c)$. The prediction thus models the
closed loop that will actually run, not a constant torque.

\paragraph{Cost and selection.} With $\bm e_k,\bm\omega_{e,k}$ the predicted errors against the reference at $t+(k+1)\Delta$,
\begin{equation}
J_c=\sum_{k=0}^{N-1}\Big[w_k\big(\textstyle\sum_i\omega_i^4e_{k,i}^2+2\sum_i\omega_i^2\omega_{e,k,i}^2\big)+\big\|\bm\tau^{(c)}_k\oslash\bm J\big\|^2\Big],
\qquad w_k=\begin{cases}1.5,&k=N-1\\1,&\text{otherwise}\end{cases}
\label{eq:nmpc}
\end{equation}
and $\bm u_0=\bm\tau^{(c^\star)}_0$, $c^\star=\arg\min_cJ_c$. The angle weight $\omega^4$ is the LQR weight $q_e/r$ of
Sec.~\ref{sec:lqr}; the effort term is the squared commanded acceleration. This is a sampling-based (grid) receding-horizon
controller in the sense of \cite{rawlings2017mpc}, not a gradient-based NMPC solver. \emph{Derived here from the code.}

\listing{experiments/nengo_nef_replication_2026-08-29/run_mujoco_nef_baselines.py}{152}{159}{candidate grid and cost weights}
\listing{experiments/nengo_nef_replication_2026-08-29/run_mujoco_nef_baselines.py}{201}{214}{rigid-body model and RK4 step (MuJoCo's body-to-world convention)}
\listing{experiments/nengo_nef_replication_2026-08-29/run_mujoco_nef_baselines.py}{247}{289}{\code{_nmpc}: rollout, cost \eqref{eq:nmpc}, selection, one-period self-check}
\listing{experiments/final_robustness_gating_2026-09-20/run_campaign_a5_nmpc_corrected.py}{25}{26}{corrected A5 settings: no update hold}

\implnote{(i) The rollout predicts the nominal family only; it does not know the residual that will be added. (ii) The first A5
NMPC run updated at 100\,Hz and held the torque between updates, a leftover from an earlier exploratory setting. The audit
(\code{AUDIT_NMPC_HOLD_RATE_2026-09-25.md}) showed that the hold, not the predictive machinery, caused that row's anomaly:
removing it restored NMPC, and imposing the same hold on MPC broke MPC identically. The paper's NMPC row is the full-rate re-run.
(iii) The class logs a one-period-ahead self-check of its own prediction (\code{nmpc_pred_err_mean_deg}); it was used to
diagnose (ii), not as a tuning signal.}

% =============================================================================
\section{Results across families (Sec.~VI-B)}\label{sec:results}
% =============================================================================
\begin{table}[H]\centering\small
\caption{Reduction of mean quaternion-angle error by the gated CMG--PES residual on each nominal family, noise-free, $\rho=3$,
45-s missions (\% , positive is better; ``5-mission'' is the reduction of geometric means). Read from
\ccode{results/a5_numbers.json}, which \ccode{analyze_a5.py} computes from the run summaries; NMPC from the corrected
full-rate run. Last column: baseline geometric mean (mrad).}\label{tab:results}
\begin{tabular}{@{}llrrrrrrr@{}}\toprule
Family & Platform & Roll & Full3 & Hard45 & Hardstep & Yaw360 & 5-mission & base (mrad)\\\midrule
\tabinput{generated/family_results.tex}
\bottomrule\end{tabular}\end{table}

These are the numbers the manuscript reports (X2/CF2: backstepping $37.3/22.4$, PID+ff $36.7/22.4$, LQR $31.1/22.4$, SMC
$43.8/26.6$, MPC $33.4/22.7$, NMPC $31.9/21.7$), and the claim verifier checks them (\code{reproduce/verify_acc2027_claims.py}).
What the table shows, without further interpretation:
\begin{itemize}[leftmargin=1.2em]
\item The five-mission reduction stays within $31$--$44\%$ on X2 and $22$--$27\%$ on CF2 across families whose baseline errors
  differ by up to $1.7\times$ (X2 geometric means $3.7$--$6.3$\,mrad), with the residual unchanged.
\item The reduction is largest on Roll and Yaw360 and smallest on Hardstep, whose step changes of tens of degrees drive every
  linear family far past the $5$--$6^\circ$ at which its proportional term reaches the box (Table~\ref{tab:gains}).
\item SMC is the only family with a clearly positive Hardstep result ($+34.2\%$ X2, $+12.2\%$ CF2). SMC is also the only
  family whose law is itself bounded below the box away from the origin (Sec.~\ref{sec:smc}); A5 did not isolate whether that
  is the cause, so it is recorded here as an open observation.
\end{itemize}

% =============================================================================
\section{Scope of the family comparison}
% =============================================================================
\begin{itemize}[leftmargin=1.2em]
\item The families are \emph{attachment tests}, not tuned competitors: all use the same bandwidth scale $\rho=3$, model
  feedforward off, and the A5 settings above. No family was retuned per mission or per platform, and the residual
  ($\beta$, $\eta$, $k_e$) is the same for all.
\item MPC is unconstrained and realized as its first-move gain; NMPC is a 27-point grid search around backstepping. Neither is a
  constrained optimal controller; the torque box acts through the shared saturation.
\item PID's integrator is weak (slow pole at $-0.32$\,s$^{-1}$) and is disclosed as undertuned in the A5 report.
\item All gains are designed on the per-axis double integrator \eqref{eq:design} with principal moments; the X2 inertia tensor is
  not diagonal, so the poles of Table~\ref{tab:poles} are design values, not measured closed-loop poles.
\end{itemize}

\appendix
\section{Source files and hashes}\label{app:hashes}
SHA-256 (first 16 hex digits) of each cited file, compared with the \code{source_sha256} field of one A5 run per family and
platform (Roll, classical arm; NMPC from the corrected run): \fhashMatch/\fhashChecked{} recorded hashes are identical.
\begin{table}[H]\centering\small
\begin{tabular}{@{}p{0.55\linewidth}ll@{}}\toprule
File & SHA-256 & vs.\ A5 manifests\\\midrule
\tabinput{generated/family_hashes.tex}
\bottomrule\end{tabular}\end{table}

\bibliographystyle{IEEEtran}
\bibliography{../v15_2026-09-29_corrected/references_ACC2027,family_refs}
\end{document}
"
\newif\ifwithbs \withbstrue
\ifdefined\releasebuild \withbsfalse \fi
% the manuscript's own figure style, plus the family extensions
% Shared TikZ style for all architecture figures. Grayscale-safe.
\tikzset{
  blk/.style   ={draw,rounded corners=2pt,align=center,font=\footnotesize,
                 minimum height=8mm,inner sep=3pt,fill=white},
  clas/.style  ={blk,fill=black!4},
  neuro/.style ={blk,fill=black!12,very thick},
  plant/.style ={blk,fill=black!20},
  meas/.style  ={blk,dashed},
  sum/.style   ={draw,circle,inner sep=0pt,minimum size=4.5mm,font=\scriptsize},
  sig/.style   ={-{Latex[length=2mm]},thick},
  sigd/.style  ={-{Latex[length=2mm]},thick,dashed},
  lbl/.style   ={font=\scriptsize,inner sep=1.5pt,fill=white},
  grp/.style   ={draw,dashed,rounded corners=3pt,inner sep=6pt},
  note/.style   ={font=\scriptsize\itshape,align=left},
}

% Extra styles for the controller-family figures; the base styles (blk, clas, neuro, plant, sig, ...) are the
% manuscript's own architectures/tikz/arch_style.tex, so every figure matches the paper's Fig. 1.
\tikzset{
  law/.style    ={blk,fill=blue!7,draw=blue!45!black,thick},
  offl/.style   ={blk,dashed,fill=black!2},
  fam/.style    ={draw=blue!45!black,dashed,rounded corners=3pt,inner sep=7pt},
  famlbl/.style ={font=\scriptsize\itshape,text=blue!45!black,inner sep=1pt},
  off/.style    ={-{Latex[length=2mm]},thick,dotted,gray},
  every node/.append style={font=\footnotesize},
}
% Common left half: reference generator and attitude error (identical for every family).
\newcommand{\commonleft}{%
  \node[clas] (ref) at (0,0) {\textbf{Mission reference}\\$q_d(t)$\\[1pt]
        $\bm\omega_d$: central difference};
  \node[clas,right=7mm of ref] (err) {\textbf{Attitude error}\\$q_e=q_d\otimes q^{-1}$\\
        $\bm e=2\,\sgnp(q_{e,0})\,\bm q_{e,v}$\\$\bm\omega_e=\bm\omega-\bm\omega_d$};
  \draw[sig] (ref) -- node[lbl,above]{$q_d,\bm\omega_d$} (err);
}
% Common right half and feedback of every family figure. Needs nodes (err) and (law-out) placed by the figure,
% and the coordinate (satx) giving the x position of the torque box.
\newcommand{\commonloop}[2]{% #1: label on u0 arrow source node, #2: text inside the torque-box node
  \node[clas] (sat) at (satx) {#2};
  \node[neuro,below=13mm of sat] (res) {\textbf{Gated CMG--PES residual}\\(Fig.~\ref{fig:interface})\\
        $\bm u=\sat_{\bm\ell}(\bm u_0+\beta\bm\ell\odot\tilde{\bm s})$};
  \node[plant,anchor=east] (plant) at ($(res.west)+(-12mm,0)$) {\textbf{Allocation and MuJoCo plant}\\
        $\bm f=\clip_{[\bm0,\bm f_{\max}]}(B^{-1}[T,\bm u\T]\T)$\\X2 / CF2 rigid body};
  \draw[sig] (#1) -- (sat);
  \draw[sig] (sat) -- node[lbl,right]{$\bm u_0\in\mathcal U$} (res);
  \draw[sig] (res) -- node[lbl,above]{$\bm u$} (plant);
  \draw[sig] (plant.west) -- ++(-4mm,0) coordinate (fb1) -| node[lbl,pos=.75,left]{measured $q,\ \bm\omega$} (err.south);
}

\newcommand{\flqrCareRelErr}{2.7e-18}
\newcommand{\flqrPminEig}{7.39}
\newcommand{\fsmcEtaXtwo}{6.00, 6.00, 5.00}
\newcommand{\fsmcPhiXtwo}{0.500, 0.500, 0.476}
\newcommand{\fsmcEtaCFtwo}{6.00, 6.00, 5.00}
\newcommand{\fsmcPhiCFtwo}{0.500, 0.500, 0.476}
\newcommand{\fpidRouthMargin}{66.5}
\newcommand{\fdareResidual}{9.5e-14}
\newcommand{\fmpcDareKone}{121.55}
\newcommand{\fmpcDareKtwo}{15.592}
\newcommand{\fmpcKone}{121.55}
\newcommand{\fmpcKtwo}{15.592}
\newcommand{\fmpcVsDare}{4.6e-10}
\newcommand{\fnmpcCandidates}{27}
\newcommand{\fnmpcHasZero}{yes}
\newcommand{\fhashMatch}{60}
\newcommand{\fhashChecked}{60}
\newcommand{\fbsC}{1.705}
\newcommand{\fbsKone}{22.80, 16.69, 10.40}
\newcommand{\fbsKtwo}{0.692, 0.527, 0.375}
\newcommand{\fbsKinErr}{$5.6\times10^{-7}$}
\newcommand{\fbsVdotErr}{0.22\%}
\newcommand{\fbsAdaptErr}{0.25\%}
\newcommand{\fbsPoleErr}{$2.7\times10^{-15}$}
\newcommand{\fbsGeotrueXtwo}{-0.4}
\newcommand{\fbsGeotrueCFtwo}{-230.0}
\newcommand{\fbsGeoadaptiveXtwo}{-4.3}
\newcommand{\fbsGeoadaptiveCFtwo}{-204.7}
\newcommand{\fbsCampaignComplete}{yes}
\newcommand{\fJXtwo}{0.06071, 0.03647, 0.02541}
\newcommand{\fellXtwo}{0.7285, 0.4376, 0.2541}
\newcommand{\fnuXtwo}{12.00, 12.00, 10.00}
\newcommand{\fJCFtwo}{2.395e-05, 2.395e-05, 3.235e-05}
\newcommand{\fellCFtwo}{0.0002874, 0.0002874, 0.0003235}
\newcommand{\fnuCFtwo}{12.00, 12.00, 10.00}


\newcommand{\BL}{\code{experiments/nengo_nef_replication_2026-08-29/run_mujoco_nef_baselines.py}}
\newcommand{\RG}{\code{experiments/final_robustness_gating_2026-09-20/run_controllers_gated.py}}
\newcommand{\archfig}[3]{% file, caption, label
  \begin{figure}[H]\centering\resizebox{\textwidth}{!}{\input{figures/#1}}\caption{#2}\label{#3}\end{figure}}

\title{Online PES Decoder Learning in Neuromorphic Residual UAV Attitude Control\\[4pt]
\large The nominal controller families of Sec.~VI-B\ifwithbs{} and a recursive backstepping design\fi:\\ equations, derivations, code, and architectures\\
\normalsize companion to the ACC~2027 manuscript (version of September 30, 2026)}
\author{Gandhi Contreras, Steven Sandoval, and Luis Rodolfo Garcia Carrillo}
\date{October 1, 2026}

\begin{document}
\maketitle

\begin{abstract}
Section~VI-B of the manuscript attaches the same gated CMG--PES residual to six nominal attitude controllers: PD with
desired-rate feedforward in backstepping form (the primary baseline), PID with rate feedforward, LQR, sliding-mode control
(SMC), finite-horizon model predictive control (MPC), and grid-search nonlinear MPC (NMPC). This document gives, for each
family, its architecture, its control law, the derivation of its gains, the exact lines of code that implement it, and the
settings used in the portability study (Amendment A5). All six families share one torque interface and are analyzed in one
coordinate, the rotation-vector error $\bm e=2\bm e_q$, so their small-error behavior can be compared directly
(Table~\ref{tab:gains}): every family reduces near the reference to a second-order loop at the same bandwidth scale, with
damping ratios from $0.71$ (LQR, MPC) to $1$ (SMC). Each gain is read from the class the runs used, and each closed-form
derivation is checked numerically.\ifwithbs{} The family the code calls ``backstepping'' is computed torque; Section~\ref{sec:bstrue}
therefore adds a recursive Lyapunov backstepping design (plain and adaptive) in the style of the EHEKATL BWB controllers,
with its stability derivation, numerical verification, and a 40-run campaign with and without the residual. This new
family is not part of the manuscript.\fi{} The source files cited here are byte-identical to the ones recorded by every A5 run
(\fhashMatch/\fhashChecked{} hashes).
\end{abstract}

\tableofcontents
\newpage

% =============================================================================
\section{How to read this document}
% =============================================================================
\begin{itemize}[leftmargin=1.2em]
\item \textbf{Paper references.} ``\Peq{15}'' is equation (15) of the manuscript; Sec.~VI-B and Proposition~1 refer to the
  manuscript. The equation-to-code companion of September 30 (\code{paper/companion_2026-09-30}) covers the quaternion
  conventions, the residual and the gate in detail; this document does not repeat them.
\item \textbf{Code.} Every listing is read from the repository file at the quoted lines, with the file's own line numbers.
  Four files carry all six families: \BL{} (class \code{AdvancedBaseline}: backstepping, LQR, SMC, MPC, NMPC), \RG{}
  (the A5 plant loop and the PID path), and the campaign drivers \code{run_campaign_a5.py} and
  \code{run_campaign_a5_nmpc_corrected.py}. Appendix~\ref{app:hashes} gives their hashes.
\item \textbf{Status tags.} \emph{Standard}: textbook result, cited. \emph{Derived here}: the steps are written out.
  \emph{Verified}: recomputed by \code{tools/build_family_tables.py} from the code itself. \emph{Implementation note}: a
  detail of the code that the paper leaves implicit, with its effect on reported numbers.
\item \textbf{Figures.} Every architecture uses the manuscript's figure style. The left part (reference and attitude error)
  and the right part (torque box, gated residual, allocation, plant) are identical in all six figures; the dashed blue group is
  the only part that changes from family to family.
\end{itemize}

% =============================================================================
\section{What all six families share}\label{sec:common}
% =============================================================================
\subsection{The residual attaches to any nominal torque}
The residual reads only the nominal request $\bm u_0$ and adds a bounded, gated correction to it (Fig.~\ref{fig:interface}).
Proposition~1 assumes only $\bm u_{0,k}\in\mathcal U=\{\bm u:|u_i|\le\ell_i\}$, so it holds for every family whose output
is saturated to the torque box. All six are: LQR, backstepping, SMC and MPC through the shared clip (line~241 of \BL),
NMPC by clipping every candidate (lines 265 and 273), and PID at line~353 of \RG. Exact recovery at $\beta=0$ was certified
by running the augmentor with $\beta=0$ against the classical run: relative difference $0$ in all five cases (LQR, SMC, MPC
and two NMPC missions; \code{certify_beta0_recovery.py}). Backstepping is covered by the 20/20 bit-exact match with the
sealed primary runs; PID has no separate $\beta=0$ run, and for it the recovery rests on Proposition~1 alone.

\begin{figure}[H]\centering
\resizebox{\textwidth}{!}{% How every family attaches to the same gated CMG--PES residual (Sec. 2). Row 1: the six u0 sources and the
% plant path; rows 2-3: the residual path of the manuscript's Fig. 1, unchanged.
\begin{tikzpicture}
  % ---- six u0 sources -----------------------------------------------------------------------------------------
  \foreach \n/\t [count=\i] in {f1/PD+ff (backstepping),f2/PID+ff,f3/LQR,f4/SMC,f5/MPC,f6/NMPC}
    \node[law,minimum height=6.5mm,minimum width=30mm] (\n) at (0,{1.75-0.9*\i}) {\t};
  \ifwithbs
  \node[law,dashed,minimum height=6.5mm,minimum width=30mm,align=center] (f7) at (0,{1.75-0.9*7-0.25}) {Recursive backstepping\\(new, not in paper)};
  \node[fam,fit=(f1)(f7)] (fams) {};
  \else
  \node[fam,fit=(f1)(f6)] (fams) {};
  \fi
  \node[famlbl,anchor=south] at (fams.north) {one family per run (A5)};
  \ifwithbs\def\nfam{7}\else\def\nfam{6}\fi
  \foreach \i in {1,...,\nfam} \draw[thick] (f\i.east) -- (f\i.east-|fams.east);

  % ---- residual path (manuscript Fig. 1) -----------------------------------------------------------------------
  \node[neuro] (enc)  at (4.9,-3.0) {\textbf{Paired currents}\\$\bar{\bm u}=\bm u_0\oslash\bm\ell$\\
        $\bm I^\pm=I_b\bm1+I_g\max(\pm\bar{\bm u},\bm0)$};
  \node[neuro] (cmg)  [right=7mm of enc] {\textbf{Six CMG units}\\$(v_j,r_j)$, no spikes\\$\bm a=\clip_{[0,1]}(G\bm r^+)$};
  \node[neuro] (dec)  [right=7mm of cmg] {\textbf{Decoder}\\$\bm s=\clip_{[-1,1]}(W\bm a)$};
  \node[clas]  (gate) [right=8mm of dec] {\textbf{Saturation-aware gate}\\$\tilde s_i=0$ if $|u_{0,i}|\ge\ell_i$ and $s_iu_{0,i}<0$\\
        else $\tilde s_i=s_i$};
  \node[neuro] (pes)  [below=11mm of cmg.south west,anchor=north west] {\textbf{PES} (online, every step)\\
        $W^+=\Pi_{\mathcal W}\big[(1-\lambda_W)W+\eta\bm\varepsilon\bm a\T\big]$,\ \ $\bm\varepsilon=\bm y^*-\bm s$};
  \node[neuro] (tgt)  [left=7mm of pes] {\textbf{PES teaching target}\\$\bm y^*=\clip_{[-1,1]}(-k_e\bm e_q)$\\from the measured attitude};

  % ---- plant path ----------------------------------------------------------------------------------------------
  \node[clas]  (sel)   at (enc|-0,0) {\textbf{Nominal torque}\\$\bm u_0=\pi_{\rm fam}(q,\bm\omega,t)\in\mathcal U$};
  \node[clas]  (box)   at (gate|-0,0) {\textbf{Bounded torque addition}\\$\bm u=\sat_{\bm\ell}(\bm u_0+\beta\bm\ell\odot\tilde{\bm s})$};
  \node[plant] (plant) [right=8mm of box] {\textbf{Allocation and plant}\\$\bm f=\clip(B^{-1}[T,\bm u\T]\T)$\\MuJoCo X2 / CF2};
  \draw[sig] (fams.east|-sel) -- (sel);
  \draw[sig] (sel) -- node[lbl,above]{$\bm u_0$} (box);
  \draw[sig] (box) -- node[lbl,above]{$\bm u$} (plant);

  % ---- signals ---------------------------------------------------------------------------------------------------
  \draw[sig] (sel) -- node[lbl,right]{$\bm u_0$} (enc);
  \draw[sig] (enc) -- node[lbl,above]{$\bm I$} (cmg);
  \draw[sig] (cmg) -- node[lbl,above]{$\bm a$} (dec);
  \draw[sig] (dec) -- node[lbl,above]{$\bm s$} (gate);
  \draw[sig] (gate) -- node[lbl,right]{$\tilde{\bm s}$} (box);
  \draw[sig] (cmg.south) -- node[lbl,left]{$\bm a$} (cmg.south|-pes.north);
  \draw[sig] ($(pes.north-|dec)+(-3mm,0)$) -- node[lbl,left]{$W$} ($(dec.south)+(-3mm,0)$);
  \draw[sig] ($(dec.south)+(3mm,0)$) -- node[lbl,right]{$\bm s$} ($(pes.north-|dec)+(3mm,0)$);
  \draw[sig] (tgt) -- node[lbl,above]{$\bm y^*$} (pes);
  \coordinate (u0tap) at ($(sel.east)+(6mm,0)$);
  \fill (u0tap) circle (1.1pt);
  \draw[sigd] (u0tap) -- ++(0,-1.25) -| node[lbl,pos=.25,above]{$\bm u_0$ (load test of the gate)} ($(gate.north)+(-12mm,0)$);

  \begin{scope}[on background layer]
    \node[fam,draw=black!55,fit=(enc)(cmg)(dec)(gate)(pes)(tgt),inner sep=9pt] (rp) {};
  \end{scope}
  \node[note,anchor=north east] at (rp.south east) {CMG--PES residual path of the manuscript, unchanged across families};
\end{tikzpicture}
}
\caption{One residual, six nominal controllers. The nominal torque comes from one family per run; the residual path is the
manuscript's Fig.~1 unchanged ($\beta=.5$ on X2, $.2$ on CF2; $\eta=8\times10^{-4}$; $k_e=8$). The gate and the outer saturation
see only $\bm u_0$, so nothing in the residual depends on how $\bm u_0$ was computed.}
\label{fig:interface}
\end{figure}

\subsection{Error coordinates and the design model}\label{sec:design}
All families act on the shortest-path error $\bm e_q=\sgnp(q_{e,0})\bm q_{e,v}$ (\Peq{4}--\Peq{6}) and the rate error
$\bm\omega_e=\bm\omega-\bm\omega_d$, with $\bm\omega_d$ the central-difference reference rate of the companion, Sec.~2.5.
\code{AdvancedBaseline} works in the rotation-vector coordinate
\begin{equation}
\bm e=2\bm e_q,\qquad \|\bm e\|=2\sin\tfrac{\theta_q}{2}=\theta_q+O(\theta_q^3),
\label{eq:e}
\end{equation}
and designs every law for the per-axis double integrator obtained near the reference. From the error kinematics
\Peq{7}--\Peq{8}, $\dot{\bm e}=q_{e,0}\bm\omega_e+\bm e_q\times\bm\omega_e\approx\bm\omega_e$; from Euler's equation \Peq{2} with
the principal moments $\bm J=\bm j_{\rm tune}$ (Table~I of the paper),
\begin{equation}
\dot e_i\approx\omega_{e,i},\qquad \dot\omega_{e,i}=a_i-\dot\omega_{d,i}-\big[(\bm\omega\times\bm J\bm\omega)\oslash\bm J\big]_i,\qquad
a_i=\tau_i/J_i .
\label{eq:design}
\end{equation}
Every family outputs a desired angular acceleration $\bm a$; the torque is $\bm\tau=\bm J\odot\bm a$ plus, optionally, the
gyroscopic term, saturated to the box:
\begin{equation}
\bm u_0=\sat_{\bm\ell}\big(\bm J\odot\bm a+\chi_{\rm ff}\,\bm\omega\times(\bm J\odot\bm\omega)\big),\qquad
\bm a=\chi_{\rm ff}\dot{\bm\omega}_d+\bm a_{\rm fam}(\bm e,\bm\omega_e).
\label{eq:interface}
\end{equation}
\emph{Derived here.} In every run of the paper $\chi_{\rm ff}=0$ (\code{--no-model-ff}): the acceleration feedforward and the
gyroscopic term are off, so each family is a feedback law on $(\bm e,\bm\omega_e)$ scaled by inertia, and the residual is the
only adaptive element. The rate feedforward $\bm\omega_d$ stays on (\code{--rate-ff left-neg}) for all families.

\listing{experiments/nengo_nef_replication_2026-08-29/run_mujoco_nef_baselines.py}{189}{199}{$\bm\omega_d$ and $\dot{\bm\omega}_d$ (the latter is zero when model feedforward is off)}
\listing{experiments/nengo_nef_replication_2026-08-29/run_mujoco_nef_baselines.py}{216}{245}{\code{AdvancedBaseline.torque}: one interface, five families, shared clip at line 241}

\subsection{Torque box and platform constants}
$\ell_i=J_i\nu_i$ with $\bm\nu=(12,12,10)$\,rad/s$^2$, the same for both platforms, so the box limits every family to the
same angular acceleration $|a_i|\le\nu_i$:

\begin{table}[H]\centering\small
\caption{Constants read from the A5 run manifests (\ccode{derived} field).}\label{tab:platform}
\begin{tabular}{@{}lll@{}}\toprule
& X2 & CF2 (1.5$\times$ rotor design)\\\midrule
$\bm J$ (kg\,m$^2$) & \fJXtwo & \fJCFtwo\\
$\bm\ell$ (N\,m) & \fellXtwo & \fellCFtwo\\
$\bm\ell\oslash\bm J$ (rad/s$^2$) & \fnuXtwo & \fnuCFtwo\\
$\rho\bm\omega_n$ (rad/s) & 12, 12, 10.5 & 12, 12, 10.5\\\bottomrule
\end{tabular}\end{table}

Because the box is defined as an acceleration, every gain in this document is the same for X2 and CF2 (the script asserts
it), and the torque differs only by the inertia factor.

\subsection{The A5 settings}
All six families ran with the same plant, missions, residual, gate, $\rho=3$ bandwidth scale and seed; only the family flags
differ:
\listing{experiments/final_robustness_gating_2026-09-20/run_campaign_a5.py}{27}{40}{per-family flags}
\listing{experiments/final_robustness_gating_2026-09-20/run_campaign_a5.py}{49}{56}{common flags: $\rho=3$, rate feedforward on, model feedforward off, gate on in the CMG arm}

\implnote{The NMPC row of the paper comes from \code{run_campaign_a5_nmpc_corrected.py}, which drops the
\code{--baseline-rate-hz 100} hold of the first A5 run so that NMPC, like every other family, updates at every $0.5$-ms step
(Sec.~\ref{sec:nmpc}). The held run is kept on disk and not used.}

\subsection{Small-error equivalence of the six families}\label{sec:equiv}
Near the reference, inside the SMC boundary layer and away from the box, every family reduces to
$a_i=-k_{1,i}e_i-k_{2,i}\omega_{e,i}\,(-k_{I,i}\!\int\! e_i)$, i.e.\ the characteristic polynomial $s^2+k_2s+k_1$ (or
$s^3+k_2s^2+k_1s+k_I$ for PID). Table~\ref{tab:gains} lists the gains each family actually uses; Table~\ref{tab:poles} gives
the resulting poles of the linearized design model \eqref{eq:design}.

\begin{table}[H]\centering\small
\caption{Equivalent small-error gains, read from \ccode{AdvancedBaseline} with the A5 settings (roll\,/\,yaw; pitch equals roll).
$\omega_{\rm eq}=\sqrt{k_1}$, $\zeta_{\rm eq}=k_2/(2\sqrt{k_1})$. ``Box at'' is the rotation angle at which the
proportional term alone reaches the torque box, $\nu_i/k_{1,i}$. \emph{Verified.}}\label{tab:gains}
\begin{tabular}{@{}lllllll@{}}\toprule
Family & $k_1$ (s$^{-2}$) & $k_2$ (s$^{-1}$) & $\omega_{\rm eq}$ (rad/s) & $\zeta_{\rm eq}$ & box at ($^\circ$) & $k_I$ (s$^{-3}$)\\\midrule
\tabinput{generated/family_gains.tex}
\bottomrule\end{tabular}\end{table}

\begin{table}[H]\centering\small
\caption{Closed-loop poles of the linearized design model \eqref{eq:design} with each family's gains. \emph{Verified.}}\label{tab:poles}
\begin{tabular}{@{}lll@{}}\toprule
Family & axis & poles (s$^{-1}$)\\\midrule
\tabinput{generated/family_poles.tex}
\bottomrule\end{tabular}\end{table}

The families therefore differ in damping and in what happens away from the origin: how each one behaves at the box (SMC
uses a bounded reaching term, PID freezes its integrator, NMPC searches around the nominal law), not in bandwidth. The
proportional term of every linear law reaches the box at a rotation error of $5$--$6^\circ$, so Hardstep's steps of tens of
degrees saturate every family, which is where the gate matters.

% =============================================================================
\section{PD with desired-rate feedforward in computed-torque form (primary baseline)}\label{sec:pdff}
In the code this family is called \code{backstepping} (\code{--baseline backstepping}); it is the computed-torque form.\ifwithbs{}
It is not a recursive backstepping design; the recursive design is Sec.~\ref{sec:bstrue}.\fi
% =============================================================================
\archfig{fig_pdff.tex}{PD+ff in computed-torque (backstepping) form, the paper's $\rho=3$ primary baseline. The model
feedforward path (dotted) exists in the code and is switched off in every run, which makes this law identical to \Peq{15}.}{fig:pdff}

\paragraph{Law.} Computed-torque backstepping \cite{krstic1995nonlinear,slotine1991applied}: choose the error dynamics
$\dot{\bm\omega}_e=-k_1\bm e-k_2\bm\omega_e$ and invert \eqref{eq:design},
\begin{equation}
\bm u_0=\sat_{\bm\ell}\Big(\bm J\odot\big(\dot{\bm\omega}_d-k_1\bm e-k_2\bm\omega_e\big)+\bm\omega\times(\bm J\odot\bm\omega)\Big),\qquad
k_1=(\rho\omega_n)^2,\quad k_2=2\zeta\rho\omega_n,\ \zeta=0.95 .
\label{eq:bs}
\end{equation}
With the feedforward on, no saturation and a diagonal inertia, the gyroscopic term and $\dot{\bm\omega}_d$ cancel and the
rate error obeys $\dot{\bm\omega}_e=-k_1\bm e-k_2\bm\omega_e$ exactly; with $\dot{\bm e}\approx\bm\omega_e$ this is
$\ddot e_i+k_2\dot e_i+k_1e_i\approx0$ near the reference, the textbook reason for this form. \emph{Standard.}

\paragraph{Reduction to the paper's PD+ff.} With $\chi_{\rm ff}=0$ and $\bm e=2\bm e_q$,
$\bm u_0=\sat_{\bm\ell}(-2\bm J(\rho\bm\omega_n)^2\odot\bm e_q-2\zeta\bm J\rho\bm\omega_n\odot\bm\omega_e)$, which is \Peq{13}--\Peq{15}
with $\bm\omega_n\mapsto\rho\bm\omega_n$. \emph{Derived here; verified bit-exact in 20/20 cells}
(\code{run_campaign_a5_backstepping.py}, relative deviation $0$ against the sealed primary runs).

\listing{experiments/nengo_nef_replication_2026-08-29/run_mujoco_nef_baselines.py}{138}{141}{gains $k_1=\omega^2$, $k_2=2\zeta\omega$, $\omega=\rho\omega_n$}
\listing{experiments/final_robustness_gating_2026-09-20/run_controllers_gated.py}{79}{86}{the frozen paper baseline: backstepping, $\rho=3$, $\zeta=.95$, \code{left-neg}, no model feedforward}

\ifwithbs
% =============================================================================
\section{Recursive Lyapunov backstepping (new family)}\label{sec:bstrue}
% =============================================================================
The family of Sec.~\ref{sec:pdff} is computed torque: it chooses linear error dynamics and inverts the model, and with
model feedforward off it is exactly PD+ff. This section adds a backstepping controller in the strict sense: a recursive
design with a virtual control on the kinematic subsystem, a cross-coupling term that cancels the interconnection in the
Lyapunov derivative, and a final Lyapunov function that certifies the closed loop
\cite{krstic1995nonlinear,khalil2002nonlinear}. It follows step by step the design used for the longitudinal and adaptive
controllers of the EHEKATL BWB (strict-feedback form, errors $z_1,z_2$, virtual control with the term $-g_1\T z_1$,
$V=\sum\frac12z_i\T z_i$, La Salle--Yoshizawa, adaptive estimates with $\frac12\tilde f\T\Gamma\tilde f$). The plain and the
adaptive versions are implemented in \code{experiments/backstepping_true_2026-10-01/true_backstepping.py} and run through the
unchanged A5 runner. \textbf{They are not part of the manuscript.}

\archfig{fig_bstrue.tex}{Recursive backstepping. Step~1 stabilizes the quaternion kinematics with a virtual rate; step~2 drives
the rate to it with a torque that cancels the gyroscopic term, feeds forward the transported desired-rate derivative, and
contains the coupling term $-c\,z_1$. The adaptive variant adds a projected estimate of an unknown body torque.}{fig:bstrue}

\subsection{Strict-feedback form}
With the passive convention of the paper, $q_e=q_d\otimes q^{-1}$ and its body-frame rate
$\bm\omega_e=\bm\omega-\bar{\bm\omega}_d$, $\bar{\bm\omega}_d=R(q_e)\T\bm\omega_d$ (\Peq{7}--\Peq{8}, with the transport
$R(q_e)\T$ kept), the attitude system is in strict-feedback form:
\begin{equation}
\dot z_1=\tfrac12T(q_e)\,(\bm\omega-\bar{\bm\omega}_d),\qquad
\bm J\dot{\bm\omega}=-\bm\omega\times\bm J\bm\omega+\bm\tau+\bm d,\qquad
z_1=\bm e_q=\sgnp(q_{e,0})\bm q_{e,v},\quad T(q_e)=|q_{e,0}|I+[z_1]_\times ,
\label{eq:sf}
\end{equation}
i.e.\ $x_1=z_1$, $x_2=\bm\omega$, $f_1=-\frac12T\bar{\bm\omega}_d$, $g_1=\frac12T$, $f_2=-\bm J^{-1}(\bm\omega\times\bm J\bm\omega)$,
$g_2=\bm J^{-1}$ in the thesis notation, with $\bm d$ an unknown body torque (zero in the plain design). There is no actuator
state: the rotor forces of the MuJoCo models act without lag, so the thesis's $z_3$ step is not needed. The first identity
follows from \Peq{8}: $\frac{d}{dt}(\sgnp(q_{e,0})\bm q_{e,v})=\frac12(|q_{e,0}|\bm\omega_e+z_1\times\bm\omega_e)$.

\subsection{Step 1: kinematic subsystem}
Take $V_1=2(1-|q_{e,0}|)\ge0$, zero exactly at $q_e=\pm1$. With $\dot q_{e,0}=-\frac12\bm q_{e,v}\T\bm\omega_e$ (\Peq{8}),
\[
\dot V_1=-2\,\sgnp(q_{e,0})\,\dot q_{e,0}=z_1\T\bm\omega_e .
\]
Treat $\bm\omega_e$ as the control of this subsystem and choose the virtual rate $\bm\omega_e^d=-K_1z_1$, $K_1=\operatorname{diag}(k_{1,i})>0$:
if $\bm\omega_e=\bm\omega_e^d$, then $\dot V_1=-z_1\T K_1z_1$.

\subsection{Step 2: rate subsystem}
Define the second error $z_2=\bm\omega_e-\bm\omega_e^d=\bm\omega_e+K_1z_1$ and $V=c\,V_1+\frac12z_2\T\bm Jz_2$, $c>0$. Using
$\bm\omega_e=z_2-K_1z_1$ and $\dot{\bm\omega}_e=\dot{\bm\omega}-\dot{\bar{\bm\omega}}_d$,
\[
\dot V=c\,z_1\T z_2-c\,z_1\T K_1z_1+z_2\T\big(-\bm\omega\times\bm J\bm\omega+\bm\tau+\bm d-\bm J\dot{\bar{\bm\omega}}_d+\bm JK_1\dot z_1\big).
\]
The control
\begin{equation}
\boxed{\ \bm\tau=\bm\omega\times\bm J\bm\omega+\bm J\big(\dot{\bar{\bm\omega}}_d-K_1\dot z_1\big)-c\,z_1-K_2z_2-\hat{\bm d}\ },\qquad
\bm u_0=\sat_{\bm\ell}(\bm\tau),
\label{eq:bstrue}
\end{equation}
with $K_2=\operatorname{diag}(k_{2,i})>0$ and $\dot z_1=\frac12T(q_e)\bm\omega_e$ computed from measured signals, gives
\begin{equation}
\dot V=-c\,z_1\T K_1z_1-z_2\T K_2z_2+z_2\T\tilde{\bm d},\qquad \tilde{\bm d}=\bm d-\hat{\bm d}.
\label{eq:vdot}
\end{equation}
The term $-c\,z_1$ is the backstepping coupling term: it cancels $c\,z_1\T z_2$, the interconnection between the two
subsystems ($-g_1\T z_1$ in the thesis). Without it the design is computed torque.

\paragraph{Plain law} ($\bm d=\hat{\bm d}=\bm0$). $\dot V=-c\,z_1\T K_1z_1-z_2\T K_2z_2<0$ for $(z_1,z_2)\ne0$; by La Salle--Yoshizawa
$z_1,z_2\to0$, hence $q_e\to\pm1$ and $\bm\omega\to\bar{\bm\omega}_d$, and $z_1,z_2\in\mathcal L_2$. Because
$c\,V_1\ge c\|z_1\|^2/2$ near the origin, the convergence is exponential there.

\paragraph{Adaptive law.} For a constant unknown $\bm d$, add $\frac12\tilde{\bm d}\T\Gamma^{-1}\tilde{\bm d}$ to $V$. Since
$\dot{\tilde{\bm d}}=-\dot{\hat{\bm d}}$, the update
\begin{equation}
\dot{\hat{\bm d}}=\operatorname{Proj}_{|\hat d_i|\le D_i}\big(\Gamma z_2\big)
\label{eq:adapt}
\end{equation}
cancels $z_2\T\tilde{\bm d}$ in \eqref{eq:vdot} and leaves $\dot V_a\le-c\,z_1\T K_1z_1-z_2\T K_2z_2$; the projection keeps
$\hat{\bm d}$ bounded without increasing $\dot V_a$ when $|d_i|\le D_i$ \cite{ioannou1996robust,krstic1995nonlinear}. This is the
thesis's $\dot{\tilde f}=-\Gamma^{-1}z$ law with a bound added. Here $\Gamma_i=0.3\cdot2J_i\omega_i^2/k_{1,i}$, which makes the
integral action on $z_1$ ($\Gamma k_1\!\int\!z_1$) equal to the PID family's, and $D_i=0.25\,\ell_i$.

\paragraph{Rate transport.} $\bar{\bm\omega}_d=R(q_e)\T\bm\omega_d$ and, since $q_e$ rotates with body rate $\bm\omega_e$
($\dot R_e=R_e[\bm\omega_e]_\times$),
\[
\dot{\bar{\bm\omega}}_d=R(q_e)\T\dot{\bm\omega}_d-\bm\omega_e\times\bar{\bm\omega}_d ,
\]
with $\bm\omega_d$, $\dot{\bm\omega}_d$ the same central differences as Sec.~\ref{sec:design}. Unlike the A5 families, this law
uses the transport, the gyroscopic term and $\dot{\bm\omega}_d$ by construction; the \code{--no-model-ff} flag does not apply.

\subsection{Gain design: same linearization as PD+ff}
About $z=0$, \eqref{eq:sf}--\eqref{eq:bstrue} give, per axis, $\ddot z_{1,i}+(k_{1,i}/2+k_{2,i}/J_i)\dot z_{1,i}+(c+k_{1,i}k_{2,i})/(2J_i)\,z_{1,i}=0$.
Writing $k_{1,i}=2\lambda_{1,i}$, $k_{2,i}=J_i\lambda_{2,i}$ and imposing
\[
\lambda_{1,i}+\lambda_{2,i}=2\zeta\omega_i,\qquad \lambda_{1,i}\lambda_{2,i}+\frac{c}{2J_i}=\omega_i^2,\qquad \omega_i=\rho\omega_{n,i},\ \zeta=0.95,
\]
makes the small-error poles exactly those of the paper's PD+ff ($-11.40\pm3.75j$ on roll). Any difference between the two
families is then due to the backstepping structure, not to bandwidth or damping. $V_1$ is a scalar, so one $c$ serves all
axes; real $\lambda$'s need $c\ge2J_i\omega_i^2(1-\zeta^2)$ on every axis, and the smallest such value is used
($\lambda_1\le\lambda_2$: the outer kinematic loop is the slower one). For X2: $c=\fbsC$\,N\,m,
$\bm k_1=(\fbsKone)$\,s$^{-1}$, $\bm k_2=(\fbsKtwo)$\,N\,m\,s.

\subsection{Verification}
\code{verify_true_backstepping.py} checks the derivation on the rigid-body model (\code{results/verification.json}). \emph{Verified}:
\begin{itemize}[leftmargin=1.2em]
\item $\dot z_1$ and $\dot{\bar{\bm\omega}}_d$ in closed form against central finite differences along rigid-body trajectories on Full3,
  Hard45 and Yaw360: worst absolute error \fbsKinErr{} (signals of order $0.1$--$1$);
\item closed loop, plain law, no torque limit, Full3 and Hard45 from a $\approx20^\circ$ initial error: $V$ decreases monotonically to
  $10^{-12}$, and the finite-difference $\dot V$ equals $-c\,z_1\T K_1z_1-z_2\T K_2z_2$ within \fbsVdotErr{} (relative, step-size limited);
\item adaptive law with a constant unknown torque of $20\%$ of the box: $\hat{\bm d}$ reaches $\bm d$ within \fbsAdaptErr{} after $20$\,s;
\item linearized poles equal the PD+ff poles to \fbsPoleErr.
\end{itemize}

\listing{experiments/backstepping_true_2026-10-01/true_backstepping.py}{44}{55}{gain design: smallest common $c$, $\lambda_1\le\lambda_2$}
\listing{experiments/backstepping_true_2026-10-01/true_backstepping.py}{84}{114}{$z_1$, $z_2$, $\dot z_1$, transport, law \eqref{eq:bstrue}, projected adaptation \eqref{eq:adapt}}
\listing{experiments/backstepping_true_2026-10-01/run_true_backstepping.py}{30}{42}{the family enters the unchanged A5 runner}

\implnote{The Lyapunov certificate holds for the model \eqref{eq:sf} without saturation. In the plant the torque box, rotor
clipping, the non-diagonal X2 inertia (the law uses the principal moments, as every family does) and the residual add terms
that the certificate does not cover; the sign selection makes $V_1$ non-smooth at $q_{e,0}=0$, the same local scope as the
paper. The output is saturated to the box, so Proposition~1 applies unchanged.}

\subsection{Results: recursive backstepping with and without the residual}
Same protocol as A5 (Table~\ref{tab:results}): 45-s missions, $\rho=3$, noise-free, seed 1, gate on, the residual unchanged. Run by \code{experiments/backstepping_true_2026-10-01/run_campaign_true_bs.py}; numbers computed from the run summaries by \code{tools/build_family_tables.py}.
\begin{table}[H]\centering\small
\caption{Reduction of mean quaternion-angle error by the gated CMG--PES residual on the recursive backstepping laws (\%, positive is better; 5-mission = reduction of geometric means; last column: baseline geometric mean, mrad).}\label{tab:bsres}
\begin{tabular}{@{}llrrrrrrr@{}}\toprule
Family & Platform & Roll & Full3 & Hard45 & Hardstep & Yaw360 & 5-mission & base (mrad)\\\midrule
Recursive backstepping & X2 & -114.8 & +7.2 & +11.8 & -2.9 & +43.5 & \textbf{-0.4} & 3.30\\
 & CF2 & -1416.5 & -18.2 & -38.1 & -0.8 & -1467.9 & \textbf{-230.0} & 0.59\\\addlinespace
Adaptive recursive backstepping & X2 & -114.7 & +6.5 & +11.7 & -0.8 & +30.9 & \textbf{-4.3} & 2.63\\
 & CF2 & -669.0 & -9.6 & -19.6 & -0.4 & -2494.5 & \textbf{-204.7} & 0.88\\\addlinespace
\bottomrule\end{tabular}\end{table}
\begin{table}[H]\centering\small
\caption{Nominal laws alone (no residual): error reduction of the recursive designs relative to the computed-torque PD+ff (the paper's primary baseline), \% (positive: the recursive law tracks better). The linearizations are identical, so these differences come from the backstepping structure.}\label{tab:bsvsct}
\begin{tabular}{@{}llrrrrrr@{}}\toprule
Platform & law & Roll & Full3 & Hard45 & Hardstep & Yaw360 & 5-mission\\\midrule
X2 & recursive & +29.6 & +30.5 & +31.0 & -2.4 & +41.5 & \textbf{+27.4}\\
X2 & adaptive & +36.5 & +30.8 & +29.1 & -7.9 & +80.7 & \textbf{+42.2}\\
CF2 & recursive & +95.0 & +64.8 & +73.3 & -0.5 & +97.9 & \textbf{+84.2}\\
CF2 & adaptive & +91.2 & +47.5 & +58.9 & -695.8 & +99.6 & \textbf{+76.7}\\
\bottomrule\end{tabular}\end{table}


\paragraph{What the two tables show.}
\begin{itemize}[leftmargin=1.2em]
\item \textbf{The recursive law is a much better nominal controller} (Table~\ref{tab:bsvsct}): with the same linearization it
  reduces the computed-torque PD+ff error on four of five missions on both platforms, most strongly on CF2. Its law contains
  $\bm J\dot{\bar{\bm\omega}}_d$ and $\bm\omega\times\bm J\bm\omega$, which every A5 family ran without (\code{--no-model-ff}). The
  computed-torque law \emph{with} model feedforward was not run, so this campaign does not separate the share of the gain due to the
  acceleration feedforward from the share due to the recursive structure (cross term and rate transport). That run is the
  natural next control.
\item \textbf{The unchanged residual does not help it} (Table~\ref{tab:bsres}). It still reduces the error on X2 Full3, Hard45
  and Yaw360, but where the recursive law's error is already very small (Roll; CF2 Roll and Yaw360, a few hundredths of a mrad) it
  raises the error to tenths of a mrad, worse in absolute terms than the residual on the paper's baseline. The residual's
  $\beta$, $\eta$ and $k_e$ were set for a lagging PD+ff; on a nominal law that already cancels the dominant error, the residual's
  own correction becomes the error. Retuning ($\beta$ in particular) was not tested. This is a limit of the unchanged residual,
  not of Proposition~1, which still holds: the output stays in the box and $\beta=0$ recovers the law.
\item \textbf{The adaptive variant fails on CF2 Hardstep} under persistent rotor saturation (79\% of samples): the projected estimate
  keeps integrating $z_2$ while the torque it requests cannot be delivered, a windup through the adaptation law that the
  Lyapunov argument (which assumes no saturation) does not exclude. Freezing the adaptation when the box or a rotor saturates, the
  rule the PID family uses, is the standard fix and was not run. On the other missions the adaptive and plain laws are close.
\end{itemize}
\fi % withbs

% =============================================================================
\section{PID with rate feedforward and rotor-level anti-windup}\label{sec:pid}
% =============================================================================
\archfig{fig_pid.tex}{PID+ff. The integrator runs on $\bm e_q$ and is frozen per axis whenever the torque box clipped that axis,
and on all axes whenever any rotor saturated at the previous step (dashed paths).}{fig:pid}

\paragraph{Law.} The PID path bypasses \code{AdvancedBaseline} (\code{--baseline pd}) and adds an integral term to \Peq{15}:
\begin{equation}
\bm u_0=\sat_{\bm\ell}\big(-\bm k_q\odot\bm e_q-\bm k_\omega\odot\bm\omega_e-\bm k_i\odot\bm z\big),\qquad
\bm k_q=2\bm J\odot(\rho\bm\omega_n)^2,\ \bm k_\omega=2\zeta\bm J\odot\rho\bm\omega_n,\ \bm k_i=0.3\,\bm k_q .
\label{eq:pid}
\end{equation}
\paragraph{Conditional integration.} With $h=5\times10^{-4}$\,s and the clip test of line 357,
\begin{equation}
z_{i,k+1}=z_{i,k}+h\,e_{q,i,k}\,m_{i,k},\qquad
m_{i,k}=\mathbf 1\big[u_{0,i,k}=\tau_{{\rm u},i,k}\big]\cdot\mathbf 1\big[\text{no rotor saturated at }k-1\big].
\label{eq:aw}
\end{equation}
This is conditional integration (integrator clamping) \cite{astrom2006advancedpid}, applied at the rotor level as well as at the
torque box: a rotor at its force limit means the realized torque differs from the command even if the box is inactive.
\emph{Standard; the rotor-level rule is this study's choice.}

\paragraph{Loop.} In the coordinate $\bm e$ the gains per unit inertia are $k_1=(\rho\omega_n)^2$, $k_2=2\zeta\rho\omega_n$ and
$k_I=0.3(\rho\omega_n)^2$ (Table~\ref{tab:gains}). The Routh condition for $s^3+k_2s^2+k_1s+k_I$ is $k_2k_1>k_I$, met with
a ratio of at least \fpidRouthMargin. The integral adds a slow real pole at $-0.32$\,s$^{-1}$ (time constant $\approx3$\,s,
Table~\ref{tab:poles}): the integrator is weak, which the A5 report discloses as ``undertuned''. \emph{Derived here; verified.}

\listing{experiments/final_robustness_gating_2026-09-20/run_controllers_gated.py}{253}{260}{gains; PID uses the plain PD path}
\listing{experiments/final_robustness_gating_2026-09-20/run_controllers_gated.py}{347}{361}{\eqref{eq:pid} and the conditional integrator \eqref{eq:aw}}
\listing{experiments/final_robustness_gating_2026-09-20/run_controllers_gated.py}{475}{475}{rotor-saturation flag for the next step}

\implnote{The integrator acts on $\bm e_q$ (not $\bm e=2\bm e_q$), so $\bm k_i=0.3\bm k_q$ means
$k_I=0.3(\rho\omega_n)^2$ in Table~\ref{tab:gains}'s coordinate. With \code{--ki-factor 0} this path is \Peq{15} exactly.}

% =============================================================================
\section{LQR}\label{sec:lqr}
% =============================================================================
\archfig{fig_lqr.tex}{LQR. The gain is computed once from the closed-form solution of the continuous algebraic Riccati
equation (CARE) for the per-axis double integrator; online it is a fixed state feedback.}{fig:lqr}

\paragraph{Design problem.} For each axis, $\bm x=[e,\omega_e]\T$, $\dot{\bm x}=A\bm x+Ba$ with
$A=\left[\begin{smallmatrix}0&1\\0&0\end{smallmatrix}\right]$, $B=[0,1]\T$, minimize
$\int_0^\infty(q_e e^2+ra^2)\,dt$ \cite{andersonmoore1990optimal}.

\paragraph{Closed-form CARE.} With $P=\left[\begin{smallmatrix}p_1&p_2\\p_2&p_3\end{smallmatrix}\right]$, the CARE
$A\T P+PA-PBr^{-1}B\T P+Q=0$ reads, entry by entry,
\[
(1,1):\ q_e-\frac{p_2^2}{r}=0,\qquad (1,2):\ p_1-\frac{p_2p_3}{r}=0,\qquad (2,2):\ 2p_2-\frac{p_3^2}{r}=0 .
\]
The positive-definite solution has $p_2=\sqrt{q_er}$, $p_3=\sqrt{2p_2r}$, $p_1=p_2p_3/r$, so
\begin{equation}
K=r^{-1}B\T P=\Big[\sqrt{q_e/r},\ \sqrt{2\sqrt{q_e/r}}\Big]\quad\xrightarrow{\ q_e/r=\omega^4\ }\quad K=[\omega^2,\ \sqrt2\,\omega],
\label{eq:lqr}
\end{equation}
a closed loop $s^2+\sqrt2\omega s+\omega^2$ with $\zeta=1/\sqrt2$ at the same $\omega=\rho\omega_n$ as the other families.
\emph{Derived here. Verified}: $P$ from these formulas leaves a CARE residual of \flqrCareRelErr{} (relative), its smallest
eigenvalue is \flqrPminEig{} $>0$, and $B\T P$ equals the code's gain.

\listing{experiments/nengo_nef_replication_2026-08-29/run_mujoco_nef_baselines.py}{132}{137}{$K=[\omega^2,\sqrt2\omega]$}
\listing{experiments/nengo_nef_replication_2026-08-29/run_mujoco_nef_baselines.py}{232}{233}{the law: $\bm a=\dot{\bm\omega}_d-k_1\bm e-k_2\bm\omega_e$}

% =============================================================================
\section{Sliding-mode control}\label{sec:smc}
% =============================================================================
\archfig{fig_smc.tex}{SMC. A linear sliding variable and a boundary-layer reaching law; inside the layer the law is a critically
damped PD, outside it the reaching term has a fixed magnitude of half the torque box.}{fig:smc}

\paragraph{Law.} Sliding variable and reaching law \cite{utkin1992sliding,slotine1991applied}, per axis:
\begin{equation}
\sigma_i=\omega_{e,i}+\lambda_ie_i,\qquad
a_i=\dot\omega_{d,i}-\lambda_i\omega_{e,i}-\eta_i\,\sat\!\big(\sigma_i/\phi_i\big),\qquad
\lambda_i=\rho\omega_{n,i},\ \ \eta_i=0.5\,\ell_i/J_i,\ \ \phi_i=\eta_i/\lambda_i .
\label{eq:smc}
\end{equation}
With the A5 setting, $\bm\eta=(\fsmcEtaXtwo)$\,rad/s$^2$ and $\bm\phi=(\fsmcPhiXtwo)$\,rad/s on both platforms.

\paragraph{Inside the boundary layer} ($|\sigma_i|<\phi_i$): $a_i=-\lambda_i\omega_{e,i}-(\eta_i/\phi_i)(\omega_{e,i}+\lambda_ie_i)$, and
$\eta_i/\phi_i=\lambda_i$, so
\[
a_i=-\lambda_i^2e_i-2\lambda_i\omega_{e,i}:\qquad k_1=\lambda^2,\quad k_2=2\lambda,\quad \zeta=1,
\]
a double pole at $-\rho\omega_n$ (Table~\ref{tab:poles}). \emph{Derived here; verified from the code's
\code{lam}, \code{eta}, \code{phi}.}

\paragraph{Outside the layer.} On the design model \eqref{eq:design} with model feedforward off,
$\dot\sigma_i=\dot\omega_{e,i}+\lambda_i\dot e_i=a_i-\dot\omega_{d,i}+\lambda_i\omega_{e,i}=-\eta_i\operatorname{sgn}(\sigma_i)-\dot\omega_{d,i}$,
hence
\[
\sigma_i\dot\sigma_i\le-(\eta_i-|\dot\omega_{d,i}|)\,|\sigma_i| ,
\]
the reaching condition, whenever the reference acceleration stays below $\eta_i$ ($6$\,rad/s$^2$ on roll and pitch). The layer
is then reached in finite time and the error decays inside it; the boundary layer trades exact sliding for a continuous,
chattering-free command. Because $J_i\eta_i=\ell_i/2$, the reaching term alone uses half the torque box, leaving the
other half for $-\lambda\bm\omega_e$. \emph{Standard argument, specialized here.}

\listing{experiments/nengo_nef_replication_2026-08-29/run_mujoco_nef_baselines.py}{142}{149}{$\lambda=\omega$, $\eta=0.5\,\bm\ell\oslash\bm J$, $\phi=\eta/\omega$}
\listing{experiments/nengo_nef_replication_2026-08-29/run_mujoco_nef_baselines.py}{234}{236}{\eqref{eq:smc}}

% =============================================================================
\section{Model predictive control (unconstrained, finite horizon)}\label{sec:mpc}
% =============================================================================
\archfig{fig_mpc.tex}{MPC. The finite-horizon quadratic program has no constraints, so its solution is linear in the state; the
first move of the optimal sequence is a fixed gain $K_0$, computed once and applied at every plant step.}{fig:mpc}

\paragraph{Prediction model.} Per axis, the exact zero-order-hold discretization of the double integrator with sample time
$T=20$\,ms,
\[
\bm x_{k+1}=A_d\bm x_k+B_da_k,\qquad A_d=\begin{bmatrix}1&T\\0&1\end{bmatrix},\quad B_d=\begin{bmatrix}T^2/2\\T\end{bmatrix}.
\]
\paragraph{Condensed problem.} Stacking $N=120$ predictions, $X=S_x\bm x_0+S_uU$ with
$S_x=[A_d;A_d^2;\dots;A_d^N]$ and the block-Toeplitz $S_u$ with blocks $A_d^{k-j}B_d$ ($j\le k$). The cost
$\sum_{k=1}^{N}\bm x_k\T Q\bm x_k+\sum_{k=0}^{N-1}Ra_k^2$ with $Q=\operatorname{diag}(\omega^4,q_\omega\omega^2)$, $q_\omega=0$, $R=1$ becomes
\[
U\T HU+2\bm x_0\T F\T U+\text{const},\qquad H=S_u\T\bar QS_u+\bar R,\quad F=S_u\T\bar QS_x,
\]
with $\bar Q=I_N\otimes Q$, $\bar R=I_N\otimes R$. Without constraints the minimizer is $U^\star=-H^{-1}F\bm x_0$, and the
receding-horizon law applies its first element, $a_0=-K_0\bm x_0$ with $K_0=[H^{-1}F]_{1,:}$
\cite{borrelli2017predictive,rawlings2017mpc}. \emph{Standard.}

\paragraph{What the horizon does.} Table~\ref{tab:mpc} recomputes $K_0$ with the code's own function for growing $N$. From
$N\approx30$ ($0.6$\,s, several closed-loop time constants) the first-move gain no longer changes: at the A5 horizon it equals
the infinite-horizon discrete LQR gain obtained by iterating the Riccati recursion to its fixed point to a relative
\fmpcVsDare{} ($K_0=[\fmpcKone,\ \fmpcKtwo]$ on roll; DARE residual \fdareResidual). The MPC row is therefore a discrete LQR
designed at $T=20$\,ms. That sample time, not the horizon, is why its proportional gain is $16\%$ below the continuous LQR
of Sec.~\ref{sec:lqr} ($121.5$ vs.\ $144$\,s$^{-2}$). \emph{Verified.}

\begin{table}[H]\centering\small
\caption{First-move MPC gain on the roll axis ($\omega=12$\,rad/s) versus horizon, from \ccode{AdvancedBaseline._mpc_first_gain};
the last two columns divide by the continuous LQR gain $[\omega^2,\sqrt2\omega]$. A5 uses $N=120$.}\label{tab:mpc}
\begin{tabular}{@{}rrrrrr@{}}\toprule
$N$ & $NT$ (s) & $K_{0,1}$ & $K_{0,2}$ & $K_{0,1}/\omega^2$ & $K_{0,2}/(\sqrt2\omega)$\\\midrule
\tabinput{generated/mpc_horizon.tex}
\bottomrule\end{tabular}\end{table}

\listing{experiments/nengo_nef_replication_2026-08-29/run_mujoco_nef_baselines.py}{167}{187}{condensed matrices $S_x,S_u,H,F$ and the first-move gain}
\listing{experiments/nengo_nef_replication_2026-08-29/run_mujoco_nef_baselines.py}{237}{238}{the online law $\bm a=\dot{\bm\omega}_d-K_0[\bm e;\bm\omega_e]$}

\implnote{The gain is designed for a 20-ms move but applied as a continuous-time state feedback at every 0.5-ms step (A5 ran
every family at full rate). This is a common simplification; it is the reason the MPC row behaves like a softer LQR
($\zeta_{\rm eq}=0.71$, $\omega_{\rm eq}=11.0$\,rad/s) and not like a sampled controller. ``MPC'' here means an unconstrained
receding-horizon gain; the torque box is enforced by the shared clip, not inside the optimization.}

% =============================================================================
\section{Nonlinear MPC (grid search around the backstepping law)}\label{sec:nmpc}
% =============================================================================
\archfig{fig_nmpc.tex}{NMPC. At every plant step, 27 candidate offsets are added to the nominal backstepping torque; each
candidate is scored by an RK4 rollout of the closed loop on the rigid-body model over $0.3$\,s, and the first torque of the
cheapest candidate is applied.}{fig:nmpc}

\paragraph{Candidates.} The nominal law is \eqref{eq:bs} ($\zeta=0.95$). A grid of $3^3=\fnmpcCandidates$ constant offsets
$\bm d_c\in\{-1,0,1\}^3\odot0.1\bm\ell$ is added to it; $\bm d_c=\bm0$ is one of them (\fnmpcHasZero), so the nominal law is always
a candidate.

\paragraph{Prediction (``policy'' mode).} From the measured state, each candidate is rolled out for $N=5$ steps of
$\Delta=60$\,ms with the rigid-body model $\dot Q=\frac12Q\otimes\tilde{\bm\omega}$,
$\dot{\bm\omega}=(\bm\tau-\bm\omega\times\bm J\bm\omega)\oslash\bm J$, integrated by classical RK4 \cite{hairer1993ode}. At
$k=0$ the candidate's first torque $\bm\tau^{(c)}_0=\sat_{\bm\ell}(\bm\tau_{\rm nom}+\bm d_c)$ is used; at $k\ge1$ the nominal law is
re-evaluated on the predicted state and the same offset added,
$\bm\tau^{(c)}_k=\sat_{\bm\ell}(\bm J\odot\bm a(\hat q_k,\hat{\bm\omega}_k,t+k\Delta)+\bm d_c)$. The prediction thus models the
closed loop that will actually run, not a constant torque.

\paragraph{Cost and selection.} With $\bm e_k,\bm\omega_{e,k}$ the predicted errors against the reference at $t+(k+1)\Delta$,
\begin{equation}
J_c=\sum_{k=0}^{N-1}\Big[w_k\big(\textstyle\sum_i\omega_i^4e_{k,i}^2+2\sum_i\omega_i^2\omega_{e,k,i}^2\big)+\big\|\bm\tau^{(c)}_k\oslash\bm J\big\|^2\Big],
\qquad w_k=\begin{cases}1.5,&k=N-1\\1,&\text{otherwise}\end{cases}
\label{eq:nmpc}
\end{equation}
and $\bm u_0=\bm\tau^{(c^\star)}_0$, $c^\star=\arg\min_cJ_c$. The angle weight $\omega^4$ is the LQR weight $q_e/r$ of
Sec.~\ref{sec:lqr}; the effort term is the squared commanded acceleration. This is a sampling-based (grid) receding-horizon
controller in the sense of \cite{rawlings2017mpc}, not a gradient-based NMPC solver. \emph{Derived here from the code.}

\listing{experiments/nengo_nef_replication_2026-08-29/run_mujoco_nef_baselines.py}{152}{159}{candidate grid and cost weights}
\listing{experiments/nengo_nef_replication_2026-08-29/run_mujoco_nef_baselines.py}{201}{214}{rigid-body model and RK4 step (MuJoCo's body-to-world convention)}
\listing{experiments/nengo_nef_replication_2026-08-29/run_mujoco_nef_baselines.py}{247}{289}{\code{_nmpc}: rollout, cost \eqref{eq:nmpc}, selection, one-period self-check}
\listing{experiments/final_robustness_gating_2026-09-20/run_campaign_a5_nmpc_corrected.py}{25}{26}{corrected A5 settings: no update hold}

\implnote{(i) The rollout predicts the nominal family only; it does not know the residual that will be added. (ii) The first A5
NMPC run updated at 100\,Hz and held the torque between updates, a leftover from an earlier exploratory setting. The audit
(\code{AUDIT_NMPC_HOLD_RATE_2026-09-25.md}) showed that the hold, not the predictive machinery, caused that row's anomaly:
removing it restored NMPC, and imposing the same hold on MPC broke MPC identically. The paper's NMPC row is the full-rate re-run.
(iii) The class logs a one-period-ahead self-check of its own prediction (\code{nmpc_pred_err_mean_deg}); it was used to
diagnose (ii), not as a tuning signal.}

% =============================================================================
\section{Results across families (Sec.~VI-B)}\label{sec:results}
% =============================================================================
\begin{table}[H]\centering\small
\caption{Reduction of mean quaternion-angle error by the gated CMG--PES residual on each nominal family, noise-free, $\rho=3$,
45-s missions (\% , positive is better; ``5-mission'' is the reduction of geometric means). Read from
\ccode{results/a5_numbers.json}, which \ccode{analyze_a5.py} computes from the run summaries; NMPC from the corrected
full-rate run. Last column: baseline geometric mean (mrad).}\label{tab:results}
\begin{tabular}{@{}llrrrrrrr@{}}\toprule
Family & Platform & Roll & Full3 & Hard45 & Hardstep & Yaw360 & 5-mission & base (mrad)\\\midrule
\tabinput{generated/family_results.tex}
\bottomrule\end{tabular}\end{table}

These are the numbers the manuscript reports (X2/CF2: backstepping $37.3/22.4$, PID+ff $36.7/22.4$, LQR $31.1/22.4$, SMC
$43.8/26.6$, MPC $33.4/22.7$, NMPC $31.9/21.7$), and the claim verifier checks them (\code{reproduce/verify_acc2027_claims.py}).
What the table shows, without further interpretation:
\begin{itemize}[leftmargin=1.2em]
\item The five-mission reduction stays within $31$--$44\%$ on X2 and $22$--$27\%$ on CF2 across families whose baseline errors
  differ by up to $1.7\times$ (X2 geometric means $3.7$--$6.3$\,mrad), with the residual unchanged.
\item The reduction is largest on Roll and Yaw360 and smallest on Hardstep, whose step changes of tens of degrees drive every
  linear family far past the $5$--$6^\circ$ at which its proportional term reaches the box (Table~\ref{tab:gains}).
\item SMC is the only family with a clearly positive Hardstep result ($+34.2\%$ X2, $+12.2\%$ CF2). SMC is also the only
  family whose law is itself bounded below the box away from the origin (Sec.~\ref{sec:smc}); A5 did not isolate whether that
  is the cause, so it is recorded here as an open observation.
\end{itemize}

% =============================================================================
\section{Scope of the family comparison}
% =============================================================================
\begin{itemize}[leftmargin=1.2em]
\item The families are \emph{attachment tests}, not tuned competitors: all use the same bandwidth scale $\rho=3$, model
  feedforward off, and the A5 settings above. No family was retuned per mission or per platform, and the residual
  ($\beta$, $\eta$, $k_e$) is the same for all.
\item MPC is unconstrained and realized as its first-move gain; NMPC is a 27-point grid search around backstepping. Neither is a
  constrained optimal controller; the torque box acts through the shared saturation.
\item PID's integrator is weak (slow pole at $-0.32$\,s$^{-1}$) and is disclosed as undertuned in the A5 report.
\item All gains are designed on the per-axis double integrator \eqref{eq:design} with principal moments; the X2 inertia tensor is
  not diagonal, so the poles of Table~\ref{tab:poles} are design values, not measured closed-loop poles.
\end{itemize}

\appendix
\section{Source files and hashes}\label{app:hashes}
SHA-256 (first 16 hex digits) of each cited file, compared with the \code{source_sha256} field of one A5 run per family and
platform (Roll, classical arm; NMPC from the corrected run): \fhashMatch/\fhashChecked{} recorded hashes are identical.
\begin{table}[H]\centering\small
\begin{tabular}{@{}p{0.55\linewidth}ll@{}}\toprule
File & SHA-256 & vs.\ A5 manifests\\\midrule
\tabinput{generated/family_hashes.tex}
\bottomrule\end{tabular}\end{table}

\bibliographystyle{IEEEtran}
\bibliography{../v15_2026-09-29_corrected/references_ACC2027,family_refs}
\end{document}
